% 移动操作与复合机器人部：顺序 = 下列 \input 先后；章号 \chapter 自动。
% 教学序：轮式底盘非完整运动学(Pfaffian/可控性) → 底盘+臂联合运动学·冗余IK·停车位姿 → 运动学最优控制 → 操作技能接口。
% 承接：约束动力学(ch:constrained-dynamics 的 Pfaffian/Frobenius)、统一动力学(ch:unified-multimodal-mpc)、
%   机械臂运动学(ch:arm-kin)、操作空间阻抗(ch:operational-space)；与 P7 足式部、P3 协作部 \cref 互链不重复。
\part{移动操作与复合机器人}
\chapter{轮式底盘运动学与非完整约束：从 Pfaffian 矩阵到可控性}\label{ch:mm-wheeled}

% ════════════════════════════════════════════════════════════════
%  P9 移动操作部首章：轮式底盘的速度层运动学。
%  融入 Tutorial「轮式运动学与 Pfaffian」之全部理论，记号归一到本书 SE(2) 约定，
%  Frobenius/可积分类判据指向 ch:constrained-dynamics 去重，此处给轮式工程实例。
%  教学层遵循 v5.0：动机→反面→历史→理论→陷阱→练习 + 符号表/速查/故障排查。
% ════════════════════════════════════════════════════════════════

\begin{practice}[前置自测]
开始之前，先试答下列问题。若已能流畅作答，可只速读本章表格与速查卡；若卡住，正是本章要为你解决的。
\begin{enumerate}
  \item 一辆后轮驱动、前轮转向的汽车，能否\emph{瞬时}向正侧方平移一厘米？经过一连串机动后，能否到达正侧方一米处？这两个回答为什么不矛盾？
  \item 平面刚体的位形空间 $\SE(2)$ 有几个自由度？给体速度 $(v_x^b,v_y^b,\omega)$，世界系下的 $\dot{x},\dot{y}$ 怎么算？
  \item “约束”在数学上是什么？把质点限制在球面上的约束，与把轮子限制为“只能纯滚动”的约束，本质区别在哪？
  \item 给定一族速度约束 $\mathbf{A}(\mathbf{q})\dot{\mathbf{q}}=\mathbf{0}$，你打算如何判断它能否积分成位置约束 $\boldsymbol{\varphi}(\mathbf{q})=\mathbf{0}$？
  \item 差速底盘“不能横移”，可平行泊车时它分明侧向挪进了车位——这个“凭空多出来的侧向运动”从哪来的？能用一个数学对象精确刻画吗？
\end{enumerate}
\end{practice}

\section*{本章目标}
学完本章，你应当能够：
\begin{enumerate}
  \item \textbf{说清}非完整约束“限制瞬时速度方向、却不限制可达位形”的反直觉本质，并与完整约束严格区分；
  \item \textbf{建立} $\SE(2)$ 上的底盘运动学，在体坐标系与世界坐标系之间精确换算速度；
  \item \textbf{独立推导}单轮圆盘、差速、阿克曼、全向/麦轮四种底盘的 Pfaffian 约束矩阵 $\mathbf{A}(\mathbf{q})\dot{\mathbf{q}}=\mathbf{0}$，并据维度/秩读出瞬时自由度；
  \item \textbf{运用} Frobenius 判据（\cref{thm:cd-frobenius} 的轮式专化）判断约束可积性，并用李括号 $[\mathbf{g}_1,\mathbf{g}_2]$ 给出非完整性的显式证明；
  \item \textbf{解释} Chow--Rashevsky 定理与 LARC 如何保证差速车“绕路可达任意位形”，并用 Ball--Box 定理估计绕路代价 $O(\varepsilon^{1/k})$；
  \item \textbf{把}底盘写成漂移自由的仿射控制系统 $\dot{\mathbf{q}}=\sum_i u_i\mathbf{g}_i(\mathbf{q})$，为运动规划与最优控制铺路；
  \item \textbf{建模}纯滚动失效时的轮胎滑移——纵向滑移比 $\kappa$、侧向滑移角 $\alpha$、Pacejka 魔术公式与摩擦锥，并理解从 $\mathbf{A}\dot{\mathbf{q}}=\mathbf{0}$ 到 $\mathbf{A}\dot{\mathbf{q}}=\mathbf{r}(\kappa,\alpha)$ 的约束残差观点及其对里程计漂移的影响。
\end{enumerate}

\subsection*{本章知识导航}
本章是一条“\emph{把轮子与地面的物理接触，翻译成可计算的速度约束、再追问它在全局意义上限制了什么}”的主线。结构如下：

\begin{center}
\small
\begin{tikzpicture}[
  node distance=4.5mm and 7mm, >=Stealth,
  blk/.style={draw, rounded corners, align=center, font=\small, inner sep=3pt, fill=main!6, draw=main!60!black},
  grp/.style={draw, rounded corners, align=center, font=\small, inner sep=3pt, fill=second!6, draw=second!60!black},
  ln/.style={->, thick, gray!60!black}]
\node[blk] (mot) {动机\\不能横移?};
\node[blk, right=of mot] (se2) {$\SE(2)$ 运动学\\体$\leftrightarrow$世界速度};
\node[blk, right=of se2] (pf) {Pfaffian 矩阵\\四种底盘};
\node[grp, below=8mm of mot] (frob) {可积性\\Frobenius};
\node[grp, right=of frob] (chow) {可控性\\Chow / LARC};
\node[grp, right=of chow] (aff) {仿射系统\\$\dot{\mathbf{q}}=\sum u_i\mathbf{g}_i$};
\node[blk, below=8mm of frob, fill=third!8, draw=third!60!black] (slip) {滑移建模\\$\kappa,\alpha$, Pacejka};
\node[blk, right=of slip, fill=third!8, draw=third!60!black] (res) {约束残差\\$\mathbf{A}\dot{\mathbf{q}}=\mathbf{r}$};
\draw[ln] (mot)--(se2); \draw[ln] (se2)--(pf);
\draw[ln] (pf.south) -- ++(0,-4.5mm) -| (frob.north);
\draw[ln] (frob)--(chow); \draw[ln] (chow)--(aff);
\draw[ln] (frob.south)-- (slip.north);
\draw[ln] (slip)--(res);
\end{tikzpicture}
\end{center}

\noindent\textbf{推荐路径}：\cref{sec:mw-motivation,sec:mw-se2} 是任何读者的主干（把“横向不滑”这一句话精确化为速度约束）；\cref{sec:mw-pfaffian} 是核心推导，四种底盘的 $\mathbf{A}(\mathbf{q})$ 一次到位；\cref{sec:mw-integrability,sec:mw-controllability} 是\emph{理论高地}，回答“这约束是否可积、系统是否可控”，并复用约束力学章（\cref{ch:constrained-dynamics}）的判据；\cref{sec:mw-affine} 把底盘写成控制系统的标准形，是后续规划/最优控制章的接口；\cref{sec:mw-slip} 面向真实地面，松弛纯滚动假设。带工程色彩的滑移在线估计与里程计影响可在读完状态估计章后回看。

\subsection*{前置知识桥接}
本章把两条线收束到轮式底盘上。其一是\textbf{约束的分类与判据}：\cref{ch:constrained-dynamics} 已把一切一阶约束统一为 Pfaffian 形式 $\mathbf{A}(\mathbf{q})\dot{\mathbf{q}}+\mathbf{b}=\mathbf{0}$（\cref{eq:cd-pfaffian}），并给出可积性的 Frobenius 判据（\cref{thm:cd-frobenius}）与完整/非完整的形式定义（\cref{def:gm-constraint}）。本章不重复这些一般理论，而是把它们落到四种底盘的\emph{具体矩阵}上，并补全约束力学章一笔带过的可控性侧（Chow 定理）。其二是\textbf{平面位姿的表示}：刚体在平面上的位姿构成李群 $\SE(2)$，是 $\SE(3)$（\cref{ch:lie}）在二维的特例——只有一个转角 $\theta$，于是 $n_q=n_v=3$，本章直接采用这一简化。

\subsection*{如果跳过本章会怎样}
\begin{itemize}
  \item 在运动规划里，你会习惯性地用 A$^\star$/Dijkstra 在栅格上搜最短路，再让差速车跟踪——结果规划出含直角拐弯的路径，底层根本执行不了；不懂\emph{非完整约束}就不知道该换成 Hybrid A$^\star$ 或 state-lattice 这类\emph{尊重运动学}的规划器。
  \item 在轮式里程计里，你会默认轮速就是地面速度；一旦上了湿滑或松软地面，纯滚动假设失效而你毫无察觉，定位会以每秒数厘米的速度漂移——不懂\emph{约束残差} $\mathbf{A}\dot{\mathbf{q}}=\mathbf{r}(\kappa,\alpha)$ 就不知道何时、为何要降低对轮速观测的信任。
  \item 进入移动操作（本部后续章）后，底盘的 $\SE(2)$ 速度要与机械臂雅可比拼成联合雅可比 $\mathbf{V}_{ee}=[\mathbf{J}_b\ \mathbf{J}_a]\boldsymbol{\nu}$；底盘那一块 $\mathbf{J}_b$ 的结构正由本章的 $\mathbf{A}(\mathbf{q})$ 零空间决定。
\end{itemize}

\begin{note}
\textbf{本章记号与约定.}\;
底盘位形 $\mathbf{q}=(x,y,\theta)\in\SE(2)$（或在含轮转角/转向角时维数更高，逐处说明），向量粗体小写、矩阵粗体大写。
体坐标系速度记 $\mathbf{v}^b=(v_x^b,v_y^b,\omega)^\top$，$v_x^b$ 沿车体前向、$v_y^b$ 沿左向、$\omega=\dot\theta$ 为偏航角速度；世界系速度记 $\dot{\mathbf{q}}=(\dot x,\dot y,\dot\theta)^\top$。
非完整约束统一写成 Pfaffian 齐次形式 $\mathbf{A}(\mathbf{q})\dot{\mathbf{q}}=\mathbf{0}$（与 \cref{eq:cd-pfaffian} 在 $\mathbf{b}=\mathbf{0}$、定常情形一致）。
约束零空间张成的\emph{允许速度分布}记 $\mathcal{D}(\mathbf{q})=\ker\mathbf{A}(\mathbf{q})$，其上一组基向量场记 $\mathbf{g}_i(\mathbf{q})$，控制输入记 $u_i$。
李括号 $[\mathbf{g}_i,\mathbf{g}_j]$ 沿用 \cref{ch:lie} 的定义（向量场的对易子）。
\end{note}

% ════════════════════════════════════════════════════════════════
\section{动机：为什么轮式机器人不能随意横移}
\label{sec:mw-motivation}

\paragraph{从停车场说起。}
设想你在驾驶一辆普通轿车（后轮驱动、前轮转向），要从当前车道横移到相邻车道。直觉立刻告诉你：\emph{不能直接横挪}——车子只会前进、后退、转弯。可经过一串“前进—打方向—后退—回正”的操作，你确实能挪进相邻车位。这个再日常不过的经验，藏着一个深刻的数学性质：\textbf{速度受限，位置却不受限}。

\begin{table}[H]\centering\small
\caption{平行泊车经验揭示的三个层次}
\begin{tabular}{ll}
\toprule
特征 & 描述 \\
\midrule
瞬时限制 & 任意时刻速度方向受约束——不能垂直于车轮平面运动 \\
全局自由 & 给足时间与空间，车辆可到达平面上任意位置与朝向 \\
表象矛盾 & 速度受限，可位置不受限——这怎么可能同时成立？ \\
\bottomrule
\end{tabular}
\end{table}

\noindent 这种“瞬时速度受限、可达位形不变”的约束，叫\emph{非完整约束}（nonholonomic constraint）；与之相对的，把系统困在低维曲面上的，叫\emph{完整约束}（holonomic constraint）。二者的形式定义已由 \cref{def:gm-constraint} 给出，这里先把直觉立住。

\paragraph{反面：把两种约束混为一谈会出什么事。}
假设你为差速机器人写路径规划器。若\emph{错把非完整约束当完整约束}：规划器会以为机器人能沿任意方向移动，于是生成含横移、含直角拐弯的轨迹；底层控制器无法执行，机器人原地打转、轮胎刺耳摩擦。反过来，若\emph{错把非完整约束当成不可逾越的禁锢}（以为差速车只能走直线）：大量本可达的位形被误判为不可达，机器人被困在狭窄走廊里直行，丧失了“多次机动到达任意位形”的能力。正确的理解只有一句：非完整约束限制的是\textbf{瞬时运动方向}（微分约束），不是\textbf{最终可达集}（积分约束）。

\paragraph{历史：从经典力学到非线性控制。}
非完整约束的研究始于 $19$ 世纪末的分析力学。Hertz 最早系统讨论它在力学中的角色，Appell 给出处理非完整系统的方程，Chaplygin 求出滚动圆盘等系统的精确解；$1939$ 年 Chow（与 $1938$ 年的 Rashevsky 独立地）证明了非完整系统的可控性判据；$20$ 世纪 $70$ 年代起 Brockett 把非完整约束接入控制论，开创非线性控制；$1994$ 年 Murray、Li 与 Sastry 的《A Mathematical Introduction to Robotic Manipulation》\cite{murray1994mathematical} 把非完整运动规划系统化；其后 LaValle\cite{lavalle2006planning} 又把 RRT 等采样方法推广到非完整系统。本章正是这条脉络在轮式底盘上的落点。

\begin{insight}[完整 vs 非完整的一句话区分]\label{ins:mw-holo}
完整约束减少系统的\textbf{位形空间维度}（质点从三维空间被压到二维球面）；非完整约束减少系统的\textbf{瞬时速度自由度}，却\textbf{不减少}位形空间维度。前者让位形空间“塌缩”，后者只是给每一步的运动方向上了把锁，门后的房间一间不少。
\end{insight}

\begin{pitfall}[概念误区：“非完整 = 不可控”]
新手常想：“轮子不能横移，所以差速车的横向位置不可控。”错。非完整约束锁的是\emph{瞬时速度方向}，不是\emph{可达位形集}。差速车靠“前进—转弯—前进”的序列能到达平面上任意 $(x,y,\theta)$；其可控性由 Chow 定理保证（\cref{sec:mw-controllability}）。“不可控”指存在永远到不了的状态；“非完整”只是说到某些状态要\emph{绕路}。务必把这两件事分开。
\end{pitfall}

\begin{practice}
\begin{enumerate}
  \item 各举三例日常的完整约束与非完整约束，并说明每个限制的是位置还是速度。
  \item 一个冰球在冰面上自由滑行，它受的约束（若有）是完整还是非完整？若再加上“只能沿当前朝向滑”的约束，情况如何变化？
  \item 平面刚体位形空间 $\SE(2)$ 维度为 $3$。若施加一个非完整约束减少一个瞬时速度自由度，剩几个瞬时速度自由度？位形空间维度变了吗？
\end{enumerate}
\end{practice}

% ════════════════════════════════════════════════════════════════
\section{\texorpdfstring{$\SE(2)$}{SE(2)} 上的运动学与速度约束}
\label{sec:mw-se2}

\paragraph{位形与速度的两套坐标。}
轮式底盘在平面上运动，位形完全由\emph{位置} $(x,y)$ 与\emph{朝向} $\theta$ 确定，构成李群
\begin{equation}
  \mathbf{q}=\begin{pmatrix}x\\y\\\theta\end{pmatrix}\in\SE(2)\cong\mathbb{R}^2\times S^1,
  \qquad
  \mathbf{T}=\begin{pmatrix}\cos\theta&-\sin\theta&x\\\sin\theta&\cos\theta&y\\0&0&1\end{pmatrix}\in\SE(2).
  \label{eq:mw-se2}
\end{equation}
与 $\SE(3)$（\cref{ch:lie}）相比，平面上朝向只剩一个角 $\theta$，无需四元数，故位形与速度同维：$n_q=n_v=3$。这是 $S^1$ 相对 $\SO(3)$ 的简化，也是后续推导都在 $3$ 维里打转的原因。

\paragraph{体速度与世界速度的换算。}
刚体速度可在两个坐标系里写。\emph{世界系速度}（空间速度）就是位形的时间导数 $\dot{\mathbf{q}}=(\dot x,\dot y,\dot\theta)^\top$；\emph{体速度} $\mathbf{v}^b=(v_x^b,v_y^b,\omega)^\top$ 则把线速度投影到随车转动的前向/左向轴上。二者经平面旋转矩阵相联：
\begin{equation}
  \begin{pmatrix}\dot x\\\dot y\end{pmatrix}
  =\mathbf{R}(\theta)\begin{pmatrix}v_x^b\\v_y^b\end{pmatrix}
  =\begin{pmatrix}\cos\theta&-\sin\theta\\\sin\theta&\cos\theta\end{pmatrix}\begin{pmatrix}v_x^b\\v_y^b\end{pmatrix},
  \qquad \dot\theta=\omega.
  \label{eq:mw-body2world}
\end{equation}
逐分量展开即
\begin{equation}
  \dot x=v_x^b\cos\theta-v_y^b\sin\theta,\quad
  \dot y=v_x^b\sin\theta+v_y^b\cos\theta,\quad
  \dot\theta=\omega.
  \label{eq:mw-body2world-comp}
\end{equation}
这组式子是本章一切 Pfaffian 推导的“坐标基建”：物理约束写在体系里最自然（“侧向速度为零”就是 $v_y^b=0$），而约束矩阵 $\mathbf{A}(\mathbf{q})$ 必须落到世界系的 $\dot{\mathbf{q}}$ 上，二者靠 \cref{eq:mw-body2world} 来回搬运。

\begin{figure}[ht]
\centering
\begin{tikzpicture}[>=Stealth, scale=1.0]
  % world frame
  \draw[->, gray!70!black] (-0.4,-0.4) -- (1.4,-0.4) node[right]{$x$};
  \draw[->, gray!70!black] (-0.4,-0.4) -- (-0.4,1.2) node[above]{$y$};
  \node[gray!70!black] at (-0.7,-0.7) {世界系};
  % chassis at (2.6,0.9), heading theta
  \begin{scope}[shift={(3.0,0.7)}, rotate=28]
    \draw[thick, main!70!black, fill=main!8, rounded corners=2pt] (-0.7,-0.45) rectangle (0.7,0.45);
    % wheels
    \fill[gray!50!black] (-0.55,-0.5) rectangle (-0.25,-0.4);
    \fill[gray!50!black] (-0.55,0.4) rectangle (-0.25,0.5);
    \fill[gray!50!black] (0.25,-0.5) rectangle (0.55,-0.4);
    \fill[gray!50!black] (0.25,0.4) rectangle (0.55,0.5);
    % body axes
    \draw[->, second!70!black, very thick] (0,0) -- (1.1,0) node[right]{$\mathbf{e}_\parallel$ (前向 $v_x^b$)};
    \draw[->, third!70!black, very thick] (0,0) -- (0,0.95) node[above]{$\mathbf{e}_\perp$ ($v_y^b$)};
    \fill (0,0) circle (1.2pt);
  \end{scope}
  % position + heading annotation
  \draw[->, dashed, gray!60!black] (-0.4,-0.4) -- (3.0,0.7);
  \node[gray!60!black] at (1.3,0.45) {$(x,y)$};
  \draw[->, gray!70!black] (3.7,0.7) arc (0:28:0.7);
  \node[gray!70!black] at (4.1,0.95) {$\theta$};
\end{tikzpicture}
\caption{$\SE(2)$ 底盘的体坐标系与世界坐标系。前向轴 $\mathbf{e}_\parallel$、左向轴 $\mathbf{e}_\perp$ 随底盘转过 $\theta$；体速度 $(v_x^b,v_y^b,\omega)$ 经 \cref{eq:mw-body2world} 变到世界系 $(\dot x,\dot y,\dot\theta)$。}
\label{fig:mw-se2-frames}
\end{figure}

\begin{insight}[反事实：若不分体系与世界系会怎样]\label{ins:mw-frame}
取一辆朝向 $\theta=90^\circ$ 的车，以体系前向速度 $v_x^b=1$\,m/s 行驶。世界系里它其实在沿 $y$ 方向走（$\dot x=0,\dot y=1$）。若把体速度直接当世界速度用，就会算出车在沿 $x$ 走——方向完全错。这类错误在仿真里表现为机器人“螃蟹横行”。\cref{eq:mw-body2world} 的旋转矩阵正是杜绝它的唯一闸门。
\end{insight}

\paragraph{Pfaffian 形式的精确定义。}
现在可以把约束精确化。\emph{Pfaffian 约束}是对广义速度的线性约束，本章取其定常齐次形式
\begin{equation}
  \mathbf{A}(\mathbf{q})\dot{\mathbf{q}}=\mathbf{0},\qquad \mathbf{A}(\mathbf{q})\in\mathbb{R}^{k\times n},
  \label{eq:mw-pfaffian}
\end{equation}
其中 $k$ 是独立约束个数、$n$ 是位形维度。这正是 \cref{eq:cd-pfaffian} 在 $\mathbf{b}=\mathbf{0}$ 时的特例。$\mathbf{A}(\mathbf{q})$ 的每一行指定一个\emph{被禁止的速度方向}：系统速度不得有沿该方向的分量。所有合规速度构成\emph{允许速度分布}
\begin{equation}
  \mathcal{D}(\mathbf{q})=\{\mathbf{v}\in\mathbb{R}^n\mid\mathbf{A}(\mathbf{q})\mathbf{v}=\mathbf{0}\}=\ker\mathbf{A}(\mathbf{q}),
  \label{eq:mw-distribution}
\end{equation}
其维度 $\dim\mathcal{D}=n-\operatorname{rank}\mathbf{A}(\mathbf{q})$，即系统在该时刻的\textbf{瞬时运动自由度}。

\begin{table}[H]\centering\small
\caption{Pfaffian 约束的四个关键量}
\begin{tabular}{lll}
\toprule
概念 & 数学表示 & 物理含义 \\
\midrule
位形空间维度 & $n=\dim(\mathbf{q})$ & 描述位置需几个坐标 \\
约束个数 & $k=\operatorname{rank}(\mathbf{A})$ & 几个独立速度方向被禁 \\
瞬时自由度 & $n-k$ & 此刻能沿几个方向运动 \\
允许速度分布 & $\mathcal{D}(\mathbf{q})=\ker\mathbf{A}(\mathbf{q})$ & 此刻所有合规速度的集合 \\
\bottomrule
\end{tabular}
\end{table}

\begin{pitfall}[编程陷阱：在世界系写约束却忘了它依赖 $\theta$]
想表达“不能横移”，新手常直接写常数矩阵 $\mathbf{A}=[0,1,0]$（禁 $y$ 方向速度）。这只在 $\theta=0$ 时对，其他朝向全错。根因是 $\mathbf{A}(\mathbf{q})$ \emph{必须}是位形（尤其 $\theta$）的函数：“禁横向速度”在体系是常数约束 $v_y^b=0$，但经 \cref{eq:mw-body2world-comp} 搬到世界系就成了 $-\dot x\sin\theta+\dot y\cos\theta=0$，矩阵 $\mathbf{A}$ 随 $\theta$ 转。正确做法永远是：先在体系写约束，再用旋转矩阵变到世界系。
\end{pitfall}

\begin{insight}[轮式约束与足端接触是同一数学对象]\label{ins:mw-unify}
足式机器人的“足端静止接触” $\mathbf{J}_c(\mathbf{q})\mathbf{v}=\mathbf{0}$（见统一动力学 \cref{eq:umm-core} 中的接触项）与轮式的 $\mathbf{A}(\mathbf{q})\dot{\mathbf{q}}=\mathbf{0}$ \emph{形式完全相同}——令 $\mathbf{A}=\mathbf{J}_c$ 即可。区别只在物理：足端约束要求接触点速度为零（完整约束，可积成“接触点位置固定”），轮式约束要求接触点不横滑（非完整约束，不可积成位置关系）。这一统一视角是把底盘嵌入浮基统一动力学、乃至轮足混合系统的数学基础。
\end{insight}

\begin{practice}
\begin{enumerate}
  \item 写出 $\SE(2)$ 上从体速度到世界速度的完整 $3\times3$ 变换矩阵，并验证其行列式为 $1$（变换保体积）。
  \item 把约束 $v_y^b=0$ 转成 Pfaffian 形式 $\mathbf{A}(\mathbf{q})\dot{\mathbf{q}}=0$，写出 $\mathbf{A}(\mathbf{q})$，并核对 $\theta=0,\pi/4,\pi/2$ 时的具体形态。
  \item 一个机器人受约束 $v_x^b=v_y^b=0$（只能原地转）。这是完整还是非完整约束？允许速度分布维度是多少？
\end{enumerate}
\end{practice}

% ════════════════════════════════════════════════════════════════
\section{四种底盘的 Pfaffian 约束矩阵}
\label{sec:mw-pfaffian}

本节是全章核心。我们从物理约束出发，逐一推导四种主流底盘的 Pfaffian 矩阵。每种推导都走同一个\textbf{三步框架}：(1) 写出体系物理约束（纯滚动 + 不横滑）；(2) 经 \cref{eq:mw-body2world-comp} 变到世界系；(3) 组装 $\mathbf{A}(\mathbf{q})$，并据秩读出瞬时自由度。

\subsection{单轮圆盘：最小非完整系统}
\label{sec:mw-disk}

竖直圆盘在平面上滚动，是非完整约束的\emph{最小模型}：四个位形坐标、两个约束。它的配置由轮心位置、轮平面朝向、绕轴转角描述：
\begin{equation}
  \mathbf{q}=(x,\ y,\ \varphi,\ \theta)^\top\in\mathbb{R}^2\times S^1\times S^1,\qquad n=4,
  \label{eq:mw-disk-q}
\end{equation}
其中 $\varphi$ 是轮平面朝向（heading），$\theta$ 是绕自身轴的滚动角。在轮的体系里定义前向单位向量 $\mathbf{e}_\parallel=(\cos\varphi,\sin\varphi)^\top$ 与横向单位向量 $\mathbf{e}_\perp=(-\sin\varphi,\cos\varphi)^\top$。

\begin{derivation}[单轮圆盘的两条纯滚动约束]
记轮心平面位置 $\mathbf{p}=(x,y)^\top$。\textbf{约束一（纵向纯滚动）}：轮心沿前向的速度等于轮缘线速度 $r\dot\theta$——轮子向前滚一圈，轮心走过的距离等于轮缘展开的弧长。若不满足，轮子要么在地上空转、要么被拖着滑：
\begin{equation}
  \mathbf{e}_\parallel^\top\dot{\mathbf{p}}-r\dot\theta=0
  \;\Longrightarrow\;
  \dot x\cos\varphi+\dot y\sin\varphi-r\dot\theta=0.
  \label{eq:mw-disk-roll}
\end{equation}
\textbf{约束二（横向不滑）}：轮心沿轮轴方向的速度为零——轮子不侧滑（否则就像车在冰面上打横）：
\begin{equation}
  \mathbf{e}_\perp^\top\dot{\mathbf{p}}=0
  \;\Longrightarrow\;
  -\dot x\sin\varphi+\dot y\cos\varphi=0.
  \label{eq:mw-disk-noslip}
\end{equation}
把两式写成矩阵形式 $\mathbf{A}(\mathbf{q})\dot{\mathbf{q}}=\mathbf{0}$：
\begin{equation}
  \underbrace{\begin{pmatrix}\cos\varphi&\sin\varphi&0&-r\\-\sin\varphi&\cos\varphi&0&0\end{pmatrix}}_{\mathbf{A}(\mathbf{q})}
  \begin{pmatrix}\dot x\\\dot y\\\dot\varphi\\\dot\theta\end{pmatrix}=\begin{pmatrix}0\\0\end{pmatrix}.
  \label{eq:mw-disk-A}
\end{equation}
\end{derivation}

\noindent\textbf{维度核验}：$\mathbf{A}\in\mathbb{R}^{2\times4}$，两行线性无关（前两列构成的 $2\times2$ 块行列式 $\cos^2\varphi+\sin^2\varphi=1\neq0$），故 $\operatorname{rank}\mathbf{A}=2$，瞬时自由度 $=4-2=2$。物理上：圆盘任意时刻只能做两种运动——前/后滚（同时改变 $\theta$ 与 $(x,y)$）和改变朝向（改变 $\varphi$）；既不能横移，也不能在不转轮的情况下平移。

求 \cref{eq:mw-disk-A} 的零空间，得允许速度的参数化
\begin{equation}
  \dot{\mathbf{q}}=u_1\underbrace{\begin{pmatrix}r\cos\varphi\\r\sin\varphi\\0\\1\end{pmatrix}}_{\mathbf{g}_1(\mathbf{q})}
  +u_2\underbrace{\begin{pmatrix}0\\0\\1\\0\end{pmatrix}}_{\mathbf{g}_2(\mathbf{q})},
  \label{eq:mw-disk-vf}
\end{equation}
其中 $u_1=\dot\theta$（轮转速）、$u_2=\dot\varphi$（转向速率）是两个输入，$\mathbf{g}_1,\mathbf{g}_2$ 是张成 $\mathcal{D}(\mathbf{q})$ 的两个\emph{控制向量场}。

\begin{insight}[Pfaffian 零空间即控制向量场]\label{ins:mw-vf}
约束 $\mathbf{A}(\mathbf{q})\dot{\mathbf{q}}=\mathbf{0}$ 的零空间给出一组控制向量场 $\{\mathbf{g}_1,\dots,\mathbf{g}_m\}$，系统运动方程随之写成 $\dot{\mathbf{q}}=\sum_i u_i\mathbf{g}_i(\mathbf{q})$——这是\emph{漂移自由（drift-free）的仿射控制系统}的标准形（\cref{sec:mw-affine}）。非完整运动规划的全部工具（Chow 定理、RRT、steering function）都建立在这个形式之上。
\end{insight}

\subsection{差速底盘}
\label{sec:mw-diff}

差速底盘（differential drive）是最简单的实用底盘：两个独立驱动轮加一或多个被动万向轮，TurtleBot、iRobot Create 与大量教学机器人皆是。其位形只需三个变量（轮转角由运动学约束隐含）：$\mathbf{q}=(x,y,\theta)^\top\in\SE(2)$，$(x,y)$ 取两轮中点、$\theta$ 为底盘朝向。两轮转速 $\omega_L,\omega_R$ 映射到底盘速度：
\begin{equation}
  v_x^b=\frac{r(\omega_R+\omega_L)}{2},\qquad \omega=\dot\theta=\frac{r(\omega_R-\omega_L)}{L},
  \label{eq:mw-diff-fk}
\end{equation}
$r$ 为轮半径、$L$ 为轮距。差速底盘的\emph{唯一}约束是横向不滑 $v_y^b=0$，经 \cref{eq:mw-body2world-comp} 变到世界系：
\begin{equation}
  -\dot x\sin\theta+\dot y\cos\theta=0
  \;\Longrightarrow\;
  \mathbf{A}(\mathbf{q})=\begin{pmatrix}-\sin\theta&\cos\theta&0\end{pmatrix}.
  \label{eq:mw-diff-A}
\end{equation}
维度分析：$n=3$，$k=\operatorname{rank}\mathbf{A}=1$，瞬时自由度 $=3-1=2$（前进/后退 $+$ 原地旋转）。注意它与单轮圆盘的关键差别：差速底盘有两个独立电机输入 $(\omega_L,\omega_R)$，恰好等于两个瞬时自由度——系统在允许速度子空间内是\emph{完全驱动}的。运动方程
\begin{equation}
  \dot{\mathbf{q}}=\underbrace{\begin{pmatrix}\cos\theta\\\sin\theta\\0\end{pmatrix}}_{\mathbf{g}_1}v
  +\underbrace{\begin{pmatrix}0\\0\\1\end{pmatrix}}_{\mathbf{g}_2}\omega,
  \label{eq:mw-diff-vf}
\end{equation}
$v=v_x^b$ 为前向线速度、$\omega=\dot\theta$ 为角速度。这两个向量场将在 \cref{sec:mw-integrability,sec:mw-controllability} 反复出场，是全章的“主角”。

\subsection{阿克曼转向}
\label{sec:mw-ackermann}

阿克曼转向（Ackermann steering）是汽车与类车机器人的标准转向方式：转弯靠前轮偏转，而非两轮差速。其位形需四个变量：
\begin{equation}
  \mathbf{q}=(x,\ y,\ \theta,\ \delta)^\top,
  \label{eq:mw-ack-q}
\end{equation}
$(x,y)$ 取后轴中点、$\theta$ 为车体朝向、$\delta$ 为前轮转向角。约束仍是后轮不横滑：
\begin{equation}
  v_y^b=0\;\Longrightarrow\;-\dot x\sin\theta+\dot y\cos\theta=0,\qquad
  \mathbf{A}(\mathbf{q})=\begin{pmatrix}-\sin\theta&\cos\theta&0&0\end{pmatrix}.
  \label{eq:mw-ack-A}
\end{equation}
阿克曼运动学额外给出转向角 $\delta$ 与转弯半径 $R$ 的几何关系 $\tan\delta=L_{\text{wb}}/R$（$L_{\text{wb}}$ 为轴距），从而把偏航角速度与前向速度焊在一起：
\begin{equation}
  \dot\theta=\frac{v}{R}=\frac{v\tan\delta}{L_{\text{wb}}}.
  \label{eq:mw-ack-yaw}
\end{equation}
维度分析与差速相同（$n=4$ 时 $k=1$，对 $(x,y,\theta)$ 子空间瞬时自由度仍为这套约束所限），但实际输入只有 $2$ 个（前向速度 $v$、转向角速率 $\dot\delta$），故系统是\emph{欠驱动}的——$\delta$ 不受 Pfaffian 约束限制，却受机械转向角上限 $|\delta|\le\delta_{\max}$ 约束。运动方程
\begin{equation}
  \dot{\mathbf{q}}=\begin{pmatrix}\cos\theta\\\sin\theta\\\tan\delta/L_{\text{wb}}\\0\end{pmatrix}v
  +\begin{pmatrix}0\\0\\0\\1\end{pmatrix}\dot\delta.
  \label{eq:mw-ack-vf}
\end{equation}

\begin{insight}[反事实：若阿克曼没有最大转向角]\label{ins:mw-ackmax}
设想 $\delta_{\max}=90^\circ$。当 $\delta\to90^\circ$ 时 $\tan\delta\to\infty$，即便前向速度 $v$ 很小，偏航角速度 $\dot\theta$ 也趋于无穷——车辆绕后轴做极小半径回转，物理上退化为差速车的原地旋转。真实汽车 $\delta_{\max}$ 通常只有 $30^\circ$–$40^\circ$，这进一步抬高了最小转弯半径，正是停车入库要反复“揉库”的根源。
\end{insight}

\subsection{全向轮与麦克纳姆轮}
\label{sec:mw-omni}

全向底盘（如 Kuka youBot 底座、RoboCup 小型组机器人）用特殊轮消去横向不滑约束，使机器人能沿任意方向瞬时运动。\emph{全向轮}（Swedish/omni wheel）在轮缘装一圈自由滚动的小辊子，使轮子可沿轴向自由滑动，于是横向不滑约束被\emph{消除}。三轮全向底盘（$120^\circ$ 对称布置）的运动方程为
\begin{equation}
  \dot{\mathbf{q}}=\mathbf{J}^{-1}(\theta)\,(r\omega_1,\ r\omega_2,\ r\omega_3)^\top,
  \label{eq:mw-omni-fk}
\end{equation}
而其 Pfaffian 约束矩阵 $\mathbf{A}(\mathbf{q})=\varnothing$——\textbf{无约束}，全向底盘是\emph{完整系统}，位形空间维度 $=$ 瞬时自由度 $=3$，可瞬时沿任意方向（前后、左右、旋转的任意组合）运动。

\emph{麦克纳姆轮}（Mecanum wheel）是全向运动的另一实现：辊子以 $45^\circ$ 角安装（全向轮是 $90^\circ$）。四轮麦轮底盘（矩形布置）的逆运动学为
\begin{equation}
  \begin{pmatrix}\omega_1\\\omega_2\\\omega_3\\\omega_4\end{pmatrix}
  =\frac{1}{r}\begin{pmatrix}1&-1&-(L_x+L_y)\\1&1&(L_x+L_y)\\1&1&-(L_x+L_y)\\1&-1&(L_x+L_y)\end{pmatrix}
  \begin{pmatrix}v_x^b\\v_y^b\\\omega\end{pmatrix},
  \label{eq:mw-mecanum}
\end{equation}
$L_x,L_y$ 为轮到底盘中心的前后/左右距离。与全向轮一样，麦轮底盘\emph{没有}非完整约束，是完整系统；但工程上辊子与地面的接触力学使其\textbf{横向力传递效率低}，故侧向运动的最大加速度远小于前向——这是\emph{动力学}限制，而非运动学约束，切勿混入 $\mathbf{A}(\mathbf{q})$。

\begin{table}[H]\centering\small
\caption{四种底盘的 Pfaffian 矩阵与瞬时自由度对照}
\label{tab:mw-chassis}
\begin{tabular}{lccc c}
\toprule
底盘类型 & $\mathbf{A}(\mathbf{q})$ & $n$ & $\operatorname{rank}\mathbf{A}$ & 瞬时自由度 / 非完整？ \\
\midrule
单轮圆盘 & $\left[\begin{smallmatrix}\cos\varphi&\sin\varphi&0&-r\\-\sin\varphi&\cos\varphi&0&0\end{smallmatrix}\right]$ & $4$ & $2$ & $2$ / 是 \\[2pt]
差速 & $\left[\begin{smallmatrix}-\sin\theta&\cos\theta&0\end{smallmatrix}\right]$ & $3$ & $1$ & $2$ / 是 \\[2pt]
阿克曼 & $\left[\begin{smallmatrix}-\sin\theta&\cos\theta&0&0\end{smallmatrix}\right]$ & $4$ & $1$ & $3$（输入 $2$）/ 是 \\[2pt]
全向 / 麦轮 & $\varnothing$ & $3$ & $0$ & $3$ / 否 \\
\bottomrule
\end{tabular}
\end{table}

\begin{pitfall}[思维陷阱：多轮约束当成各轮约束的简单叠加]
四轮系统看似每个轮各两个约束、共 $8$ 个。但 $\SE(2)$ 只有 $3$ 个自由度，这 $8$ 个约束至多 $1$–$3$ 个独立，其余线性相关——是哪些独立，由轮距、轴距等\emph{轮间几何}决定。正确做法：先写出所有轮的约束，再分析堆叠后 $\mathbf{A}(\mathbf{q})$ 的秩，只保留独立行。直接“每个轮各自满足约束就行”会得到一个秩亏、自相矛盾的约束系统。
\end{pitfall}

\begin{pitfall}[概念误区：“全向底盘无约束，所以规划最简单”]
全向底盘\emph{无运动学约束}，但仍有\emph{动力学约束}（最大加速度、轮打滑上限）与\emph{环境约束}（避障）。况且其横向力传递效率低，实际最大横向加速度往往只有前向的 $50\%$–$70\%$。运动规划的难度不只来自运动学约束，还来自动力学与环境约束——“能横移”不等于“随便画直线就能跟踪”。
\end{pitfall}

\begin{practice}
\begin{enumerate}
  \item 为差速底盘推导正运动学：给 $\omega_L,\omega_R$，算 $(\dot x,\dot y,\dot\theta)$；验证 $\omega_L=\omega_R$ 时直行、$\omega_L=-\omega_R$ 时原地旋转。
  \item 推导自行车模型（阿克曼的简化，$\mathbf{q}=(x,y,\theta,\delta)$，$(x,y)$ 取后轮接地点）的 Pfaffian 约束矩阵，与汽车阿克曼矩阵 \cref{eq:mw-ack-A} 对比。
  \item 用符号计算（如 SymPy）求四轮麦轮底盘的正运动学矩阵，验证四轮同速同向时前进、对角轮反向时原地旋转。
\end{enumerate}
\end{practice}

% ════════════════════════════════════════════════════════════════
\section{可积性：Frobenius 判据的轮式专化}
\label{sec:mw-integrability}

\paragraph{悬而未决的问题。}
上一节推了四种底盘的 Pfaffian 约束。一个关键问题仍未回答：\textbf{如何判断一个 Pfaffian 约束是非完整的？}也就是说，$\mathbf{A}(\mathbf{q})\dot{\mathbf{q}}=\mathbf{0}$ 能否积分成位置约束 $\boldsymbol{\varphi}(\mathbf{q})=\mathbf{0}$？这不是数学好奇心——它直接决定规划策略：约束若可积，可达空间降维，规划在低维流形上做；约束若不可积，可达空间不变，但每步运动方向受限，规划要“绕路”。判错类型，规划器就会给出执行不了的路径。

一般判据已由约束力学章给出，本节只做轮式专化，不重复证明：

\begin{theorem}[可积判据 = Frobenius（轮式专化，引自 \cref{thm:cd-frobenius}）]\label{thm:mw-frobenius}
设允许速度分布 $\mathcal{D}(\mathbf{q})=\ker\mathbf{A}(\mathbf{q})=\operatorname{span}\{\mathbf{g}_1,\dots,\mathbf{g}_m\}$ 光滑。则 $\mathcal{D}$ 可积（存在 $n-m$ 个独立函数 $\varphi_i$ 使 $\nabla\varphi_i\perp\mathcal{D}$，约束完整）当且仅当 $\mathcal{D}$ \emph{对合}（involutive）：
\begin{equation}
  \forall\,i,j:\quad [\mathbf{g}_i,\mathbf{g}_j](\mathbf{q})\in\mathcal{D}(\mathbf{q}),
  \label{eq:mw-involutive}
\end{equation}
其中 $[\mathbf{g}_i,\mathbf{g}_j]=\dfrac{\partial\mathbf{g}_j}{\partial\mathbf{q}}\mathbf{g}_i-\dfrac{\partial\mathbf{g}_i}{\partial\mathbf{q}}\mathbf{g}_j$ 为李括号。等价地，以 $\omega^i$ 记 $\mathbf{A}$ 的第 $i$ 行（约束 $1$-形式），可积当且仅当 $\mathrm{d}\omega^i\wedge\omega^1\wedge\cdots\wedge\omega^k=0\ (\forall i)$；单约束（$k=1$）时简化为 $\mathrm{d}\omega\wedge\omega=0$。
\end{theorem}

\begin{insight}[李括号的几何含义]\label{ins:mw-bracket}
$[\mathbf{g}_i,\mathbf{g}_j]$ 度量“先沿 $\mathbf{g}_i$ 走一小步、再沿 $\mathbf{g}_j$ 走一小步”与“先 $\mathbf{g}_j$ 再 $\mathbf{g}_i$”两条路径终点之差。若两终点重合（括号为零或落在 $\mathcal{D}$ 内），约束可积；若终点不同（括号逸出 $\mathcal{D}$），说明交替使用两个输入能\emph{生成新的运动方向}——这正是非完整系统“绕路到达”的数学本质。它也类比微积分里的判据：向量场 $\mathbf{F}$ 是某标量梯度当且仅当 $\nabla\times\mathbf{F}=\mathbf{0}$（无旋）；Frobenius 定理把“无旋”推广成“对合”。
\end{insight}

\subsection{差速底盘：李括号显式计算}
\label{sec:mw-diff-bracket}

用 \cref{thm:mw-frobenius} 验证差速底盘的约束确实非完整。两个控制向量场取自 \cref{eq:mw-diff-vf}：
\begin{equation}
  \mathbf{g}_1=(\cos\theta,\ \sin\theta,\ 0)^\top,\qquad \mathbf{g}_2=(0,\ 0,\ 1)^\top.
  \label{eq:mw-diff-g}
\end{equation}

\begin{derivation}[{$[\mathbf{g}_1,\mathbf{g}_2]$ 恰是体系侧向方向}]
两个雅可比为
\[
  \frac{\partial\mathbf{g}_1}{\partial\mathbf{q}}=\begin{pmatrix}0&0&-\sin\theta\\0&0&\cos\theta\\0&0&0\end{pmatrix},\qquad
  \frac{\partial\mathbf{g}_2}{\partial\mathbf{q}}=\mathbf{0}.
\]
于是
\begin{equation}
  [\mathbf{g}_1,\mathbf{g}_2]
  =\frac{\partial\mathbf{g}_2}{\partial\mathbf{q}}\mathbf{g}_1-\frac{\partial\mathbf{g}_1}{\partial\mathbf{q}}\mathbf{g}_2
  =\mathbf{0}-\begin{pmatrix}-\sin\theta\\\cos\theta\\0\end{pmatrix}
  =\begin{pmatrix}\sin\theta\\-\cos\theta\\0\end{pmatrix}.
  \label{eq:mw-diff-bracket}
\end{equation}
检验它是否落在 $\mathcal{D}=\operatorname{span}\{\mathbf{g}_1,\mathbf{g}_2\}$ 内：若存在 $\alpha,\beta$ 使 $(\sin\theta,-\cos\theta,0)^\top=\alpha(\cos\theta,\sin\theta,0)^\top+\beta(0,0,1)^\top$，由第三行得 $\beta=0$，前两行给 $\alpha\cos\theta=\sin\theta$ 且 $\alpha\sin\theta=-\cos\theta$，相除得 $\tan\theta=-1/\tan\theta$，即 $\tan^2\theta=-1$，\emph{无实解}。故 $[\mathbf{g}_1,\mathbf{g}_2]\notin\mathcal{D}$，分布非对合，约束\textbf{不可积}，系统\textbf{非完整}。
\end{derivation}

\noindent 也可走微分形式一路核验：差速约束 $1$-形式 $\omega=-\sin\theta\,\mathrm{d}x+\cos\theta\,\mathrm{d}y$，其外微分 $\mathrm{d}\omega=-\cos\theta\,\mathrm{d}\theta\wedge\mathrm{d}x-\sin\theta\,\mathrm{d}\theta\wedge\mathrm{d}y$，与 $\omega$ 楔乘并用 $\mathrm{d}\theta\wedge\mathrm{d}y\wedge\mathrm{d}x=-\mathrm{d}\theta\wedge\mathrm{d}x\wedge\mathrm{d}y$ 合并，得
\begin{equation}
  \mathrm{d}\omega\wedge\omega=(\cos^2\theta+\sin^2\theta)\,\mathrm{d}\theta\wedge\mathrm{d}x\wedge\mathrm{d}y=\mathrm{d}\theta\wedge\mathrm{d}x\wedge\mathrm{d}y\neq0,
  \label{eq:mw-diff-form}
\end{equation}
两条路殊途同归：约束不可积。这与约束力学章对独轮车的同一结论（\cref{ex:cd-three}）一致。

\begin{insight}[平行泊车的数学本质]\label{ins:mw-parking}
$[\mathbf{g}_1,\mathbf{g}_2]=(\sin\theta,-\cos\theta,0)^\top$ 恰是体系的\emph{侧向}方向。这意味着：交替使用前进（$\mathbf{g}_1$）与转向（$\mathbf{g}_2$），差速车能凑出侧向运动——这正是平行泊车的内核。沿“前进 $\varepsilon$ → 转 $\varepsilon$ → 后退 $\varepsilon$ → 反转 $\varepsilon$”走一圈，机器人不回到原点，而是净侧移 $\Delta\mathbf{q}\approx\varepsilon^2[\mathbf{g}_1,\mathbf{g}_2]+O(\varepsilon^3)$。每轮机动产生 $O(\varepsilon^2)$ 的侧移，故要反复“揉”多次才挪进车位。李括号不是抽象记号，它精确地说出了“靠已有运动方向的交替，能造出哪个新方向”。
\end{insight}

\begin{pitfall}[编程陷阱：数值算李括号时步长选大了]
用有限差分近似 $\partial\mathbf{g}/\partial\mathbf{q}$ 时取步长 $h=0.1$（过大），算出的数值李括号与解析结果差异显著，尤其当 $\theta$ 接近 $0$ 或 $\pi$。根因是 $\mathbf{g}_1=(\cos\theta,\sin\theta,0)^\top$ 对 $\theta$ 的导数是三角函数，有限差分需很小步长才精确。正确做法：用符号微分（SymPy），或自动微分（如 JAX 的前向模式）。
\end{pitfall}

\begin{practice}
\begin{enumerate}
  \item 计算单轮圆盘（\cref{eq:mw-disk-vf} 的两个向量场）的李括号 $[\mathbf{g}_1,\mathbf{g}_2]$，验证它不在 $\mathcal{D}=\operatorname{span}\{\mathbf{g}_1,\mathbf{g}_2\}$ 内。
  \item 系统 $\mathbf{q}=(x,y,\theta)$ 受约束 $\dot x\sin\theta-\dot y\cos\theta+\ell\dot\theta=0$（$\ell$ 常数，拖车约束）。用 Frobenius 判据判断它是否可积。
  \item 证明全向底盘的约束平凡可积（因无约束）；再给全向底盘挂一节拖车（底盘 $(x,y,\theta)$ 无约束、拖车角 $\phi$ 受约束），写出整系统的 Pfaffian 矩阵并判可积性。
\end{enumerate}
\end{practice}

% ════════════════════════════════════════════════════════════════
\section{可控性：Chow--Rashevsky 定理与绕路代价}
\label{sec:mw-controllability}

\paragraph{从“不可积”到“可达”。}
Frobenius 判据告诉我们差速车的约束非完整——瞬时不能横移。但“瞬时不能横移”绝不等于“到不了侧方位置”。Chow--Rashevsky 定理回答最关键的问题：非完整系统到底能到达哪些位形？答案靠李代数秩条件给出。

\begin{theorem}[李代数秩条件（LARC）与 Chow--Rashevsky\cite{murray1994mathematical}]\label{thm:mw-chow}
设漂移自由控制系统 $\dot{\mathbf{q}}=\sum_{i=1}^m u_i\mathbf{g}_i(\mathbf{q})$，$\mathbf{g}_i$ 光滑。若在每一点 $\mathbf{q}$，向量场 $\{\mathbf{g}_i\}$ 及其各阶李括号张成的李代数维数等于位形空间维数 $n$（即\emph{LARC 成立}），则系统\textbf{小时间局部可控}（small-time locally controllable, STLC）——从任意位形出发，可在任意短时间内到达其邻域内任意位形。
\end{theorem}

对差速底盘，把 \cref{eq:mw-diff-bracket} 生成的 $\mathbf{g}_3=[\mathbf{g}_1,\mathbf{g}_2]=(\sin\theta,-\cos\theta,0)^\top$ 与原两场并列，检查秩：
\begin{equation}
  \operatorname{rank}\begin{pmatrix}\cos\theta&0&\sin\theta\\\sin\theta&0&-\cos\theta\\0&1&0\end{pmatrix}=3,\qquad
  \det=-1\cdot\big(\cos\theta\cdot(-\cos\theta)-\sin\theta\cdot\sin\theta\big)=\cos^2\theta+\sin^2\theta=1\neq0,
  \label{eq:mw-larc}
\end{equation}
处处满秩。即：$\mathbf{g}_1,\mathbf{g}_2$ 只张成 $2$ 维分布（不含横向），但加上一阶李括号 $\mathbf{g}_3$ 后张满整个 $\SE(2)$ 的 $3$ 维。LARC 成立 $\Rightarrow$ 差速车可从\emph{任意}位形到达\emph{任意}位形。

\begin{table}[H]\centering\small
\caption{Frobenius 与 Chow：两个判据回答两个不同问题}
\begin{tabular}{llll}
\toprule
定理 & 回答的问题 & 条件 & 结论 \\
\midrule
Frobenius & 约束是否可积？ & 李括号落在分布内 & 约束完整（可积） \\
Chow--Rashevsky & 系统是否可控？ & 李括号张满全空间 & 系统完全可控 \\
\bottomrule
\end{tabular}
\end{table}

\begin{insight}[同一个李括号，两种解读]\label{ins:mw-twoside}
有趣的是 Frobenius 与 Chow 用的是\emph{同一个}李括号 $[\mathbf{g}_1,\mathbf{g}_2]$，结论却互补。Frobenius 问它\emph{在不在}分布内：在，则约束可积、运动被困在子流形；Chow 问它\emph{补没补齐}全空间：补齐，则系统可控。对差速车，$[\mathbf{g}_1,\mathbf{g}_2]$ 逸出 $2$ 维分布（故非完整），又恰好补成 $3$ 维（故可控）——“非完整”与“可控”在同一个向量上达成了统一。
\end{insight}

\subsection{绕路代价：Ball--Box 定理}
\label{sec:mw-ballbox}

LARC 只保证\emph{存在}到达路径，不保证路径短。横向运动靠李括号“凑”出来，效率远低于直接可控的方向。Sub-Riemannian 几何里的\textbf{Ball--Box 定理}给出量化关系：若生成方向 $\mathbf{d}$ 需要 $k$ 阶李括号，则沿 $\mathbf{d}$ 移动距离 $\varepsilon$ 的最短路径长约为
\begin{equation}
  \ell(\varepsilon)\sim O\!\left(\varepsilon^{1/k}\right).
  \label{eq:mw-ballbox}
\end{equation}

\begin{table}[H]\centering\small
\caption{李括号阶数决定绕路代价}
\begin{tabular}{lcll}
\toprule
系统 & LARC 所需阶数 & 横移 $\varepsilon$ 的路径长 & 规划难度 \\
\midrule
全向底盘 & $0$（无需李括号） & $\varepsilon$（直线） & 低 \\
差速底盘 & $1$ & $O(\sqrt{\varepsilon})$ & 中 \\
阿克曼底盘 & $1$（含最小转弯半径） & $O(\sqrt{\varepsilon})$，下界更大 & 高 \\
拖车系统 & $2$ & $O(\varepsilon^{1/3})$ & 很高 \\
\bottomrule
\end{tabular}
\end{table}

\noindent 这把一条直觉事实量化了：差速车倒库比全向底盘难得多（要反复进退），拖着拖车倒库更是难上加难（要更多更细的机动）。阶数越高，“绕路”越长，规划越棘手。它对规划器选型有直接指导：完整系统（全向）可用 A$^\star$/RRT 任意路径搜索；非完整且 LARC 成立的系统须用 Hybrid A$^\star$、RRT + steering 或 state-lattice 这类尊重运动学的规划器。

\begin{pitfall}[概念误区：Chow 定理保证“快速到达”]
新手以为“差速车满足 LARC，所以能快速到任意位置”。Chow 定理只保证\emph{存在}到达路径，不保证路径长度或时间。靠李括号凑出的横移要反复机动，效率远低于直接方向：差速车横移 $1$\,m 的路径长远大于前进 $1$\,m。路径效率与李括号阶数挂钩——阶数越高，绕路越长。
\end{pitfall}

\begin{practice}
\begin{enumerate}
  \item 对阿克曼底盘算出所需的全部李括号，验证 LARC 是否成立；若成立，需要几阶？
  \item 差速车拖一节被动拖车，$\mathbf{q}=(x,y,\theta_1,\theta_2)$（$\theta_1$ 为拖车朝向）。写出控制向量场，验证 LARC 是否成立、需几阶李括号。
  \item （开放题）若某系统 LARC 需 $k$ 阶李括号，直觉上“绕路”到达侧向位置的路径长约是直线距离的多少倍？（提示：Ball--Box 定理 \cref{eq:mw-ballbox}。）
\end{enumerate}
\end{practice}

% ════════════════════════════════════════════════════════════════
\section{控制向量场视角：漂移自由仿射系统}
\label{sec:mw-affine}

把前面四种底盘的运动方程统一抽象，会发现它们都长成同一副模样：位形的时间导数是控制输入的\emph{线性组合}，且不含与输入无关的“漂移项”。这就是\textbf{漂移自由仿射控制系统}的标准形：
\begin{equation}
  \dot{\mathbf{q}}=\sum_{i=1}^{m}u_i\,\mathbf{g}_i(\mathbf{q})=\mathbf{G}(\mathbf{q})\,\mathbf{u},
  \qquad \mathbf{G}(\mathbf{q})=[\mathbf{g}_1(\mathbf{q})\ \cdots\ \mathbf{g}_m(\mathbf{q})],
  \label{eq:mw-affine}
\end{equation}
其中列向量 $\mathbf{g}_i$ 张成允许速度分布 $\mathcal{D}(\mathbf{q})=\ker\mathbf{A}(\mathbf{q})$，$\mathbf{u}=(u_1,\dots,u_m)^\top$ 是降维后的控制输入，$m=n-\operatorname{rank}\mathbf{A}$。差速底盘的 $\mathbf{G}=[\mathbf{g}_1\ \mathbf{g}_2]$（\cref{eq:mw-diff-vf}）、$\mathbf{u}=(v,\omega)^\top$；阿克曼的 $\mathbf{G}$ 与 $\mathbf{u}=(v,\dot\delta)^\top$ 见 \cref{eq:mw-ack-vf}；全向底盘 $\mathbf{A}=\varnothing$ 时 $\mathbf{G}=\mathbf{I}_3$、$\mathbf{u}=(v_x,v_y,\omega)$。

\begin{insight}[“漂移自由”的两重意义]\label{ins:mw-driftfree}
\cref{eq:mw-affine} 没有形如 $\mathbf{f}_0(\mathbf{q})$ 的漂移项，意味着：输入全为零时机器人静止（无惯性滑行），且运动方向完全由当下输入张成的 $\mathcal{D}(\mathbf{q})$ 决定。这带来两个好处：其一，正运动学就是把 $\mathbf{u}$ 经 $\mathbf{G}(\mathbf{q})$ 投到 $\dot{\mathbf{q}}$，规划器只需统一接口 $\dot{\mathbf{q}}=\mathbf{G}(\mathbf{q})\mathbf{u}$，换底盘只换 $\mathbf{G}$；其二，Chow 定理（\cref{thm:mw-chow}）正是为这类系统量身定制的可控性判据，本章一切“绕路可达”的结论都挂靠在此。注意这是\emph{运动学}模型；一旦把质量、滑移惯性纳入，就会出现漂移项，进入\emph{动力学}范畴（本部后续与 \cref{ch:constrained-dynamics}）。
\end{insight}

\noindent 把约束零空间显式参数化为 $\mathbf{G}(\mathbf{q})$，正是“速度命令投影到可行空间”的工程操作之源：给定一个可能违反约束的速度命令 $\mathbf{v}_{\text{cmd}}$，其在允许速度分布上的正交投影为
\begin{equation}
  \mathbf{v}_{\text{proj}}=(\mathbf{I}-\mathbf{A}^{\dagger}\mathbf{A})\,\mathbf{v}_{\text{cmd}},
  \label{eq:mw-project}
\end{equation}
$\mathbf{A}^\dagger$ 为 $\mathbf{A}$ 的 Moore--Penrose 伪逆，$\mathbf{I}-\mathbf{A}^\dagger\mathbf{A}$ 恰是到 $\ker\mathbf{A}(\mathbf{q})=\mathcal{D}(\mathbf{q})$ 的正交投影。差速底盘代入即把任意 $(v_x,v_y,\omega)$ 的横向分量抹掉、保留前向与转动——这与零空间投影在冗余机械臂里的用法（\cref{sec:umm-nullspace}）是同一套数学。

% ════════════════════════════════════════════════════════════════
\section{轮胎滑移与约束松弛：当纯滚动假设失效}
\label{sec:mw-slip}

\paragraph{纯滚动何时失效。}
前几节的推导都压在一个假设上：\emph{纯滚动 $+$ 不横滑}。它在低速（$<1$\,m/s）、干硬地面、轻载、直线运动时近似成立；但在湿滑地砖、松软沙土、急加速/急刹、高速过弯时严重失效。本节量化“滑了多少”，并把理想约束 $\mathbf{A}\dot{\mathbf{q}}=\mathbf{0}$ 松弛为带残差的 $\mathbf{A}\dot{\mathbf{q}}=\mathbf{r}(\kappa,\alpha)$。

\begin{insight}[滑移 = Pfaffian 约束的残差]\label{ins:mw-residual}
$\mathbf{A}(\mathbf{q})\dot{\mathbf{q}}=\mathbf{0}$ 描述的是\emph{理想运动学}——假设轮地刚性接触、零滑动。真实轮胎是弹性体，接触区有有限面积，轮胎力靠变形产生而非理想约束。滑移建模就是在 Pfaffian 框架里引入“约束的偏差”：从 $\mathbf{A}\dot{\mathbf{q}}=\mathbf{0}$ 变为
\begin{equation}
  \mathbf{A}(\mathbf{q})\dot{\mathbf{q}}=\mathbf{r}(\kappa,\alpha),
  \label{eq:mw-A-residual}
\end{equation}
残差 $\mathbf{r}$ 由滑移参数 $\kappa,\alpha$ 决定，$\kappa=\alpha=0$ 时退回理想纯滚动。
\end{insight}

\paragraph{两个滑移参数。}
滑移用两个无量纲量刻画。\emph{纵向滑移比} $\kappa$：
\begin{equation}
  \kappa=\frac{r\omega-v_x}{\max(|v_x|,\,v_{\text{low}})},
  \label{eq:mw-kappa}
\end{equation}
$\kappa=0$ 为纯滚动（$r\omega=v_x$），$\kappa>0$ 为驱动打滑（轮缘快于地面，如起步空转），$\kappa<0$ 为制动打滑（地面快于轮缘，如急刹），$\kappa=1$ 为完全空转（$v_x=0,\omega\neq0$），$\kappa=-1$ 为完全抱死（$\omega=0,v_x\neq0$）。分母里的 $v_{\text{low}}$（典型 $0.05$–$0.1$\,m/s）是\emph{低速保护}：低速时分母趋零会让 $\kappa$ 剧烈震荡、失去物理意义。\emph{侧向滑移角} $\alpha$：
\begin{equation}
  \alpha=\arctan\!\left(\frac{v_y}{|v_x|+v_{\text{low}}}\right),
  \label{eq:mw-alpha}
\end{equation}
$\alpha=0$ 无侧滑，$\alpha>0$ 向左侧滑、$\alpha<0$ 向右侧滑。

\paragraph{Pacejka 魔术公式。}
Hans Pacejka 提出的“魔术公式”（Magic Formula）是汽车工业最广用的轮胎力模型。称“魔术”是因为它\emph{没有直接的物理推导}，纯属经验拟合，却惊人地适配各种轮胎与工况。其纯纵向或纯侧向的简化形式为
\begin{equation}
  F=F_z\,D\,\sin\!\Big(C\,\arctan\!\big(B\,s-E\,(B\,s-\arctan(B\,s))\big)\Big),
  \label{eq:mw-pacejka}
\end{equation}
$s$ 是滑移参数（纵向取 $\kappa$、侧向取 $\alpha$），$B,C,D,E$ 是四个经验参数（见 \cref{tab:mw-pacejka}）。无论轮胎模型多复杂，接触力终究受\emph{摩擦锥}约束
\begin{equation}
  \sqrt{F_x^2+F_y^2}\le\mu F_z,
  \label{eq:mw-friction-cone}
\end{equation}
这是一切轮式接触约束的最终边界，与足式接触力的摩擦锥（\cref{eq:cd-signorini-foot} 所在的接触约束族）同源。

\begin{table}[H]\centering\small
\caption{Pacejka 魔术公式的四个参数}
\label{tab:mw-pacejka}
\begin{tabular}{lll}
\toprule
参数 & 名称 & 物理含义 \\
\midrule
$B$ & 刚度因子 & 小滑移时力--滑移斜率 \\
$C$ & 形状因子 & 曲线形状（纵向约 $1.65$，侧向约 $1.3$） \\
$D$ & 峰值因子 & 最大摩擦系数 $\mu_{\max}$ \\
$E$ & 曲率因子 & 峰值后的下降速率 \\
\bottomrule
\end{tabular}
\end{table}

\begin{figure}[ht]
\centering
\begin{tikzpicture}[>=Stealth, scale=1.0]
  \draw[->] (0,0) -- (5.4,0) node[right, font=\small]{滑移 $s$};
  \draw[->] (0,0) -- (0,3.0) node[above, font=\small]{$F/F_z$};
  % magic-formula-like curve: rise, peak, slight fall (control points, no trig units)
  \draw[very thick, main!70!black]
    plot[smooth, tension=0.7]
    coordinates {(0,0) (0.5,1.35) (1.0,2.25) (1.55,2.55) (2.3,2.4) (3.4,2.18) (5.0,2.05)};
  % peak marker
  \draw[dashed, gray!60!black] (0,2.55) -- (1.55,2.55) -- (1.55,0);
  \node[gray!60!black, font=\scriptsize, left] at (0,2.55) {$\mu_{\max}$};
  % region labels
  \draw[<->, second!70!black] (0,-0.55) -- (0.95,-0.55);
  \node[second!70!black, font=\scriptsize] at (0.5,-0.85) {线性区};
  \draw[<->, third!70!black] (0.95,-0.55) -- (2.3,-0.55);
  \node[third!70!black, font=\scriptsize] at (1.6,-0.85) {峰值区};
  \draw[<->, red!60!black] (2.3,-0.55) -- (5.0,-0.55);
  \node[red!60!black, font=\scriptsize] at (3.6,-0.85) {滑动区};
  % linear tangent at origin
  \draw[dotted, second!70!black] (0,0) -- (1.2,2.85);
\end{tikzpicture}
\caption{Pacejka 魔术公式 \cref{eq:mw-pacejka} 的典型 S 形力--滑移曲线：小滑移时力近似线性增长（斜率 $\propto B$），到峰值区达最大摩擦 $\mu_{\max}F_z$，进入滑动区后力反而下降。线性区可用 $F\approx C_s\,s$ 近似。}
\label{fig:mw-pacejka}
\end{figure}

\begin{table}[H]\centering\small
\caption{Pacejka 曲线的三个特征区}
\begin{tabular}{llll}
\toprule
区域 & 滑移范围 & 物理行为 & 控制意义 \\
\midrule
线性区 & $|s|<0.03$ & 力--滑移近似线性 $F\approx C_s s$ & 正常操作区，PD 控制有效 \\
峰值区 & $|s|\approx0.05$–$0.1$ & 摩擦达最大 $F_{\max}=DF_z$ & 最大制/驱动力，ABS/TCS 目标 \\
滑动区 & $|s|>0.15$ & 摩擦力下降（$F<F_{\max}$） & 失控区，须减速或切换模式 \\
\bottomrule
\end{tabular}
\end{table}

\paragraph{轮足/小轮场景：线性模型足够。}
对轮足机器人或小轮（半径约 $0.05$\,m），接触面积小、轮胎力学与汽车差异大，通常\emph{不必}精拟 Pacejka 参数——但力--滑移曲线的\emph{定性形状}仍适用（小滑移线性、大滑移饱和）。控制器设计常用线性轮胎模型即可：
\begin{equation}
  F_x=C_\kappa\,\kappa,\qquad F_y=C_\alpha\,\alpha,
  \label{eq:mw-linear-tire}
\end{equation}
$C_\kappa$ 为纵向刚度（滑移比 $\kappa$ 无量纲，故 $C_\kappa$ 单位为 N；小轮典型 $500$–$2000$\,N \pz{量级依轮径/材质/法向载荷而异，乘用车整轮的纵向刚度可达数万 N}）、$C_\alpha$ 为侧偏刚度（$\alpha$ 为弧度，故单位 N/rad；小轮典型 $300$–$1500$\,N/rad \pz{同样依平台而异：乘用车整轮 $C_\alpha$ 量级在 $10^4$–$10^5$\,N/rad，此处取值面向轻载小轮}），在 $|\kappa|<0.05$、$|\alpha|<3^\circ$ 内精度足够。

\paragraph{残差与里程计漂移。}
把滑移接回约束残差 \cref{eq:mw-A-residual}：当存在滑移时 $\mathbf{A}(\mathbf{q})\dot{\mathbf{q}}=\mathbf{r}(\kappa,\alpha)\neq\mathbf{0}$。线性/小角近似下，纵向残差 $r_{\text{lon}}\approx r\dot\theta\cdot\kappa$、侧向残差 $r_{\text{lat}}\approx v\sin\alpha$。$\kappa=\alpha=0$ 时残差归零，回到理想纯滚动；滑移越大，残差越大，纯滚动运动学的预测越不准。基于纯滚动假设的轮式里程计（wheel odometry）因此会积累定位误差，误差增长率近似
\begin{equation}
  \dot{e}_{\text{pos}}\approx v\,\sqrt{\kappa^2+\alpha^2}.
  \label{eq:mw-odom-drift}
\end{equation}
取 $v=1$\,m/s、$\kappa=0.05$、$\alpha=3^\circ$，得 $\dot{e}_{\text{pos}}\approx0.07$\,m/s——每秒漂移约 $7$\,cm，$10$\,秒就累积约 $70$\,cm，对“在门口精确放包裹”的配送机器人完全不可接受。这正是实际导航不能只靠轮式里程计、必须融合 IMU/视觉/LiDAR 的原因：滑移大时调高轮速观测的协方差、调低对里程计的信任，是约束残差观点给状态估计的直接指导。

\begin{pitfall}[编程陷阱：低速时直接套滑移比公式]
直接写 $\kappa=(r\omega-v_x)/v_x$、不加低速保护：机器人低速启动时分母趋零，$\kappa$ 突然飙到 $\pm100$，触发安全保护或令控制器震荡。正确做法：分母用 $\max(|v_x|,v_{\text{low}})$，并在 $|v_x|<v_{\text{low}}$ 时改用速度残差 $r\omega-v_x$（绝对量）而非比值 $\kappa$。
\end{pitfall}

\begin{pitfall}[概念误区：轮足机器人必须用 Pacejka 模型]
Pacejka 模型有 $6$–$20$ 个经验参数，要大量轮胎测试数据标定。轮足机器人的“轮胎”多是橡胶包覆的小轮，与汽车轮胎差异巨大。工程上轮足更常用\emph{线性力--滑移模型}（\cref{eq:mw-linear-tire}）或\emph{简单摩擦锥}（\cref{eq:mw-friction-cone}），因为参数少、可标定、算得快。按精度需求选模型复杂度：户外高速可考虑简化 Pacejka，室内低速线性模型足矣。
\end{pitfall}

\begin{practice}
\begin{enumerate}
  \item 实现带低速保护的纵向滑移比函数（输入 $v_{\text{wheel}},v_{\text{ground}},v_{\text{low}}$），对 $v_{\text{ground}}\in\{0.01,0.1,1.0\}$、$v_{\text{wheel}}\in\{0,0.5,1.0,2.0\}$ 各组合画热力图。
  \item 用 Pacejka 简化公式（取 $B=10,C=1.5,D=0.8,E=-0.5$）画侧向力 vs 滑移角曲线，标出峰值摩擦对应的滑移角，并解释峰值为何不在最大滑移角处。
  \item 设计“滑移感知权重调度器”：输入四轮滑移比，输出四个权重 $w_i\in[0,1]$ 调节状态估计对各轮速观测的信任，要求 $|\kappa_i|<0.05$ 时 $w_i=1$、$|\kappa_i|>0.3$ 时 $w_i=0$、中间平滑过渡。
\end{enumerate}
\end{practice}

% ════════════════════════════════════════════════════════════════
\section{承上启下}
\label{sec:mw-bridge}

本章把“轮子与地面的接触”翻译成了速度层的 Pfaffian 约束，并追问了它的两面：可积性（Frobenius，复用 \cref{thm:cd-frobenius}）与可控性（Chow，本章补全），最后用约束残差把理想纯滚动与真实滑移连了起来。这套\emph{速度约束 $\to$ 控制向量场 $\to$ 可控性}的链条，是移动操作一切运动学的地基。

下一章 \cref{ch:mm-kinematics} 将沿两个方向推进：其一，把本章 $\SE(2)$ 底盘的速度 $\mathbf{v}^b$ 与机械臂雅可比拼成移动操作的\emph{联合雅可比} $\mathbf{V}_{ee}=[\mathbf{J}_b\ \mathbf{J}_a]\boldsymbol{\nu}$，其中底盘块 $\mathbf{J}_b$ 的列正是本章 $\mathbf{G}(\mathbf{q})$（允许速度分布的基）；其二，处理底盘非完整约束与机械臂冗余（\cref{sec:ak-redundancy}）叠加后的零空间分配与奇异性（\cref{sec:ak-singularity}）。而本章把底盘写成漂移自由仿射系统 $\dot{\mathbf{q}}=\mathbf{G}(\mathbf{q})\mathbf{u}$ 的标准形，也正是后续移动操作最优控制（轨迹优化把它当作动力学约束）的接口。

% ════════════════════════════════════════════════════════════════
\section*{本章小结}

\begin{table}[H]\centering\small
\caption{符号表}
\begin{tabular}{lll}
\toprule
符号 & 含义 & 首次出现 \\
\midrule
$\mathbf{q}=(x,y,\theta)$ & 底盘位形（$\SE(2)$） & \cref{eq:mw-se2} \\
$\mathbf{v}^b=(v_x^b,v_y^b,\omega)$ & 体坐标系速度（前向/左向/偏航） & \cref{eq:mw-body2world} \\
$\mathbf{R}(\theta)$ & 平面旋转矩阵（体$\to$世界） & \cref{eq:mw-body2world} \\
$\mathbf{A}(\mathbf{q})$ & Pfaffian 约束矩阵 & \cref{eq:mw-pfaffian} \\
$\mathcal{D}(\mathbf{q})=\ker\mathbf{A}$ & 允许速度分布 & \cref{eq:mw-distribution} \\
$\mathbf{g}_i(\mathbf{q}),\ u_i$ & 控制向量场 / 输入 & \cref{eq:mw-disk-vf} \\
$[\mathbf{g}_i,\mathbf{g}_j]$ & 李括号 & \cref{eq:mw-involutive} \\
$\mathbf{G}(\mathbf{q})=[\mathbf{g}_1\,\cdots\,\mathbf{g}_m]$ & 仿射系统输入矩阵 & \cref{eq:mw-affine} \\
$r,\ L,\ L_{\text{wb}},\ \delta$ & 轮半径/轮距/轴距/转向角 & \cref{eq:mw-diff-fk,eq:mw-ack-q} \\
$\kappa,\ \alpha$ & 纵向滑移比 / 侧向滑移角 & \cref{eq:mw-kappa,eq:mw-alpha} \\
$B,C,D,E$ & Pacejka 四参数 & \cref{eq:mw-pacejka} \\
$\mathbf{r}(\kappa,\alpha)$ & 约束残差（滑移） & \cref{eq:mw-A-residual} \\
$\mu,\ F_z$ & 摩擦系数 / 法向力 & \cref{eq:mw-friction-cone} \\
\bottomrule
\end{tabular}
\end{table}

\begin{table}[H]\centering\small
\caption{定理 / 公式速查表}
\begin{tabular}{lll}
\toprule
定理 / 公式 & 一句话说明 & 对应 \\
\midrule
体$\to$世界速度 \cref{eq:mw-body2world} & 旋转矩阵搬运速度，约束推导的坐标基建 & \cref{sec:mw-se2} \\
Pfaffian 形式 \cref{eq:mw-pfaffian} & $\mathbf{A}\dot{\mathbf{q}}=\mathbf{0}$，零空间即瞬时自由度 & \cref{eq:mw-distribution} \\
四底盘 $\mathbf{A}(\mathbf{q})$ & 三步框架推导，秩读出自由度 & \cref{tab:mw-chassis} \\
Frobenius（专化） & $[\mathbf{g}_i,\mathbf{g}_j]\in\mathcal{D}$ $\Leftrightarrow$ 可积 & \cref{thm:mw-frobenius} \\
差速李括号 \cref{eq:mw-diff-bracket} & $[\mathbf{g}_1,\mathbf{g}_2]$ = 侧向 = 泊车本质 & \cref{ins:mw-parking} \\
Chow / LARC \cref{eq:mw-larc} & 李括号张满 $\Rightarrow$ 可控（绕路可达） & \cref{thm:mw-chow} \\
Ball--Box \cref{eq:mw-ballbox} & $k$ 阶括号 $\Rightarrow$ 绕路 $O(\varepsilon^{1/k})$ & \cref{sec:mw-ballbox} \\
仿射系统 \cref{eq:mw-affine} & $\dot{\mathbf{q}}=\mathbf{G}(\mathbf{q})\mathbf{u}$，漂移自由标准形 & \cref{sec:mw-affine} \\
滑移残差 \cref{eq:mw-A-residual} & $\mathbf{A}\dot{\mathbf{q}}=\mathbf{r}(\kappa,\alpha)$，纯滚动松弛 & \cref{ins:mw-residual} \\
Pacejka \cref{eq:mw-pacejka} & 经验力--滑移 S 曲线，四参数 & \cref{fig:mw-pacejka} \\
里程计漂移 \cref{eq:mw-odom-drift} & $\dot e\approx v\sqrt{\kappa^2+\alpha^2}$，融合动机 & \cref{sec:mw-slip} \\
\bottomrule
\end{tabular}
\end{table}

\begin{table}[H]\centering\small
\caption{知识点总表}
\begin{tabular}{clll}
\toprule
\# & 知识点 & 核心要点 & 难度 \\
\midrule
1 & 完整 vs 非完整 & 前者降位形维、后者锁瞬时速度方向 & $\bigstar\bigstar$ \\
2 & $\SE(2)$ 运动学 & 体$\leftrightarrow$世界速度，$n_q=n_v=3$ & $\bigstar\bigstar$ \\
3 & Pfaffian 推导 & 三步框架，四底盘 $\mathbf{A}(\mathbf{q})$ & $\bigstar\bigstar\bigstar$ \\
4 & Frobenius 可积性 & 李括号在不在分布内 & $\bigstar\bigstar\bigstar$ \\
5 & Chow / LARC & 李括号张满 $\Rightarrow$ 可控 & $\bigstar\bigstar\bigstar$ \\
6 & Ball--Box 绕路代价 & $O(\varepsilon^{1/k})$，阶数定难度 & $\bigstar\bigstar\bigstar\bigstar$ \\
7 & 仿射控制系统 & $\dot{\mathbf{q}}=\mathbf{G}(\mathbf{q})\mathbf{u}$ 标准形 & $\bigstar\bigstar\bigstar$ \\
8 & 滑移建模 & $\kappa,\alpha$、Pacejka、摩擦锥 & $\bigstar\bigstar\bigstar$ \\
9 & 约束残差与漂移 & $\mathbf{A}\dot{\mathbf{q}}=\mathbf{r}$，里程计融合 & $\bigstar\bigstar\bigstar$ \\
\bottomrule
\end{tabular}
\end{table}

\section*{延伸阅读}
\begin{itemize}
  \item \textbf{非完整运动规划的系统化源头}：Murray、Li 与 Sastry《A Mathematical Introduction to Robotic Manipulation》\cite{murray1994mathematical}——李括号、Chow 定理、非完整 steering 的经典论述（本章 \cref{sec:mw-integrability,sec:mw-controllability} 的主要依据）。
  \item \textbf{采样规划与非完整系统}：LaValle《Planning Algorithms》\cite{lavalle2006planning}——把 RRT 等推广到非完整系统、运动原语与 lattice 规划。
  \item \textbf{约束力学与可积性判据}：本书 \cref{ch:constrained-dynamics}（Pfaffian 统一形式、Frobenius/Signorini 判型、DAE 与接触）与 \cref{ch:geometric-mechanics}（完整/非完整的几何定义 \cref{def:gm-constraint}）——本章的一般理论出处。
  \item \textbf{现代运动学教材视角}：Lynch 与 Park《Modern Robotics》\cite{lynch2017modern}——$\SE(2)$/$\SE(3)$ 旋量理论与轮式机器人运动学的统一处理。
  \item \textbf{轮胎力学}：Pacejka《Tire and Vehicle Dynamics》\cite{pacejka2012tire}——魔术公式 \cref{eq:mw-pacejka} 的权威出处与参数标定；非完整力学几何见 Bloch《Nonholonomic Mechanics and Control》\cite{bloch2015nonholonomic}。
\end{itemize}

\section*{本章与后续章节的关系}
\begin{table}[H]\centering\small
\begin{tabular}{lll}
\toprule
后续章节 & 关系 & 本章铺垫 \\
\midrule
\cref{ch:mm-kinematics} 移动操作运动学 & 底盘 $+$ 臂联合雅可比 & $\mathbf{G}(\mathbf{q})$ 即底盘雅可比块 $\mathbf{J}_b$ \\
\cref{ch:constrained-dynamics} 约束动力学 & 把约束写进动力学 KKT & Pfaffian 形式、Frobenius 判据 \\
\cref{ch:unified-multimodal-mpc} 统一多模态 MPC & 接触/滚动约束的统一处理 & 接触=滚动的残差松弛视角 \\
\cref{sec:umm-nullspace} 零空间投影 & 命令投影到可行子空间 & \cref{eq:mw-project} 的零空间投影 \\
\bottomrule
\end{tabular}
\end{table}

\section*{故障排查手册}
\begin{enumerate}
  \item \textbf{症状}：差速/阿克曼机器人跟踪规划路径时\emph{频繁偏离、原地打转、轮胎打滑异响}。\textbf{排查}：规划器是否用了 A$^\star$/Dijkstra 等\emph{完整系统}搜索（\cref{sec:mw-motivation} 反面）？应换成 Hybrid A$^\star$/state-lattice 等尊重非完整约束的规划器。
  \item \textbf{症状}：约束矩阵\emph{只在 $\theta=0$ 时正确}、其他朝向运动方向错乱。\textbf{原因}：在世界系写了常数 $\mathbf{A}$、忘了它依赖 $\theta$（\cref{sec:mw-se2} 陷阱）。\textbf{对策}：先在体系写约束，再经 \cref{eq:mw-body2world-comp} 变到世界系。
  \item \textbf{症状}：数值李括号与解析结果\emph{差异大}、$\theta$ 近 $0$ 或 $\pi$ 时尤甚。\textbf{原因}：有限差分步长过大（\cref{sec:mw-diff-bracket} 陷阱）。\textbf{对策}：改用符号微分或自动微分。
  \item \textbf{症状}：机器人低速启动时滑移比\emph{突然飙到 $\pm100$}、触发保护或控制震荡。\textbf{原因}：滑移比分母趋零（\cref{sec:mw-slip} 陷阱）。\textbf{对策}：分母用 $\max(|v_x|,v_{\text{low}})$，极低速时改用速度残差而非比值。
  \item \textbf{症状}：纯轮式里程计\emph{定位持续漂移}，湿滑/松软地面尤甚。\textbf{原因}：滑移使 $\mathbf{A}\dot{\mathbf{q}}=\mathbf{r}\neq\mathbf{0}$、纯滚动假设失效（\cref{eq:mw-odom-drift}）。\textbf{对策}：融合 IMU/视觉/LiDAR，按估计滑移调高轮速观测协方差、调低里程计信任。
  \item \textbf{症状}：多轮底盘的约束系统\emph{秩亏、自相矛盾}。\textbf{原因}：把各轮约束简单叠加、忽略轮间几何耦合（\cref{tab:mw-chassis} 后陷阱）。\textbf{对策}：堆叠全部约束后分析 $\mathbf{A}(\mathbf{q})$ 的秩，只保留独立行。
\end{enumerate}
        % ch:mm-wheeled
\chapter{移动操作：底盘+臂联合运动学、冗余 IK 与停车位姿}\label{ch:mm-kinematics}

% ════════════════════════════════════════════════════════════════
%  移动操作（mobile manipulation）入门核心章：
%  把固定臂的微分运动学/冗余 IK 推广到「可移动且可能非完整」的浮动基座。
%  记号沿全书：向量粗体小写、矩阵粗体大写；右扰动 R·Exp(deltaφ)；
%  se(3) ξ=[rho;φ]；统一动力学 M v̇ + h = Sᵀtau + Σ Jᵀlambda + Jₑₑᵀf（\cref{eq:umm-core}）。
%  本波 5 章：ch:mm-wheeled / ch:mm-kinematics / ch:mm-ocp /
%  ch:mm-skill-interface / ch:arm-dualarm-planning。
% ════════════════════════════════════════════════════════════════

\begin{practice}[前置自测]
开始之前，先试答下列问题。能流畅作答，可只速读本章表格与符号表；若卡住，正是本章要为你解决的。
\begin{enumerate}
  \item 一台 $SE(2)$ 差速底盘，构型是 $3$ 维 $(x,y,\theta)$，可控输入却只有前进与偏航 $2$ 个量。它为什么\emph{最终}能停到平面上任意位姿、却\emph{瞬时}不能侧移？
  \item 固定臂的几何雅可比 $\mathbf{J}(\mathbf{q})$ 把关节速度映为末端 twist。若把基座装到一台会动的小车上，末端速度还应当加上哪几项？
  \item 给定期望末端 twist $\mathbf{V}_d$，为什么不能直接对“底盘 $3$ 速度 $+$ 臂 $n_a$ 速度”拼成的雅可比求伪逆、再把横向分量置零？
  \item “冗余度”在固定臂里指 $n-6$；移动操作平台多了底盘，冗余度怎么数？多出来的自由度能拿去做什么？
  \item 抓桌上一只杯子，底盘停太远够不到、停太近撞桌、朝向不对 IK 进奇异。有没有办法\emph{离线}把“底盘该停哪”预计算好，在线只查表？
\end{enumerate}
\end{practice}

\section*{本章目标}
学完本章，你应当能够：
\begin{enumerate}
  \item \textbf{建立}底盘 $+$ 臂的联合构型空间 $\mathbf{q}=(\mathbf{q}_b;\mathbf{q}_a)\in\SE(2)\times\mathbb{R}^{n_a}$，并写出末端位姿 $\mathbf{T}_{we}=\mathbf{T}_{wb}\mathbf{T}_{be}$；
  \item \textbf{推导}联合雅可比 $\mathbf{V}_{ee}=[\mathbf{J}_b\ \mathbf{J}_a]\boldsymbol{\nu}$，包括差速底盘的 $\mathbf{J}_b\in\mathbb{R}^{6\times2}$ 与全向底盘的 $\mathbf{J}_b\in\mathbb{R}^{6\times3}$；
  \item \textbf{解释}约简雅可比 $\bar{\mathbf{J}}=\mathbf{J}_{\text{full}}\bar{\mathbf{S}}$ 如何把非完整约束\emph{吸收进变量定义}，而非事后裁剪；
  \item \textbf{运用}阻尼伪逆与零空间投影，把“底盘少动 / 臂回中 / 朝向目标”等工程偏好写成数学项；
  \item \textbf{构造}任务优先级与分层 IK（递归零空间投影、含不等式约束的 HQP）来组织多级目标；
  \item \textbf{设计}停车位姿（base placement）与反向可达图 IRM，把在线 IK 搜索换成离线采样 $+$ 查表；
  \item \textbf{判断}差速底盘的“朝向”为何会进入可操作性，从而成为停车决策的关键量。
\end{enumerate}

\subsection*{本章知识导航}
本章是一条“\emph{把固定臂的微分运动学，搬到一个会跑、且可能不能横移的基座上}”的主线。五个知识块层层依赖：

\begin{center}
\small
\begin{tikzpicture}[
  node distance=5mm and 9mm, >=Stealth,
  blk/.style={draw, rounded corners, align=center, font=\small, inner sep=3pt, fill=main!6, draw=main!60!black},
  grp/.style={draw, rounded corners, align=center, font=\small, inner sep=3pt, fill=second!6, draw=second!60!black},
  ln/.style={->, thick, gray!60!black}]
\node[blk] (cfg) {建模\\ $\mathbf{q}=(\mathbf{q}_b;\mathbf{q}_a)$\\ $\mathbf{T}_{we}=\mathbf{T}_{wb}\mathbf{T}_{be}$};
\node[blk, right=of cfg] (jac) {联合雅可比\\ $[\mathbf{J}_b\ \mathbf{J}_a]$};
\node[blk, right=of jac] (red) {约简雅可比\\ $\bar{\mathbf{J}}=\mathbf{J}_{\text{full}}\bar{\mathbf{S}}$};
\node[grp, below=7mm of cfg] (ik) {冗余 IK\\ 阻尼伪逆 $+$ 零空间};
\node[grp, right=of ik] (hqp) {任务优先级\\ 分层 IK / HQP};
\node[blk, right=of hqp, fill=third!8, draw=third!60!black] (irm) {停车位姿\\ IRM / base placement};
\draw[ln] (cfg)--(jac); \draw[ln] (jac)--(red);
\draw[ln] (red.south) -- ++(0,-3.5mm) -| (ik.north);
\draw[ln] (ik)--(hqp); \draw[ln] (hqp)--(irm);
\end{tikzpicture}
\end{center}

\noindent\textbf{推荐路径}：\cref{sec:mmk-config,sec:mmk-jacobian,sec:mmk-reduced} 是任何读者的主干（建立“联合构型—联合雅可比—约简映射”三角）；\cref{sec:mmk-redundancy,sec:mmk-priority} 是\emph{求解所需}的核心工具（阻尼伪逆、零空间、分层），它们直接复用 \cref{thm:ak-dls} 的阻尼最小二乘与 \cref{sec:umm-nullspace} 的浮基零空间投影；\cref{sec:mmk-irm} 面向落地，把前面的可达性概念“烤”进离线表，引出 \cref{ch:mm-ocp} 的统一最优控制。

\subsection*{前置知识桥接}
本章站在四块前置内容之上。\cref{ch:mm-wheeled} 给了轮式底盘的运动学与\emph{准速度}（quasi-velocity）——本章把那里的差速/全向模型直接当作 $\mathbf{J}_b$ 的来源；\cref{ch:constrained-dynamics}（\cref{eq:cd-pfaffian} 的 Pfaffian 形式 $\mathbf{A}(\mathbf{q})\dot{\mathbf{q}}=0$ 与 \cref{thm:cd-frobenius} 的可积判据）给了非完整约束的判据，本章用它解释“约简”为何合法；\cref{ch:arm-kin}（几何雅可比、\cref{thm:ak-dls} 的 DLS 伪逆、\cref{sec:ak-singularity} 奇异、\cref{sec:ak-redundancy} 冗余）给了固定臂的全套微分运动学，本章把单臂的 $\mathbf{J}(\mathbf{q})$ 扩成 $[\mathbf{J}_b\ \mathbf{J}_a]$；\cref{ch:lie} 给了 $\SE(3)$、伴随与对数映射，本章在拼接 $\mathbf{T}_{we}$ 与定义末端误差时反复使用。本章要解决的\textbf{新问题}是：当“基座”不再固定、而是一个可移动且可能受非完整约束的 $SE(2)$ 物体时，如何统一描述、求解并落地“底盘 $+$ 臂共同完成末端任务”。

\subsection*{如果跳过本章会怎样}
\begin{itemize}
  \item 进入 \cref{ch:mm-ocp} 的统一最优控制时，你会看不懂状态为什么写成 $(x_b,y_b,\theta_b,\mathbf{q}_a)$、末端代价里的雅可比从哪来——因为联合雅可比与约简映射正是本章建立的；
  \item 部署真机移动抓取时，你会反复踩“导航到位却抓不到”“底盘 footprint 安全但臂撞货架”的坑，而这些坑的根都在本章的可达性、停车位姿与误差传播里。
\end{itemize}

\begin{note}[本章记号与约定]\label{note:mmk-notation}
底盘位姿 $\mathbf{q}_b=(x_b,y_b,\theta_b)\in\SE(2)$，臂关节 $\mathbf{q}_a\in\mathbb{R}^{n_a}$，联合构型 $\mathbf{q}=(\mathbf{q}_b;\mathbf{q}_a)$。变换记号统一用\emph{下标对}：$\mathbf{T}_{ab}$ 表“把 $b$ 系中的点变换到 $a$ 系”，等价于“$\{b\}$ 在 $\{a\}$ 中的位姿”，故 $\mathbf{p}_a=\mathbf{T}_{ab}\mathbf{p}_b$，$\mathbf{T}_{bw}=\mathbf{T}_{wb}^{-1}$。末端 twist 记 $\mathbf{V}_{ee}=[\mathbf{v}_{ee};\boldsymbol{\omega}_{ee}]\in\mathbb{R}^6$（线速度在前、角速度在后，与本书 $\se(3)$ 排序 $\boldsymbol{\xi}=[\boldsymbol{\rho};\boldsymbol{\phi}]$ 一致）；$\mathbf{e}_x,\mathbf{e}_y,\mathbf{e}_z$ 为标准基，$\mathbf{p}^\wedge$ 为 $\mathbb{R}^3$ 的反对称化算子（\cref{ch:lie}），故 $\mathbf{a}\times\mathbf{b}=\mathbf{a}^\wedge\mathbf{b}$。粗体小写为向量、粗体大写为矩阵。\textbf{可控速度 $\boldsymbol{\nu}$}（准速度，见 \cref{ch:mm-wheeled}）不是位形导数 $\dot{\mathbf{q}}$：差速底盘 $\boldsymbol{\nu}=(v,\omega,\dot{\mathbf{q}}_a)\in\mathbb{R}^{2+n_a}$，全向底盘 $\boldsymbol{\nu}\in\mathbb{R}^{3+n_a}$。
\end{note}

% ════════════════════════════════════════════════════════════════
\section{动机：会动的工作空间}
\label{sec:mmk-motivation}

\paragraph{固定臂的世界很清楚。}
一台螺栓固定在工作台上的机械臂，基座位姿是常量，末端能到达的位姿集合——\emph{工作空间}（workspace）——也是一个固定的几何体。逆运动学（IK）问题因此干净利落：给定目标位姿 $\mathbf{T}_g$，要么它落在工作空间内、有解，要么落在外、无解；规划就是“在这个固定集合里找一条无碰路径”。

\paragraph{反面：把臂装上小车，工作空间就开始“漂移”。}
现在把同一台臂装到一台可移动底盘上（如 Stretch、Ridgeback+UR5、TIAGo、Mobile ALOHA、Go2-W$+$小臂）。底盘一动，整条臂的工作空间随之平移、旋转；反过来，臂伸出去又会改变底盘“停在哪儿合适”。一个看似省事的拆分是：用导航栈把底盘开到桌前，再用 IK/规划栈抓杯子。它工程上成熟，却有一串绕不开的失败模式：
\begin{itemize}
  \item 底盘停\emph{太远}，末端够不到目标；
  \item 底盘停\emph{太近}，臂或负载与桌子碰撞；
  \item 底盘\emph{朝向}不合适，IK 落进奇异位形（\cref{sec:ak-singularity}），关节速度爆炸；
  \item 导航路径可行，但到位后的操作位姿不可达；
  \item 操作路径可行，但底盘在执行中自己挡住了臂。
\end{itemize}
这些都不是“导航”或“抓取”单独的 bug，而是\emph{二者没有联合考虑}的必然结果。

\begin{insight}[移动操作的本质洞察：工作空间从“集合”变成“集合的函数”]
固定臂：末端任务 $=$“在一个固定工作空间 $\mathcal{W}$ 内找解”。移动操作：底盘位姿 $\mathbf{q}_b$ 把工作空间变成一个随 $\mathbf{q}_b$ 平移/旋转的族 $\mathcal{W}(\mathbf{q}_b)$。于是规划问题升级为两步耦合的问题——\emph{先把工作空间移动到合适位置，再在其中找解}。本章的全部技术，无非是把这句话翻译成可计算的雅可比、伪逆与查表。
\end{insight}

\paragraph{两条主线。}
应对“会动的工作空间”，本章给两条互补的主线。其一是\emph{微分运动学}：把底盘速度与臂速度统一成一个联合雅可比，在速度级同时驱动两者完成末端任务（\cref{sec:mmk-jacobian,sec:mmk-reduced,sec:mmk-redundancy,sec:mmk-priority}）——这是“边到位边操作”的连续视角。其二是\emph{停车位姿}：在动手之前先回答“底盘该停哪儿”，用离线预计算把昂贵的在线搜索压成一次查表（\cref{sec:mmk-irm}）——这是“先到位再操作”的离散视角。两条线在 \cref{ch:mm-ocp} 的统一最优控制里合流。

% ════════════════════════════════════════════════════════════════
\section{联合构型空间与末端位姿}
\label{sec:mmk-config}

\paragraph{把底盘也算成“关节”。}
要统一描述底盘与臂，第一步是把底盘位姿当成系统状态的一部分。底盘位姿 $\mathbf{q}_b=(x_b,y_b,\theta_b)\in\SE(2)$；臂关节 $\mathbf{q}_a\in\mathbb{R}^{n_a}$。二者拼成\emph{联合构型}

\begin{definition}[移动操作的联合构型空间]\label{def:mmk-config}
移动操作平台的联合构型为
\begin{equation}
  \mathbf{q}=\begin{bmatrix}\mathbf{q}_b\\\mathbf{q}_a\end{bmatrix}\in\SE(2)\times\mathbb{R}^{n_a},\qquad
  \dim(\mathbf{q})=3+n_a.
  \label{eq:mmk-config}
\end{equation}
末端执行器在世界系的位姿由两段变换\emph{串联}给出：
\begin{equation}
  \boxed{\;\mathbf{T}_{we}(\mathbf{q})=\mathbf{T}_{wb}(\mathbf{q}_b)\,\mathbf{T}_{be}(\mathbf{q}_a)\in\SE(3),\;}
  \label{eq:mmk-fk}
\end{equation}
其中 $\mathbf{T}_{wb}$ 是底盘在世界系的位姿（由 $\mathbf{q}_b$ 嵌入 $\SE(3)$：平面平移 $+$ 绕 $z$ 偏航），$\mathbf{T}_{be}$ 是臂的正运动学（末端在底盘系的位姿，\cref{ch:arm-kin}）。
\end{definition}

以一台 $7$ 自由度臂为例，$\dim(\mathbf{q})=3+7=10$。固定臂只优化 $\mathbf{q}_a$；移动操作平台要\emph{同时}优化 $\mathbf{q}_b$ 与 $\mathbf{q}_a$——\cref{eq:mmk-fk} 里 $\mathbf{T}_{wb}$ 不再是常量，这正是后面一切复杂性的源头。

\begin{pitfall}[概念误区：$\mathbf{T}_{wb}$ 与 $\mathbf{T}_{bw}$ 反着传]
新手最常见的工程 bug 是“一会儿用 $\mathbf{T}_{wb}$、一会儿实际传进 $\mathbf{T}_{bw}$”。本章约定 $\mathbf{T}_{ab}$ 永远是“$\{b\}$ 在 $\{a\}$ 中的位姿”，于是 $\mathbf{p}_w=\mathbf{T}_{wb}\mathbf{p}_b=\mathbf{T}_{wb}\mathbf{T}_{be}\mathbf{p}_e$ 一路右乘、方向自洽。后面的联合雅可比与 $\SE(3)$ 误差全部沿用这个方向；写代码时核对一次下标对，能省下大量“末端位姿差一个逆”的调试。
\end{pitfall}

\paragraph{$SE(2)$ 底盘的速度变换。}
底盘在\emph{车体系}的速度 $(v_x^B,v_y^B,\omega)$ 与\emph{世界系}的位形速度 $\dot{\mathbf{q}}_b$ 之间，差一个绕 $z$ 的平面旋转 $\mathbf{R}(\theta_b)\in\SO(2)$：
\begin{equation}
  \begin{bmatrix}\dot x_b\\\dot y_b\end{bmatrix}=\mathbf{R}(\theta_b)\begin{bmatrix}v_x^B\\v_y^B\end{bmatrix},\qquad
  \mathbf{R}(\theta_b)=\begin{bmatrix}\cos\theta_b&-\sin\theta_b\\\sin\theta_b&\cos\theta_b\end{bmatrix},\qquad
  \dot\theta_b=\omega.
  \label{eq:mmk-se2vel}
\end{equation}
这是 $\SE(2)$ 上 body 与 world 速度互换的最简形式（完整推导见 \cref{ch:mm-wheeled}）。差速底盘还要再叠一层非完整约束（不能独立命令 $v_y^B$），下一节会看到它如何改变雅可比的列数。

% ════════════════════════════════════════════════════════════════
\section{联合雅可比}
\label{sec:mmk-jacobian}

\paragraph{末端速度 = 底盘的贡献 + 臂的贡献。}
末端 twist $\mathbf{V}_{ee}\in\mathbb{R}^6$ 由底盘与臂\emph{共同}产生。臂关节速度的贡献正是固定臂的几何雅可比 $\mathbf{J}_a\dot{\mathbf{q}}_a$（\cref{ch:arm-kin}）；底盘运动的贡献记为 $\mathbf{J}_b\mathbf{u}_b$，其中 $\mathbf{u}_b$ 是底盘的可控输入。叠加并按可控速度 $\boldsymbol{\nu}=(\mathbf{u}_b;\dot{\mathbf{q}}_a)$ 合并：

\begin{definition}[联合雅可比]\label{def:mmk-jac}
\begin{equation}
  \mathbf{V}_{ee}=\mathbf{J}_b(\mathbf{q})\,\mathbf{u}_b+\mathbf{J}_a(\mathbf{q})\,\dot{\mathbf{q}}_a
  =\underbrace{\begin{bmatrix}\mathbf{J}_b(\mathbf{q})&\mathbf{J}_a(\mathbf{q})\end{bmatrix}}_{\mathbf{J}(\mathbf{q})}\,\boldsymbol{\nu},
  \qquad \boldsymbol{\nu}=\begin{bmatrix}\mathbf{u}_b\\\dot{\mathbf{q}}_a\end{bmatrix}.
  \label{eq:mmk-joint-jac}
\end{equation}
\end{definition}

\noindent 关键在于 $\mathbf{J}_b$——它把底盘的平面运动“抬升”成末端的 $6$ 维 twist。下面分两类底盘显式推导。

\subsection{差速底盘的 \texorpdfstring{$\mathbf{J}_b$}{Jb}}
\label{sec:mmk-jac-diff}

\paragraph{几何推导。}
差速底盘的可控输入是 $\mathbf{u}_b=(v,\omega)$（前进线速度、偏航角速度）。设末端相对底盘的位置在世界系为 $\mathbf{p}_{be}^w=\mathbf{R}_{wb}\,\mathbf{p}_{be}^b$。逐项考察底盘运动对末端的速度贡献：

\begin{derivation}[差速底盘对末端 twist 的贡献]
\emph{前进 $v$ 的贡献.}\;底盘沿车体 $x$ 轴平移，世界系下该方向是 $\mathbf{R}_{wb}\mathbf{e}_x$。它是\emph{纯平移}，只贡献末端\textbf{线}速度 $\mathbf{R}_{wb}\mathbf{e}_x\,v$，不贡献角速度。

\emph{偏航 $\omega$ 的贡献.}\;底盘绕过其旋转中心、轴向 $\mathbf{e}_z$ 的瞬时转动。对一个刚体绕过原点、角速度 $\boldsymbol{\omega}_b=\omega\mathbf{e}_z$ 的转动，位于 $\mathbf{p}_{be}^w$ 处的点线速度为 $\boldsymbol{\omega}_b\times\mathbf{p}_{be}^w=\omega\,(\mathbf{e}_z\times\mathbf{p}_{be}^w)$；同时末端整体随底盘转，角速度贡献为 $\omega\mathbf{e}_z$。

\emph{合并.}\;把“线速度行 / 角速度行”按 $\mathbf{V}_{ee}=[\mathbf{v}_{ee};\boldsymbol{\omega}_{ee}]$ 排好，得
\begin{equation}
  \mathbf{V}_{ee}\big|_{\text{base}}=
  \begin{bmatrix}\mathbf{R}_{wb}\mathbf{e}_x & \mathbf{e}_z\times\mathbf{p}_{be}^w\\[2pt]
  \mathbf{0}_{3\times1} & \mathbf{e}_z\end{bmatrix}
  \begin{bmatrix}v\\\omega\end{bmatrix}.
  \label{eq:mmk-Jb-diff-deriv}
\end{equation}
\end{derivation}

\begin{proposition}[差速底盘的 $\mathbf{J}_b$]\label{prop:mmk-Jb-diff}
\begin{equation}
  \boxed{\;\mathbf{J}_b^{\text{diff}}=
  \begin{bmatrix}\mathbf{R}_{wb}\mathbf{e}_x & \mathbf{e}_z\times\mathbf{p}_{be}^w\\[2pt]
  \mathbf{0}_{3\times1} & \mathbf{e}_z\end{bmatrix}\in\mathbb{R}^{6\times2}.\;}
  \label{eq:mmk-Jb-diff}
\end{equation}
其中交叉积也可写成反对称矩阵形式 $\mathbf{e}_z\times\mathbf{p}_{be}^w=\mathbf{e}_z^\wedge\,\mathbf{p}_{be}^w=-(\mathbf{p}_{be}^w)^\wedge\mathbf{e}_z$；具体符号取决于 twist 约定，工程实现中必须与所用库（如 Pinocchio）的旋量排列严格统一。
\end{proposition}

\subsection{全向底盘的 \texorpdfstring{$\mathbf{J}_b$}{Jb}}
\label{sec:mmk-jac-omni}

\paragraph{多一个横向列。}
全向底盘（麦克纳姆/全向轮）可独立命令车体横向速度 $v_y^B$，输入变为 $\mathbf{u}_b=(v_x^B,v_y^B,\omega)$。横向平移在世界系沿 $\mathbf{R}_{wb}\mathbf{e}_y$，同样只贡献线速度。于是 $\mathbf{J}_b$ 比差速多一列：

\begin{proposition}[全向底盘的 $\mathbf{J}_b$]\label{prop:mmk-Jb-omni}
\begin{equation}
  \mathbf{J}_b^{\text{omni}}=
  \begin{bmatrix}\mathbf{R}_{wb}\mathbf{e}_x & \mathbf{R}_{wb}\mathbf{e}_y & \mathbf{e}_z\times\mathbf{p}_{be}^w\\[2pt]
  \mathbf{0} & \mathbf{0} & \mathbf{e}_z\end{bmatrix}\in\mathbb{R}^{6\times3}.
  \label{eq:mmk-Jb-omni}
\end{equation}
\end{proposition}

\noindent 这一列直接把冗余度提高 $1$（\cref{sec:mmk-redundancy}），也使全向底盘在“横向微调末端”上比差速底盘多一个\emph{直接}手段：差速底盘要横移末端 $5\,$cm，必须“转弯—前进—转回”几秒钟的折返，全向底盘则一步侧移到位。代价是机械复杂度与地面打滑——麦克纳姆轮在地毯、碎石上横向滑移严重，实际运动会偏离模型。Ackermann（阿克曼转向）底盘的运动学与停车约束更强，本章不展开建模，统一交给 \cref{ch:mm-wheeled}；它对 base placement 的影响在 \cref{sec:mmk-irm} 末尾点出。

\begin{figure}[ht]
\centering
\begin{tikzpicture}[>=Stealth, scale=1.0]
  % world frame
  \draw[->, gray!70] (-0.4,-0.4) -- (1.1,-0.4) node[right]{$x_w$};
  \draw[->, gray!70] (-0.4,-0.4) -- (-0.4,1.0) node[above]{$y_w$};
  % base body (a rounded chassis)
  \begin{scope}[rotate around={25:(2.3,0.6)}]
    \draw[thick, fill=main!8, draw=main!60!black, rounded corners=3pt]
      (1.7,0.1) rectangle (2.9,1.1);
    \node[font=\small] at (2.3,0.6) {底盘};
    % base frame
    \draw[->, second, thick] (2.3,0.6) -- (3.05,0.6) node[right]{$x_b\,(v)$};
    \draw[->, second, thick] (2.3,0.6) -- (2.3,1.35);
  \end{scope}
  % yaw arrow
  \draw[->, third, thick] (2.55,1.45) arc (60:120:0.55) node[midway, above]{$\omega$};
  % arm: base -> ee
  \draw[very thick, draw=black!70] (2.55,0.95) -- (3.6,1.7) -- (4.45,1.35);
  \fill (4.45,1.35) circle (1.6pt) node[right]{$\{e\}$ 末端};
  % p_be^w vector
  \draw[->, red!70!black, thick, dashed] (2.45,0.75) -- (4.4,1.32);
  \node[red!70!black, font=\small] at (3.35,0.75) {$\mathbf{p}_{be}^w$};
\end{tikzpicture}
\caption{差速底盘 $+$ 臂的速度学。底盘前进 $v$ 沿 $\mathbf{R}_{wb}\mathbf{e}_x$ 给末端纯线速度；偏航 $\omega$ 既给末端角速度 $\omega\mathbf{e}_z$，又通过力臂 $\mathbf{p}_{be}^w$ 给线速度 $\omega(\mathbf{e}_z\times\mathbf{p}_{be}^w)$。末端离底盘越远，偏航对末端线速度的杠杆越大（\cref{eq:mmk-Jb-diff}）。}
\label{fig:mmk-base-arm}
\end{figure}

% ════════════════════════════════════════════════════════════════
\section{约简雅可比：把非完整约束吸收进变量}
\label{sec:mmk-reduced}

\paragraph{差速的 $\mathbf{J}_b$ 与全向的 $\mathbf{J}_b$ 不只是“少一列”。}
\cref{eq:mmk-Jb-diff} 与 \cref{eq:mmk-Jb-omni} 看上去只差一列，容易让人忽略一个本质区别：\textbf{全向底盘的 $\mathbf{J}_b$ 是“把 $\dot{\mathbf{q}}_b$ 抬成末端 twist”，差速底盘的 $\mathbf{J}_b$ 已经把非完整约束\emph{吸收}进去了}。理解这一点，是把非完整底盘正确放进联合规划的关键，也是工程里最容易写错的地方。

\subsection{动机：为什么不能对 \texorpdfstring{$\dot{\mathbf{q}}_b$}{qb-dot} 求伪逆再裁剪}
\label{sec:mmk-naive}

一个自然但错误的做法是：先写“全自由度”几何雅可比
\begin{equation}
  \mathbf{J}_{\text{full}}(\mathbf{q})=\begin{bmatrix}\mathbf{J}_{b,\text{full}}(\mathbf{q})&\mathbf{J}_a(\mathbf{q})\end{bmatrix}\in\mathbb{R}^{6\times(3+n_a)},
  \label{eq:mmk-Jfull}
\end{equation}
把 $\dot{\mathbf{q}}_b=(\dot x_b,\dot y_b,\dot\theta_b)$ 三个分量都当可控量，再对 $\mathbf{J}_{\text{full}}$ 求伪逆得到 $\dot{\mathbf{q}}$。问题在于：差速底盘\emph{根本不能}独立命令车体横向速度 $v_y^B$。伪逆解一般含非零横向分量，丢给底盘执行时，要么被忽略（末端任务无法满足），要么用滑移近似（产生与模型不符的误差）。

\begin{pitfall}[概念误区：把非完整约束当“事后裁剪”]
新手想法：“先按全向求伪逆，再把横向速度分量置零就行。”实际上：把横向分量直接置零\textbf{不是}约束流形上的投影——置零后的 $\dot{\mathbf{q}}_b$ 配上原来的 $\dot{\mathbf{q}}_a$，乘回 $\mathbf{J}_{\text{full}}$ 已不再等于 $\mathbf{V}_d$，末端任务的满足度被破坏。正确做法：在“求解之前”就把可控速度限制在非完整约束允许的子空间内，即使用\textbf{约简雅可比}。这与 \cref{eq:cd-pfaffian} 把 Pfaffian 约束写进速度空间是同一思想。
\end{pitfall}

\subsection{准速度与约简映射}
\label{sec:mmk-quasivel}

回顾差速底盘运动学的矩阵形式（\cref{ch:mm-wheeled}）：
\begin{equation}
  \dot{\mathbf{q}}_b=\mathbf{S}(\mathbf{q}_b)\,\mathbf{u}_b,\qquad
  \mathbf{S}(\mathbf{q}_b)=\begin{bmatrix}\cos\theta_b&0\\\sin\theta_b&0\\0&1\end{bmatrix}\in\mathbb{R}^{3\times2}.
  \label{eq:mmk-S}
\end{equation}
$\mathbf{S}$ 的两列张成的，正是非完整约束 $\mathbf{A}(\mathbf{q}_b)\dot{\mathbf{q}}_b=0$ 的零空间——满足 $\mathbf{A}\mathbf{S}=\mathbf{0}$，其中 $\mathbf{A}=(\sin\theta_b,-\cos\theta_b,0)$ 即“不能瞬时侧移”（\cref{eq:cd-pfaffian} 的具体实例）。把 $\mathbf{u}_b=(v,\omega)$ 称为\textbf{准速度}（quasi-velocity）：它不是任何位形坐标的时间导数（由 \cref{thm:cd-frobenius}，分布对李括号不封闭，不可积，故不存在 $h$ 使 $v=\dot h$），却完整参数化了底盘所有可行的瞬时运动。这正是 \cref{ch:mm-wheeled} 强调“$\boldsymbol{\nu}$ 不是简单的 $\dot{\mathbf{q}}$”的数学根源。

把 $\mathbf{S}$ 推广到全身，定义\textbf{约简选择矩阵}
\begin{equation}
  \bar{\mathbf{S}}(\mathbf{q})=\begin{bmatrix}\mathbf{S}(\mathbf{q}_b)&\mathbf{0}\\\mathbf{0}&\mathbf{I}_{n_a}\end{bmatrix}
  \in\mathbb{R}^{(3+n_a)\times(2+n_a)},\qquad
  \dot{\mathbf{q}}=\bar{\mathbf{S}}(\mathbf{q})\,\boldsymbol{\nu}.
  \label{eq:mmk-Sbar}
\end{equation}

\subsection{约简雅可比}
\label{sec:mmk-Jbar}

\begin{theorem}[约简雅可比]\label{thm:mmk-Jbar}
把 \cref{eq:mmk-Sbar} 代入“全自由度”几何雅可比，得\textbf{约简雅可比}
\begin{equation}
  \boxed{\;\bar{\mathbf{J}}(\mathbf{q})=\mathbf{J}_{\text{full}}(\mathbf{q})\,\bar{\mathbf{S}}(\mathbf{q})\in\mathbb{R}^{6\times(2+n_a)}.\;}
  \label{eq:mmk-Jbar}
\end{equation}
它的底盘部分与 \cref{prop:mmk-Jb-diff} 直接推出的 $\mathbf{J}_b^{\text{diff}}$ \emph{完全一致}。
\end{theorem}
\begin{proof}
只需验证底盘块 $\mathbf{J}_{b,\text{full}}\,\mathbf{S}=\mathbf{J}_b^{\text{diff}}$。取“底座 twist 到末端 twist”的伴随形式 $\mathbf{J}_{b,\text{full}}=\begin{bsmallmatrix}\mathbf{I}_3&-(\mathbf{p}_{be}^w)^\wedge\\\mathbf{0}&\mathbf{I}_3\end{bsmallmatrix}$，而 $\mathbf{S}$ 把 $(v,\omega)$ 嵌成底座 twist $(v\cos\theta_b,v\sin\theta_b,0,0,0,\omega)$，即线速度部分为 $\mathbf{R}_{wb}\mathbf{e}_x\,v$、角速度部分为 $\mathbf{e}_z\,\omega$。代入相乘：线速度行得 $\mathbf{R}_{wb}\mathbf{e}_x\,v-(\mathbf{p}_{be}^w)^\wedge\mathbf{e}_z\,\omega=\mathbf{R}_{wb}\mathbf{e}_x\,v+(\mathbf{e}_z\times\mathbf{p}_{be}^w)\,\omega$，角速度行得 $\mathbf{e}_z\,\omega$，恰为 \cref{eq:mmk-Jb-diff}。$\qed$
\end{proof}

\noindent 这说明 \cref{eq:mmk-Jb-diff} 那个看似“凑”出来的 $\mathbf{J}_b$，其实是 $\mathbf{J}_{\text{full}}\bar{\mathbf{S}}$ 的严格结果——非完整约束已经通过\emph{右乘 $\mathbf{S}$} 被结构性地吸收。

\begin{insight}[约简的本质洞察：约束从“求解器负担”变成“变量定义”]
一旦用 $\bar{\mathbf{J}}$ 替代 $\mathbf{J}_{\text{full}}$，后面所有伪逆、零空间、QP 公式都可\emph{原封不动}套用，而求出的 $\boldsymbol{\nu}$ 天然可行——不需要任何事后裁剪。代价是：求出 $\boldsymbol{\nu}$ 后若要还原 $\dot{\mathbf{q}}_b$，必须再左乘 $\mathbf{S}$；横向速度永远是耦合量 $v\sin\theta_b$，而非独立自由度。一句话：把“约束”从运行时的额外方程，前移成了变量怎么取。
\end{insight}

\subsection{非完整可操作性}
\label{sec:mmk-wnh}

把固定臂的 Yoshikawa 可操作度（\cref{sec:ak-singularity} 的 $w=\sqrt{\det(\mathbf{J}\mathbf{J}^\top)}$）推广到移动操作时，\emph{必须}用约简雅可比而非全自由度雅可比：
\begin{equation}
  w_{\text{nh}}(\mathbf{q})=\sqrt{\det\!\big(\bar{\mathbf{J}}\bar{\mathbf{J}}^\top\big)}.
  \label{eq:mmk-wnh}
\end{equation}
正是这个“约简可操作性”——而非固定臂的 $w(\mathbf{q}_a)$——才正确刻画移动平台在当前位形下产生末端运动的能力。一个直接推论：\textbf{差速底盘的朝向 $\theta_b$ 会进入可操作性}（因为 $\bar{\mathbf{J}}$ 含 $\mathbf{R}_{wb}$），这是固定臂场景没有的现象，也是 \cref{sec:mmk-irm} 的停车位姿必须把 $\theta_b$ 当关键决策量的根本原因。

\begin{proposition}[全向不劣于差速]\label{prop:mmk-omni-ge-diff}
在\emph{同一臂构型}下，全向底盘的约简可操作性恒不低于差速底盘：
\begin{equation}
  w_{\text{nh}}^{\text{omni}}(\mathbf{q})\ \ge\ w_{\text{nh}}^{\text{diff}}(\mathbf{q}).
  \label{eq:mmk-omni-ge}
\end{equation}
\end{proposition}
\begin{proof}
全向底盘的约简雅可比 $\bar{\mathbf{J}}_{\text{omni}}$ 比差速的 $\bar{\mathbf{J}}_{\text{diff}}$ 多一列（横向 $\mathbf{R}_{wb}\mathbf{e}_y$），即 $\bar{\mathbf{J}}_{\text{diff}}$ 的列空间是 $\bar{\mathbf{J}}_{\text{omni}}$ 列空间的子集。于是 $\bar{\mathbf{J}}_{\text{omni}}\bar{\mathbf{J}}_{\text{omni}}^\top=\bar{\mathbf{J}}_{\text{diff}}\bar{\mathbf{J}}_{\text{diff}}^\top+\mathbf{r}\mathbf{r}^\top$，其中 $\mathbf{r}$ 为新增列在末端空间的贡献，$\mathbf{r}\mathbf{r}^\top\succeq\mathbf{0}$。半正定矩阵叠加一个秩一半正定项，特征值逐个不减，故行列式不减，开方得 \cref{eq:mmk-omni-ge}。$\qed$
\end{proof}

\noindent 这把“全向底盘对操作更友好”从定性印象升格为可证不等式：全向的“友好”本质上是它在\emph{任意}位形下都拥有不小于差速的末端运动能力——前提仍是轮子不打滑。

% ════════════════════════════════════════════════════════════════
\section{冗余度与冗余 IK}
\label{sec:mmk-redundancy}

\paragraph{为什么冗余是移动操作的核心概念。}
固定 $6$ 自由度臂给定末端位姿后只有有限解：常见的球腕（三轴交于一点、可位置/姿态解耦，Pieper 准则）臂最多 $8$ 组解（肘/肩/腕组合，见 \cref{ch:arm-kin}），一般 $6R$ 臂（任意几何）最多 $16$ 组解（Lee--Liang 的 $16$ 次多项式）；冗余的处理见 \cref{sec:ak-redundancy}。移动操作平台却有“多余”自由度——底盘 $+$ 臂加起来超过末端任务所需的 $6$ 维。这些多余自由度不是浪费，而是巨大的工程优势：它们让系统在完成末端任务的\emph{同时}，还能满足额外偏好与约束。

\begin{definition}[移动操作的冗余度]\label{def:mmk-redundancy}
设末端任务维度为 $m$（完整位姿任务 $m=6$），可控准速度维度为 $\dim(\boldsymbol{\nu})$。\textbf{冗余度}
\begin{equation}
  r=\dim(\boldsymbol{\nu})-m.
  \label{eq:mmk-r}
\end{equation}
差速底盘 $+$ $7$ 轴臂、$6$ 维任务：$r=(2+7)-6=3$；全向底盘 $+$ $7$ 轴臂：$r=(3+7)-6=4$。
\end{definition}

\begin{note}[冗余计的是准速度，不是位形速度]\label{note:mmk-r-quasivel}
\cref{eq:mmk-r} 用 $\dim(\boldsymbol{\nu})$ 而非 $\dim(\mathbf{q})$。差速底盘虽有 $3$ 个位形坐标，但只贡献 $2$ 个可控方向，故冗余“少 $1$”。这是\emph{速度级}现象，不影响它\emph{最终}能停到任意位姿（长期可达性，由 \cref{thm:cd-frobenius} 的李括号生成保证）——\cref{sec:mmk-irm} 的 base placement 之所以仍能枚举任意朝向 $\theta_b$，正依赖这种长期可达性。差速底盘在“纯横向”微调上比全向“少一个直接手段”，体现为它的约简雅可比在该方向上必须借道偏航 $\omega$ 与臂关节耦合实现。
\end{note}

\subsection{阻尼伪逆与零空间投影}
\label{sec:mmk-dls}

给定期望末端 twist $\mathbf{V}_d$，求联合准速度。最朴素的目标是最小二乘 $\min_{\boldsymbol{\nu}}\|\mathbf{J}\boldsymbol{\nu}-\mathbf{V}_d\|^2$，但在奇异附近 $\mathbf{J}\mathbf{J}^\top$ 病态、解会爆炸。沿用 \cref{thm:ak-dls} 的阻尼最小二乘（DLS），加入阻尼正则：
\begin{equation}
  \boldsymbol{\nu}^\star=\arg\min_{\boldsymbol{\nu}}\ \|\mathbf{J}\boldsymbol{\nu}-\mathbf{V}_d\|^2+\lambda^2\|\boldsymbol{\nu}\|^2
  \;\Longrightarrow\;
  \boldsymbol{\nu}^\star=\mathbf{J}^\top(\mathbf{J}\mathbf{J}^\top+\lambda^2\mathbf{I})^{-1}\mathbf{V}_d.
  \label{eq:mmk-dls}
\end{equation}
这与浮基统一框架的阻尼伪逆 \cref{eq:umm-dls} 是同一表达式。冗余度 $r>0$ 时，解不唯一，可加一个\emph{零空间}项表达次级偏好：
\begin{equation}
  \boxed{\;\boldsymbol{\nu}^\star=\mathbf{J}^{+}\mathbf{V}_d+(\mathbf{I}-\mathbf{J}^{+}\mathbf{J})\,\boldsymbol{\nu}_0,\;}
  \label{eq:mmk-nullspace}
\end{equation}
其中 $\boldsymbol{\nu}_0$ 是“期望速度”，投影 $(\mathbf{I}-\mathbf{J}^{+}\mathbf{J})$ 把它压进 $\mathbf{J}$ 的零空间——故 $\boldsymbol{\nu}_0$ 对末端任务\emph{一阶}无影响（与 \cref{sec:umm-nullspace} 的浮基零空间投影完全同构）。

\begin{insight}[冗余的本质洞察：零空间是工程偏好的入口]
移动操作的冗余不是“多余自由度”，而是把“我更希望谁动”写进数学问题的入口。同一个末端速度，可以由底盘动、臂动或两者共同完成；零空间项 $\boldsymbol{\nu}_0$ 就是在这些等价解里挑一个最合心意的。
\end{insight}

\paragraph{三个常用偏好。}
把 $\boldsymbol{\nu}=(\mathbf{u}_b;\dot{\mathbf{q}}_a)$ 分块，三种典型偏好写法如下（$\mathrm{wrap}(\cdot)$ 表把角度差归一到 $(-\pi,\pi]$）：
\begin{align}
  \text{底盘少动：}&\quad\boldsymbol{\nu}_0=\big(-k_b\,\mathbf{u}_b;\ \mathbf{0}\big), \label{eq:mmk-pref-base}\\
  \text{臂回到舒适位姿：}&\quad\boldsymbol{\nu}_0=\big(\mathbf{0};\ -k_a(\mathbf{q}_a-\mathbf{q}_a^{\text{nom}})\big), \label{eq:mmk-pref-arm}\\
  \text{底盘朝向目标：}&\quad\boldsymbol{\nu}_0=\big(0,\ k_\theta\,\mathrm{wrap}(\theta_{\text{tgt}}-\theta_b);\ \mathbf{0}\big). \label{eq:mmk-pref-yaw}
\end{align}
直觉是：阻尼伪逆 \cref{eq:mmk-dls} 天然取\emph{最小范数}解——底盘一个小幅前进就能整体平移末端，臂关节则要多关节协调才能产生同样位移，故最小范数解倾向多用“效率最高”的底盘自由度。若希望反过来（如狭窄空间里底盘少动），就用 \cref{eq:mmk-pref-base} 把底盘速度往零拉，零空间投影会在\emph{不动末端}的前提下把任务重新分配给臂。

\begin{pitfall}[概念误区：阻尼因子越小越好]
$\lambda\to0$ 时 \cref{eq:mmk-dls} 退化为 Moore--Penrose 伪逆，在雅可比秩亏（奇异位形，\cref{sec:ak-singularity}）时速度发散到无穷。工程上 $\lambda$ 通常取 $10^{-3}$ 到 $10^{-1}$；更精细的做法让 $\lambda$ 随最小奇异值\emph{自适应}——奇异值越小、$\lambda$ 越大，避免过度放大（\cref{thm:ak-dls} 给出这一自适应的标准形式）。
\end{pitfall}

% ════════════════════════════════════════════════════════════════
\section{任务优先级与分层 IK}
\label{sec:mmk-priority}

\paragraph{两级不够用。}
\cref{sec:mmk-dls} 的“主任务 $+$ 零空间次级目标”只有\emph{两级}：末端 twist 是硬的，其余偏好挤在零空间里软性叠加。真实移动操作往往需要\emph{多级严格优先级}，例如“末端位置 $\succ$ 末端姿态 $\succ$ 避关节限位 $\succ$ 底盘少动 $\succ$ 回中”。

\subsection{加权和的根本缺陷}
\label{sec:mmk-weighted}

把 $k$ 个任务写成一个加权最小二乘 $\boldsymbol{\nu}^\star=\arg\min_{\boldsymbol{\nu}}\sum_{i=1}^k w_i\|\mathbf{J}_i\boldsymbol{\nu}-\dot{\mathbf{x}}_i^d\|^2$，当两个任务\emph{冲突}时，解落在由权重决定的折中点上。这有两个问题：(1) 折中比例对权重数值极敏感，标定困难；(2) 没有“硬优先级”——再小的低优先任务也会以 $O(w_i/w_1)$ 量级污染高优先任务。Siciliano 与 Slotine 的\textbf{任务优先级}框架用嵌套零空间投影替代加权，从结构上杜绝低优先级对高优先级的污染。

\subsection{递归零空间投影（等式任务）}
\label{sec:mmk-recursive}

\begin{theorem}[任务优先级的递归零空间投影]\label{thm:mmk-recursive}
设任务按优先级排序 $1\succ2\succ\cdots\succ k$，第 $i$ 个任务有雅可比 $\mathbf{J}_i$ 与期望速度 $\dot{\mathbf{x}}_i^d$。递归构造
\begin{align}
  \boldsymbol{\nu}_1&=\mathbf{J}_1^{+}\dot{\mathbf{x}}_1^d,\qquad \mathbf{N}_1=\mathbf{I}-\mathbf{J}_1^{+}\mathbf{J}_1, \label{eq:mmk-rec-init}\\
  \boldsymbol{\nu}_i&=\boldsymbol{\nu}_{i-1}+\big(\mathbf{J}_i\mathbf{N}_{i-1}\big)^{+}\big(\dot{\mathbf{x}}_i^d-\mathbf{J}_i\boldsymbol{\nu}_{i-1}\big),\quad i=2,\dots,k, \label{eq:mmk-rec-step}\\
  \mathbf{N}_i&=\mathbf{N}_{i-1}-\big(\mathbf{J}_i\mathbf{N}_{i-1}\big)^{+}\big(\mathbf{J}_i\mathbf{N}_{i-1}\big), \label{eq:mmk-rec-N}
\end{align}
其中 $\mathbf{N}_{i-1}$ 是“前 $i-1$ 个任务联合零空间”的投影矩阵。最终联合准速度为 $\boldsymbol{\nu}_k$。
\end{theorem}

\noindent 每一级只在“不破坏所有更高优先级任务”的子空间里求解自己的目标。当某级与高优先级\emph{完全冲突}时，$\mathbf{J}_i\mathbf{N}_{i-1}$ 秩亏，该级在可行子空间内做最小二乘近似，而高优先级\emph{严格不受影响}——这正是加权和做不到的。

\begin{insight}[优先级 = 子空间包含关系]
零空间投影把“优先级”编码成几何包含：第 $i$ 级的解永远活在前 $i-1$ 级的零空间里，故它对高优先级任务的一阶影响\emph{恒等于零}，与它自己的“权重”无关。加权和是“用数值大小排座次”，任务优先级是“用子空间维度排座次”——后者才是真正的硬优先级。
\end{insight}

\subsection{为什么移动操作更需要 HQP}
\label{sec:mmk-hqp}

\cref{thm:mmk-recursive} 只能处理\emph{等式任务}，无法直接表达“关节限位 $\mathbf{q}_{\min}\le\mathbf{q}\le\mathbf{q}_{\max}$”“底盘速度限幅”“避障距离 $\ge d_{\text{safe}}$”这类\emph{不等式}约束。而移动操作几乎每一级都带不等式。Escande、Mansard 与 Wieber 的\textbf{分层二次规划}（Hierarchical Quadratic Programming, HQP）把每一级写成一个带等式/不等式的 QP，按优先级\emph{逐级}求解，并把上一级的最优\emph{残差}作为下一级的硬约束传下去。第 $i$ 级的 QP（概念形式）：
\begin{equation}
\begin{aligned}
  \min_{\boldsymbol{\nu},\,\mathbf{w}_i}\quad& \tfrac12\|\mathbf{w}_i\|^2\\
  \text{s.t.}\quad& \mathbf{J}_i\boldsymbol{\nu}-\dot{\mathbf{x}}_i^d=\mathbf{w}_i &&\text{（本级松弛）}\\
  & \mathbf{A}_i\boldsymbol{\nu}\le\mathbf{b}_i &&\text{（本级不等式）}\\
  & \mathbf{J}_j\boldsymbol{\nu}-\dot{\mathbf{x}}_j^d=\mathbf{w}_j^\star,\ j<i &&\text{（锁定更高优先级的最优残差）}\\
  & \mathbf{A}_j\boldsymbol{\nu}\le\mathbf{b}_j,\ j<i &&\text{（继承更高优先级不等式）}
\end{aligned}
\label{eq:mmk-hqp}
\end{equation}
求得本级最优松弛 $\mathbf{w}_i^\star$ 后冻结，进入下一级。它能在任意层级插入不等式，并在全身规模问题上以控制频率运行——这使它成为人形/移动操作整机控制的主流后端。其思想与足式 WBC（whole-body control）的分层求解一致，区别仅在那里是\emph{动力学级}（力/力矩，统一形式见 \cref{eq:umm-core}），这里是\emph{运动学级}（速度）。

\subsection{典型任务栈}
\label{sec:mmk-stack}

把上述框架落到一台差速底盘 $+$ $7$ 轴臂上，一个常见的任务栈如下：

\begin{table}[H]\centering\small
\caption{移动操作的典型优先级任务栈（用约简雅可比 $\bar{\mathbf{J}}$，非完整约束已在变量层满足）。}
\label{tab:mmk-stack}
\begin{tabular}{cllc}
\toprule
优先级 & 任务 & 类型 & 雅可比 / 约束 \\
\midrule
P0 & 关节/底盘速度限幅、避障距离 & 不等式（最硬） & $\mathbf{A}\boldsymbol{\nu}\le\mathbf{b}$ \\
P1 & 末端位置跟踪 & 等式 & $\bar{\mathbf{J}}_{\text{pos}}\boldsymbol{\nu}=\dot{\mathbf{p}}^d$ \\
P2 & 末端姿态跟踪 & 等式 & $\bar{\mathbf{J}}_{\text{ori}}\boldsymbol{\nu}=\boldsymbol{\omega}^d$ \\
P3 & 远离关节限位 / 提升可操作性 & 等式（梯度方向） & $\nabla_{\boldsymbol{\nu}}(\cdot)$ \\
P4 & 底盘少动 / 臂回中 & 等式（最软） & $\boldsymbol{\nu}\to\boldsymbol{\nu}_0$ \\
\bottomrule
\end{tabular}
\end{table}

\noindent 把避障与限幅放在最高优先级 P0，保证任何情况下安全约束都不会被末端任务“挤掉”——契合“宁可够不到，也不能撞”的工程直觉。

\begin{pitfall}[思维陷阱：层数越多越好]
新手想法：“每个偏好都给一层严格优先级，系统最讲道理。”实际上：层数越多，可行子空间逐级收缩越快，低优先级可能根本没有自由度可用（$\mathbf{N}_{i-1}$ 已退化为零维）；更糟的是\emph{优先级切换}（任务进出栈）会引起速度不连续，导致抖动。正确思维：移动操作通常 $3$--$4$ 级足够（安全 / 末端 / 限位 / 名义），其余偏好用同级加权或零空间软项处理。“多层”本身是有工程代价的。
\end{pitfall}

\subsection{可操作性梯度作为零空间目标}
\label{sec:mmk-wgrad}

\cref{sec:mmk-irm} 用可操作性做停车位姿的\emph{离线}打分；在\emph{在线}速度 IK 里，更常见的是把“提升可操作性、远离奇异”写成一个零空间速度，方向取 $w_{\text{nh}}$ 对位形的梯度：
\begin{equation}
  \boldsymbol{\nu}_0^{(w)}=k_w\,\nabla_{\mathbf{q}}\,w_{\text{nh}}(\mathbf{q}),\qquad
  \frac{\partial w_{\text{nh}}}{\partial q_j}=\tfrac12\,w_{\text{nh}}\cdot\operatorname{tr}\!\Big(\big(\bar{\mathbf{J}}\bar{\mathbf{J}}^\top\big)^{-1}\frac{\partial(\bar{\mathbf{J}}\bar{\mathbf{J}}^\top)}{\partial q_j}\Big).
  \label{eq:mmk-wgrad}
\end{equation}
（系数 $\tfrac12$ 来自 $w_{\text{nh}}=\sqrt{\det(\bar{\mathbf{J}}\bar{\mathbf{J}}^\top)}$ 对 $\sqrt{\cdot}$ 求导，再用 Jacobi 行列式公式 $\partial\det\mathbf{D}=\det\mathbf{D}\cdot\operatorname{tr}(\mathbf{D}^{-1}\partial\mathbf{D})$；等价地也可写成 $\partial w_{\text{nh}}/\partial q_j=w_{\text{nh}}\operatorname{tr}\big((\partial\bar{\mathbf{J}}/\partial q_j)\,\bar{\mathbf{J}}^{+}\big)$，此式中 $\tfrac12$ 已被吸收。）把 $\boldsymbol{\nu}_0^{(w)}$ 投进零空间 $(\mathbf{I}-\bar{\mathbf{J}}^{+}\bar{\mathbf{J}})\boldsymbol{\nu}_0^{(w)}$，系统会在完成末端任务的同时自动朝“更灵巧”的位形漂移、远离奇异。纯反应式（purely-reactive）控制器把这一项与避障、限位一起写成一个 QP，无需全局规划也能让移动平台持续保持高可操作性——这是把停车位姿的“离线选位”思想下沉为“在线连续调整”的代表性做法。

\begin{pitfall}[概念误区：可操作性梯度需要手推解析二阶导]
新手想法：“$\partial(\bar{\mathbf{J}}\bar{\mathbf{J}}^\top)/\partial q_j$ 要手推每个关节的雅可比导数，工程里没法用。”实际上：绝大多数实现用\emph{数值差分}或自动微分得到该梯度，根本不手推。在 $10$ 自由度移动操作上，一次梯度评估只是若干次正运动学，远低于一个控制周期的预算。手推解析式只在追求极致实时性时才必要。
\end{pitfall}

\paragraph{速度级 IK 算法。}
把上述零件组装成一个最常用的“主任务 $+$ 软偏好”单级求解器（多级时换成 \cref{eq:mmk-hqp} 的逐级 QP）：

\begin{algo}[移动操作速度级冗余 IK（单级 $+$ 软偏好）]\label{alg:mmk-velik}
输入：约简雅可比 $\bar{\mathbf{J}}$、期望末端 twist $\mathbf{V}_d$、名义速度 $\boldsymbol{\nu}_0$、阻尼 $\lambda$、底盘/臂正则权重对角阵 $\mathbf{R}$；输出：联合准速度 $\boldsymbol{\nu}^\star$。
\begin{enumerate}
  \item 组装正规方程 $\mathbf{H}=\bar{\mathbf{J}}^\top\bar{\mathbf{J}}+\mathbf{R}+\lambda^2\mathbf{I}$，$\quad\mathbf{g}=\bar{\mathbf{J}}^\top\mathbf{V}_d+\mathbf{R}\,\boldsymbol{\nu}_0$；
  \item 解对称正定系统 $\mathbf{H}\boldsymbol{\nu}^\star=\mathbf{g}$（Cholesky / LDLᵀ）；
  \item （可选）若任一分量越速度限幅，则把违限分量钳到边界、固定为常量，回到第 1 步对剩余自由度重解（饱和处理）；
  \item 还原底盘位形速度 $\dot{\mathbf{q}}_b=\mathbf{S}(\mathbf{q}_b)\mathbf{u}_b$，积分一步 $\mathbf{q}\leftarrow\mathbf{q}\oplus\boldsymbol{\nu}^\star\Delta t$。
\end{enumerate}
该算法对应优化问题 $\min_{\boldsymbol{\nu}}\tfrac12\|\bar{\mathbf{J}}\boldsymbol{\nu}-\mathbf{V}_d\|^2+\tfrac12(\boldsymbol{\nu}-\boldsymbol{\nu}_0)^\top\mathbf{R}(\boldsymbol{\nu}-\boldsymbol{\nu}_0)+\tfrac12\lambda^2\|\boldsymbol{\nu}\|^2$；$\mathbf{R}$ 的底盘对角项调大则求解器少用底盘速度、臂对角项调大则少用关节速度。它只解速度级任务，不含碰撞与关节限位的硬约束——后者须升级为 \cref{eq:mmk-hqp} 的 QP，或交给 \cref{ch:mm-ocp} 的带约束最优控制。
\end{algo}

% ════════════════════════════════════════════════════════════════
\section{停车位姿与反向可达图}
\label{sec:mmk-irm}

\paragraph{底盘停在哪里，决定臂能不能抓。}
前面几节是“边走边动”的连续视角；本节切到“先到位再操作”的离散视角。给定目标抓取位姿 $\mathbf{T}_g^w$，要找底盘位姿 $\mathbf{q}_b$，使存在臂构型 $\mathbf{q}_a$ 满足 $\mathbf{T}_{we}(\mathbf{q}_b,\mathbf{q}_a)=\mathbf{T}_g^w$，同时底盘/臂不碰撞、关节不超限、IK 远离奇异、视野不被遮挡、底盘可由导航到达。逐个候选解 IK 太慢——\textbf{反向可达图}（Inverse Reachability Map, IRM）把这件事压成离线采样 $+$ 在线查表。

\subsection{从可达图到反向可达图}
\label{sec:mmk-rm-irm}

\begin{definition}[可达图 RM]\label{def:mmk-rm}
把\emph{基座系}下的末端位姿空间离散成 $6$D 体素（voxel）：位置 $(x,y,z)$ 按分辨率 $\Delta_p$（典型 $5\,$cm）网格化，姿态按 $\SO(3)$ 上的均匀采样 $\Delta_R$ 离散。对每个体素 $V$，离线随机采样大量臂构型 $\mathbf{q}_a$、正运动学得 $\mathbf{T}_{be}(\mathbf{q}_a)$，落入哪个体素就为其累加一次命中。每个体素存\textbf{可达性质量}
\begin{equation}
  \rho(V)=\max_{\mathbf{q}_a:\,\mathbf{T}_{be}(\mathbf{q}_a)\in V}\ w(\mathbf{q}_a),\qquad
  w(\mathbf{q}_a)=\sqrt{\det\!\big(\mathbf{J}_a\mathbf{J}_a^\top\big)},
  \label{eq:mmk-rho}
\end{equation}
即落入该体素的所有构型中\emph{最高的 Yoshikawa 可操作度}（也可用命中率或最小到限位距离）。$\rho$ 高的体素表示“臂在这里既够得到、又灵巧、远离奇异”。
\end{definition}

RM 回答“给定底座，末端能到哪”；base placement 要的是反问题“给定末端目标，底座该在哪”。关键观察是这对变换\emph{互为逆}：若末端相对底座的位姿是 $\mathbf{T}_{be}$，则\textbf{底座相对末端}的位姿就是
\begin{equation}
  \mathbf{T}_{eb}=\mathbf{T}_{be}^{-1}.
  \label{eq:mmk-Teb}
\end{equation}

\begin{definition}[反向可达图 IRM]\label{def:mmk-irm}
把 RM 中每个体素 $\mathbf{T}_{be}$ 取逆、连同其质量 $\rho$ 一起存进“以末端为原点”的新体素表，即得 IRM。在线给定世界系目标 $\mathbf{T}_g^w$，候选底座世界位姿直接由
\begin{equation}
  \boxed{\;\mathbf{T}_b^w=\mathbf{T}_g^w\,\mathbf{T}_{eb}=\mathbf{T}_g^w\,\mathbf{T}_{be}^{-1}\;}
  \label{eq:mmk-Tbw}
\end{equation}
枚举 IRM 中所有高质量 $\mathbf{T}_{eb}$ 得到——\emph{不需要在线解任何 IK}。再把 $\mathbf{T}_b^w$ 投影到 $\SE(2)$（取 $x,y$ 与绕 $z$ 偏航，忽略高度与滚转俯仰，因底盘只在平面运动）即候选 $\mathbf{q}_b$。
\end{definition}

\begin{figure}[ht]
\centering
\begin{tikzpicture}[>=Stealth, node distance=6mm]
  % RM box
  \node[draw, rounded corners, fill=main!6, draw=main!60!black, align=center, font=\small, inner sep=4pt] (rm)
    {可达图 RM\\“给定底座，末端能到哪”\\体素键 $=\mathbf{T}_{be}$，值 $=\rho$};
  % inverse op
  \node[draw, rounded corners, fill=third!8, draw=third!60!black, right=14mm of rm, align=center, font=\small, inner sep=4pt] (inv)
    {逐体素取逆\\$\mathbf{T}_{eb}=\mathbf{T}_{be}^{-1}$};
  % IRM box
  \node[draw, rounded corners, fill=second!6, draw=second!60!black, right=14mm of inv, align=center, font=\small, inner sep=4pt] (irm)
    {反向可达图 IRM\\“给定末端，底座该在哪”\\体素键 $=\mathbf{T}_{eb}$，值 $=\rho$};
  \draw[->, thick, gray!60!black] (rm)--(inv);
  \draw[->, thick, gray!60!black] (inv)--(irm);
  % online query
  \node[below=9mm of irm, align=center, font=\small] (q)
    {在线：$\mathbf{T}_b^w=\mathbf{T}_g^w\,\mathbf{T}_{eb}$\\$\to$ 投影到 $\SE(2)$ 得候选 $\mathbf{q}_b$};
  \draw[->, thick, second!60!black] (irm)--(q);
\end{tikzpicture}
\caption{从 RM 到 IRM 的取逆与在线查表。离线大量 FK 采样建 RM，逐体素取逆得 IRM；在线一次矩阵乘 $\mathbf{T}_g^w\mathbf{T}_{eb}$ 即枚举候选底座位姿，免在线 IK。}
\label{fig:mmk-irm}
\end{figure}

\begin{insight}[IRM 的本质洞察：把在线 IK 搜索换成离线 FK 采样 $+$ 查表]
正运动学是廉价的单向计算，逆运动学是昂贵的搜索。IRM 用大量离线 FK 采样把 IK 的答案“预烤”进表里，在线只剩一次矩阵求逆与查表。这与图形学里用预计算光照贴图替代实时光线追踪是同一种工程哲学：把贵的搜索搬到离线。
\end{insight}

\begin{pitfall}[概念误区：IRM 可以直接给出可执行停车点]
新手想法：“从 IRM 查到高质量 $\mathbf{T}_{eb}$，转成 $\mathbf{q}_b$ 就能去停。”实际上：IRM 只编码了\emph{臂的运动学可达性}，它不知道现场障碍、底盘能否导航到该点、相机是否被遮挡。正确做法：IRM 给出的是\textbf{按可达质量排序的候选集}，必须再经 \cref{sec:mmk-baseplace} 的多项代价（costmap、IK 复核、视野）过滤——IRM 负责“快速缩小搜索范围”，不负责“最终拍板”。
\end{pitfall}

\subsection{离散化精度与内存的权衡}
\label{sec:mmk-discretize}

体素分辨率是精度与内存的权衡：$\Delta_p$ 太粗，候选停车点与真实最优偏差大、可能 IK 复核全失败；太细，$6$D 体素数量按 $O(\Delta_p^{-3}\Delta_R^{-3})$ 爆炸，离线构建与内存都吃不消。工程上常只对\emph{任务相关}的位姿子集（如桌面高度 $\pm10\,$cm、末端朝下的抓取姿态）建图，把 $6$D 退化成 $3$--$4$D，显著压缩规模。这也解释了为何通用 IRM 少见、而“面向某类抓取任务的定制 IRM”更实用。

\subsection{停车位姿的多项代价}
\label{sec:mmk-baseplace}

IRM 给出候选集后，最终选点要综合多项代价。一个停车候选的代价可写为
\begin{equation}
  J(\mathbf{q}_b,\mathbf{q}_a)=
  w_d\,d_{\text{nav}}(\mathbf{q}_b)
  +w_c\,c_{\text{costmap}}(\mathbf{q}_b)
  +w_i\,c_{\text{ik}}(\mathbf{q}_a)
  +w_l\,c_{\text{limit}}(\mathbf{q}_a)
  +w_v\,c_{\text{view}}(\mathbf{q}_b),
  \label{eq:mmk-baseplace-cost}
\end{equation}
其中导航距离 $d_{\text{nav}}$、障碍代价 $c_{\text{costmap}}$、IK 残差/奇异 $c_{\text{ik}}$、关节接近限位 $c_{\text{limit}}$、相机视角/遮挡 $c_{\text{view}}$ 各占一项。奇异项用可操作度刻画：
\begin{equation}
  c_{\text{ik}}=\frac{1}{w(\mathbf{q}_a)+\epsilon},\qquad w(\mathbf{q}_a)=\sqrt{\det(\mathbf{J}_a\mathbf{J}_a^\top)}.
  \label{eq:mmk-cik}
\end{equation}
按 \cref{eq:mmk-baseplace-cost} 对 IRM 候选打分排序、取最优即停车位姿。

\begin{algo}[停车位姿求解：IRM 查表 $+$ 多代价过滤]\label{alg:mmk-baseplace}
输入：目标抓取位姿 $\mathbf{T}_g^w$、离线 IRM、costmap、IK 求解器；输出：综合代价最低的停车位姿 $\mathbf{q}_b^\star$（及对应臂构型）。
\begin{enumerate}
  \item 从 IRM 取出可达质量 $\rho$ 高的若干 $\mathbf{T}_{eb}$ 候选；
  \item 对每个候选算 $\mathbf{T}_b^w=\mathbf{T}_g^w\mathbf{T}_{eb}$（\cref{eq:mmk-Tbw}），投影到 $\SE(2)$ 得候选 $\mathbf{q}_b$；
  \item 用 costmap 滤掉碰撞/不可导航的候选；
  \item 对幸存候选解一次臂 IK，复核可达、无碰、远离奇异（更新 $c_{\text{ik}},c_{\text{limit}}$）；
  \item 按 \cref{eq:mmk-baseplace-cost} 打分，返回代价最低者。
\end{enumerate}
\end{algo}

\begin{note}[非完整底盘对 base placement 的影响]\label{note:mmk-nonholo-baseplace}
\cref{def:mmk-irm} 把 $\mathbf{T}_b^w$ 投到 $\SE(2)$ 时枚举了任意朝向 $\theta_b$，这\emph{合法}——由 \cref{note:mmk-r-quasivel} 的长期可达性，差速底盘虽瞬时不能侧移，却能停到任意位姿。但 Ackermann（阿克曼）底盘停车约束更强（最小转弯半径、不能原地转向），可达的最终位姿受限，故其 base placement 更难，候选还需经一道“运动学可停性”过滤。这正是 \cref{sec:mmk-wnh} 强调“朝向 $\theta_b$ 进入可操作性”的下游后果：停车位姿不是只挑位置，朝向同样是一等决策量。
\end{note}

\paragraph{承上启下。}
本章把“底盘 $+$ 臂”统一成了联合雅可比与约简映射，并给出冗余 IK、分层求解与停车位姿两条互补路线。它们都停在\emph{运动学/速度级}：要么逐时刻反应（速度 IK），要么静态选点（base placement）。当任务需要\emph{显式地、在一段时间窗内}权衡“底盘怎么走、臂怎么动、何时到位”，并把碰撞、限位、非完整约束作为硬约束统一处理时，就要把本章的联合运动学搬进一个带约束的最优控制问题——这正是 \cref{ch:mm-ocp} 的主题，本章的 $\bar{\mathbf{J}}$、零空间与 $\SE(3)$ 末端误差将在那里成为代价与约束的零件。

% ════════════════════════════════════════════════════════════════
\section*{本章常见误解汇总}
\begin{table}[H]\centering\small
\begin{tabular}{p{0.46\textwidth} p{0.46\textwidth}}
\toprule
\textbf{常见误解} & \textbf{正确理解} \\
\midrule
差速底盘“冗余少 $1$”意味着停不到某些位姿 & 那是\emph{速度级}现象；长期仍可达全 $SE(2)$（\cref{note:mmk-r-quasivel}）\\
先按全向求伪逆、再把横向速度置零即可 & 置零非约束投影，破坏末端任务；须用约简雅可比（\cref{thm:mmk-Jbar}）\\
$\boldsymbol{\nu}$ 就是位形导数 $\dot{\mathbf{q}}$ & $\boldsymbol{\nu}$ 是准速度，$\dot{\mathbf{q}}=\bar{\mathbf{S}}\boldsymbol{\nu}$；横向恒为耦合量 $v\sin\theta_b$ \\
阻尼因子 $\lambda$ 越小越好 & $\lambda\to0$ 在奇异处速度发散；宜随奇异值自适应（\cref{thm:ak-dls}）\\
多级优先级层数越多越讲道理 & 子空间逐级塌缩、切换抖动；$3$--$4$ 级足够（\cref{tab:mmk-stack}）\\
固定臂可操作度 $w(\mathbf{q}_a)$ 足以刻画移动平台 & 须用约简版 $w_{\text{nh}}=\sqrt{\det(\bar{\mathbf{J}}\bar{\mathbf{J}}^\top)}$，含朝向 $\theta_b$（\cref{eq:mmk-wnh}）\\
IRM 查到的高质量点就能直接去停 & 仅编码臂可达性；须再过 costmap/IK/视野（\cref{alg:mmk-baseplace}）\\
$\mathbf{T}_{wb}$ 与 $\mathbf{T}_{bw}$ 可混用 & 方向相反；下标对 $\mathbf{T}_{ab}=\{b\}$ 在 $\{a\}$ 中的位姿，须一路右乘 \\
\bottomrule
\end{tabular}
\end{table}

% ════════════════════════════════════════════════════════════════
\section*{本章小结}

本章把固定臂的微分运动学与冗余 IK 推广到“可移动、可能非完整”的浮动基座，建立了移动操作的运动学地基。主线是：联合构型 $\mathbf{q}=(\mathbf{q}_b;\mathbf{q}_a)\in\SE(2)\times\mathbb{R}^{n_a}$ 与末端 $\mathbf{T}_{we}=\mathbf{T}_{wb}\mathbf{T}_{be}$（\cref{def:mmk-config}）$\to$ 联合雅可比 $\mathbf{V}_{ee}=[\mathbf{J}_b\ \mathbf{J}_a]\boldsymbol{\nu}$，差速 $6\times2$、全向 $6\times3$（\cref{def:mmk-jac,prop:mmk-Jb-diff,prop:mmk-Jb-omni}）$\to$ 约简雅可比 $\bar{\mathbf{J}}=\mathbf{J}_{\text{full}}\bar{\mathbf{S}}$ 把非完整约束吸收进变量定义（\cref{thm:mmk-Jbar}）$\to$ 阻尼伪逆 $+$ 零空间表达工程偏好（\cref{eq:mmk-dls,eq:mmk-nullspace}）$\to$ 任务优先级与 HQP 组织多级目标（\cref{thm:mmk-recursive,eq:mmk-hqp}）$\to$ IRM 与 base placement 把在线 IK 换成离线查表（\cref{def:mmk-irm,alg:mmk-baseplace}）。两个反复出现的核心命题是：差速底盘的“朝向”会进入可操作性（\cref{eq:mmk-wnh}），全向不劣于差速（\cref{prop:mmk-omni-ge-diff}）。

\begin{table}[H]\centering\small
\caption{符号表}
\begin{tabular}{lll}
\toprule
符号 & 含义 & 首次出现 \\
\midrule
$\mathbf{q}=(\mathbf{q}_b;\mathbf{q}_a)$ & 联合构型（底盘 $\SE(2)$ $+$ 臂 $\mathbb{R}^{n_a}$） & \cref{eq:mmk-config} \\
$\mathbf{T}_{wb},\mathbf{T}_{be},\mathbf{T}_{we}$ & 底座/末端相对底座/末端在世界系的位姿 & \cref{eq:mmk-fk} \\
$\mathbf{u}_b$ & 底盘可控输入（差速 $(v,\omega)$ / 全向 $(v_x^B,v_y^B,\omega)$） & \cref{eq:mmk-joint-jac} \\
$\boldsymbol{\nu}=(\mathbf{u}_b;\dot{\mathbf{q}}_a)$ & 可控准速度（非位形导数） & \cref{eq:mmk-joint-jac} \\
$\mathbf{V}_{ee}=[\mathbf{v}_{ee};\boldsymbol{\omega}_{ee}]$ & 末端 twist（线速度在前） & \cref{eq:mmk-joint-jac} \\
$\mathbf{J}=[\mathbf{J}_b\ \mathbf{J}_a]$ & 联合雅可比 & \cref{eq:mmk-joint-jac} \\
$\mathbf{p}_{be}^w$ & 末端相对底座的位置（世界系） & \cref{eq:mmk-Jb-diff} \\
$\mathbf{S},\bar{\mathbf{S}}$ & 底盘 / 全身约简选择矩阵 & \cref{eq:mmk-S,eq:mmk-Sbar} \\
$\bar{\mathbf{J}}=\mathbf{J}_{\text{full}}\bar{\mathbf{S}}$ & 约简雅可比 & \cref{eq:mmk-Jbar} \\
$\mathbf{A}(\mathbf{q}_b)$ & 非完整 Pfaffian 约束行向量，$\mathbf{A}\mathbf{S}=\mathbf{0}$ & \cref{eq:mmk-S} \\
$w_{\text{nh}}$ & 非完整可操作性 $\sqrt{\det(\bar{\mathbf{J}}\bar{\mathbf{J}}^\top)}$ & \cref{eq:mmk-wnh} \\
$r$ & 冗余度 $\dim(\boldsymbol{\nu})-m$ & \cref{eq:mmk-r} \\
$\lambda$ & 阻尼最小二乘阻尼因子 & \cref{eq:mmk-dls} \\
$\boldsymbol{\nu}_0$ & 零空间名义/偏好速度 & \cref{eq:mmk-nullspace} \\
$\mathbf{N}_i$ & 前 $i$ 个任务联合零空间投影 & \cref{eq:mmk-rec-N} \\
$\mathbf{T}_{eb}=\mathbf{T}_{be}^{-1}$ & 底座相对末端的位姿（IRM 键） & \cref{eq:mmk-Teb} \\
$\rho(V)$ & 体素可达性质量 & \cref{eq:mmk-rho} \\
$\mathbf{T}_b^w=\mathbf{T}_g^w\mathbf{T}_{eb}$ & 由目标反查的候选底座世界位姿 & \cref{eq:mmk-Tbw} \\
\bottomrule
\end{tabular}
\end{table}

\begin{table}[H]\centering\small
\caption{知识点总表}
\begin{tabular}{clll}
\toprule
\# & 知识点 & 核心要点 & 难度 \\
\midrule
1 & 联合构型与 FK & $\SE(2)\times\mathbb{R}^{n_a}$；$\mathbf{T}_{we}=\mathbf{T}_{wb}\mathbf{T}_{be}$ & \(\bigstar\bigstar\) \\
2 & 联合雅可比 & 差速 $6\times2$、全向 $6\times3$；偏航给 $\mathbf{e}_z\times\mathbf{p}$ & \(\bigstar\bigstar\bigstar\) \\
3 & 约简雅可比 & $\bar{\mathbf{J}}=\mathbf{J}_{\text{full}}\bar{\mathbf{S}}$ 吸收非完整约束 & \(\bigstar\bigstar\bigstar\) \\
4 & 非完整可操作性 & $w_{\text{nh}}$ 含朝向；全向不劣于差速 & \(\bigstar\bigstar\bigstar\) \\
5 & 冗余 IK & 阻尼伪逆 $+$ 零空间偏好 & \(\bigstar\bigstar\bigstar\) \\
6 & 任务优先级 / HQP & 递归零空间投影；不等式逐级 QP & \(\bigstar\bigstar\bigstar\bigstar\) \\
7 & IRM / base placement & 取逆查表免在线 IK；多代价过滤 & \(\bigstar\bigstar\bigstar\) \\
\bottomrule
\end{tabular}
\end{table}

\section*{延伸阅读}
\begin{itemize}
  \item \textbf{冗余与零空间地基}：Khatib 的操作空间控制（operational space control）——末端任务、冗余度与零空间控制的奠基思想，对应本书 \cref{ch:operational-space}。
  \item \textbf{移动操作协调控制}：Yamamoto 与 Yun 关于 mobile manipulator 的协调控制工作——底盘与臂如何共同贡献末端运动，是 \cref{sec:mmk-jacobian} 联合雅可比的经典出处。
  \item \textbf{任务优先级与分层求解}：Siciliano--Slotine 的任务优先级框架，以及 Escande--Mansard--Wieber 的分层二次规划（HQP），对应 \cref{sec:mmk-priority}。
  \item \textbf{反向可达图}：Vahrenkamp 等关于 reachability/inverse reachability map 的工作——\cref{sec:mmk-irm} 的体素化与取逆查表的来源。
  \item \textbf{统一最优控制落点}：\cref{ch:mm-ocp} 把本章联合运动学放进带约束的统一 OCP；多臂/双臂协同的运动学耦合见 \cref{ch:arm-dualarm-planning} 与 \cref{sec:ma-coopkin}。
\end{itemize}

\section*{故障排查手册}
\begin{enumerate}
  \item \textbf{症状}：导航到位但 IK 失败 / IK 成功却路径碰撞。\textbf{排查}：base placement 是否\emph{只}检查了导航可达而漏查臂可达（\cref{alg:mmk-baseplace} 第 4 步）；停车容差是否宽于实际定位精度；目标物位姿是否在导航后更新。\textbf{对策}：用底盘\emph{实际}最终位姿重解 IK；收紧容差；停车后加视觉伺服微调。
  \item \textbf{症状}：底盘 footprint 安全，臂或负载却撞货架。\textbf{原因}：$2$D footprint 不知道臂伸出。\textbf{对策}：导航阶段收臂、用 $2$D footprint；操作阶段底盘静止、用 $3$D 碰撞检查；连续移动操作用整机包络/距离场统一两套碰撞。
  \item \textbf{症状}：末端速度 IK 在某位形突然发散。\textbf{原因}：$\bar{\mathbf{J}}$ 秩亏（奇异，\cref{sec:ak-singularity}）、$\lambda$ 太小。\textbf{对策}：用 DLS 并让 $\lambda$ 随最小奇异值自适应（\cref{eq:mmk-dls}）；在零空间加 $w_{\text{nh}}$ 梯度项主动远离奇异（\cref{eq:mmk-wgrad}）。
  \item \textbf{症状}：差速底盘执行末端任务时横向“漂”或末端跟不上。\textbf{原因}：误把差速当全向、用 $\mathbf{J}_{\text{full}}$ 伪逆后把横向置零。\textbf{对策}：改用约简雅可比 $\bar{\mathbf{J}}$（\cref{thm:mmk-Jbar}），从变量层满足非完整约束。
  \item \textbf{症状}：抓取前底盘反复小范围微调、任务时间过长。\textbf{原因}：base placement 代价对底盘路径过敏、目标检测抖动、容差小于定位精度。\textbf{对策}：对目标位姿时间滤波；停车位姿加滞回（新候选提升小则保持原候选）；停车后改用臂 $+$ 视觉伺服微调、不再频繁挪底盘。
\end{enumerate}
 % ch:mm-kinematics
\chapter{移动操作的运动学最优控制}
\label{ch:mm-ocp}

% ════════════════════════════════════════════════════════════════
%  本章吸收 Tutorial 05_运动控制/30_复合/130_OCS2_mobile_manipulator（理论部分）。
%  纯工程（OCS2 容器/task.info/C++/ROS2/FCL）按规范跳过；只取运动学 OCP 的
%  建模、代价、约束方法论。SE(3) 误差/松弛障碍/自碰撞梯度去重指向 ch:unified-multimodal-mpc。
% ════════════════════════════════════════════════════════════════

\begin{practice}[前置自测]
开始之前，先试答下列问题。若已能流畅作答，可只速读本章的代价/约束速查表；若卡住，正是本章要为你解决的。
\begin{enumerate}
  \item 一台「差速底盘 $+$ 六轴臂」的移动操作平台，要让末端去抓桌上一个杯子。你会把\emph{底盘的力矩、轮胎摩擦}也写进优化器吗？为什么这台机器人「低速搬运」的场景，往往\emph{不必}用完整刚体动力学？
  \item 状态 $\mathbf x=(x_b,y_b,\theta_b,\mathbf q_a)$、输入 $\mathbf u=(v,\omega,\dot{\mathbf q}_a)$ 的「运动学流图」$\dot{\mathbf x}=\mathbf f(\mathbf x,\mathbf u)$ 长什么样？它对状态、对输入的雅可比为何\emph{极其稀疏}？
  \item 末端跟踪误差用「位置差 $+$ 欧拉角差」会出什么事？为什么要换成 SE(3) 对数？（若忘了，回看 \cref{sec:umm-se3}。）
  \item 一台\emph{差速}底盘不能横向平移。当末端目标在底盘侧方时，优化器会怎样「绕」过去？这是 bug 还是非完整约束的必然？
  \item 关节速度限位、关节位置限位、自碰撞——这三类约束分别该写成对\emph{输入}的边界、对\emph{状态}的边界、还是对\emph{状态}的非线性不等式？
\end{enumerate}
\end{practice}

\section*{本章目标}
学完本章，你应当能够：
\begin{enumerate}
  \item \textbf{论证}为何低速移动操作用\emph{运动学}模型即足够，并能给出「放进全动力学更真但不划算」的反事实分析；
  \item \textbf{写出}差速底盘 $+$ $n$ 轴臂的状态、输入与流图 $\dot{\mathbf x}=\mathbf f(\mathbf x,\mathbf u)$，并手推其线性化矩阵 $\mathbf A=\partial\mathbf f/\partial\mathbf x$、$\mathbf B=\partial\mathbf f/\partial\mathbf u$ 的稀疏结构；
  \item \textbf{把}末端 SE(3) 跟踪代价放进运动学 OCP，给出其一阶梯度与 Gauss--Newton 二阶近似 $\boldsymbol\ell_{xx}\approx\mathbf J_e^\top\mathbf Q\mathbf J_e$；
  \item \textbf{设计}三类约束——关节速度边界（输入）、关节位置松弛障碍（状态）、自碰撞距离（状态非线性）——并解释各自进入 OCP 的恰当形态；
  \item \textbf{厘清}本章 OCP 与上游 base placement（\cref{ch:mm-kinematics}）、分层导航栈的职责边界，知道「谁给初值、谁管全局、谁管局部连续」。
\end{enumerate}

\subsection*{本章知识导航}
本章是一条「\emph{把底盘臂联合运动学，塞进一个有限时域最优控制问题}」的主线。它\emph{不}重新推导 SE(3) 误差与松弛障碍（那是 \cref{ch:unified-multimodal-mpc} 的核心结论），而是把它们\emph{安放}进运动学 OCP 这一具体语境，并补齐运动学模型特有的部分：低维稀疏流图、Gauss--Newton 末端代价、差速非完整约束的后果。

\begin{center}
\small
\begin{tikzpicture}[
  node distance=4.5mm and 7mm, >=Stealth,
  blk/.style={draw, rounded corners, align=center, font=\small, inner sep=3pt, fill=main!6, draw=main!60!black},
  grp/.style={draw, rounded corners, align=center, font=\small, inner sep=3pt, fill=second!6, draw=second!60!black},
  ln/.style={->, thick, gray!60!black}]
\node[blk] (why) {为何运动学足够\\（低速·下层闭环）};
\node[blk, right=of why] (model) {状态/输入/流图\\稀疏线性化};
\node[grp, right=of model] (cost) {末端 SE(3) 代价\\Gauss--Newton};
\node[grp, below=8mm of why] (vlim) {关节/底盘速度\\$=$ 输入边界};
\node[grp, right=of vlim] (qlim) {关节位置\\松弛障碍};
\node[grp, right=of qlim] (col) {自碰撞\\$\hat{\mathbf n}^\top(\mathbf J_{p2}-\mathbf J_{p1})$};
\node[blk, below=8mm of qlim, fill=third!8, draw=third!60!black] (stack) {与 base placement /\\ 分层导航的边界};
\draw[ln] (why)--(model); \draw[ln] (model)--(cost);
\draw[ln] (model.south) -- ++(0,-4mm) -| (vlim.north);
\draw[ln] (vlim)--(qlim); \draw[ln] (qlim)--(col);
\draw[ln] (qlim.south)--(stack.north);
\end{tikzpicture}
\end{center}

\noindent\textbf{推荐路径}：\cref{sec:mmocp-why}（动机 $+$ 反事实）是本章的「为什么」，决定了后面一切简化的合法性；\cref{sec:mmocp-model}（状态/输入/流图/线性化）是任何读者的主干；\cref{sec:mmocp-cost}（末端代价）与 \cref{sec:mmocp-constraint}（三类约束）是把 \cref{ch:unified-multimodal-mpc} 的几何结论\emph{落到}运动学 OCP 的两块工作面；\cref{sec:mmocp-stack}（分层）面向系统集成，说明本章在「全局规划—局部 OCP—底层伺服」三明治里的位置。

\subsection*{前置知识桥接}
\cref{ch:mm-kinematics} 已建立移动操作平台的联合运动学：底盘 SE(2) 位姿 $\mathbf q_b=(x_b,y_b,\theta_b)$、body 与 world 速度的变换、联合雅可比 $\mathbf V_{\text{ee}}=[\mathbf J_b\ \mathbf J_a]\boldsymbol\nu$，以及 base placement（先把底盘停在何处再伸臂）。本章把这套\emph{瞬时}运动学，沿时间展开成一条\emph{轨迹}，并用最优控制为它配上代价与约束。\cref{ch:unified-multimodal-mpc} 则给出了将在此复用的两块几何利器：SE(3) 末端误差（\cref{sec:umm-se3}、\cref{eq:umm-xierr}）与自碰撞 $+$ 松弛障碍（\cref{sec:umm-collision}、\cref{eq:umm-relaxed}）。差速底盘的非完整本质，承自 \cref{ch:constrained-dynamics} 的 Pfaffian 约束（\cref{eq:cd-pfaffian}、\cref{thm:cd-frobenius}）。

\subsection*{如果跳过本章会怎样}
\begin{itemize}
  \item 你会知道「末端误差用 SE(3) 对数」，却不知道把它放进一个\emph{运动学} OCP 后，Hessian 该用完整二阶还是 Gauss--Newton，以及为什么后者更稳；
  \item 你会以为「底盘绕圈」是权重 bug，反复调权重——而真相往往是差速底盘\emph{不能横移}这一非完整约束的必然几何后果；
  \item 在系统集成时，你会分不清这台 OCP 与上游 base placement、全局导航栈各自负责什么，把全局避障的责任错压给一个短时域局部优化器。
\end{itemize}

\begin{note}
\textbf{本章记号与约定.}\;
底盘 SE(2) 位姿 $\mathbf q_b=(x_b,y_b,\theta_b)$，臂关节 $\mathbf q_a\in\mathbb R^{n_a}$；状态 $\mathbf x=(\mathbf q_b,\mathbf q_a)\in\mathbb R^{3+n_a}$。底盘线/角速度 $v,\omega$，臂关节速度 $\dot{\mathbf q}_a$；输入 $\mathbf u=(v,\omega,\dot{\mathbf q}_a)\in\mathbb R^{2+n_a}$。
末端期望位姿 $\mathbf T_d\in\SE(3)$，当前 $\mathbf T=\mathbf T_{\text{ee}}(\mathbf x)$，SE(3) 误差与权重沿用 \cref{sec:umm-se3}。
向量粗体小写、矩阵粗体大写；右扰动 $\mathbf R\,\mathrm{Exp}(\delta\boldsymbol\phi)$；$\se(3)$ 元素 $\boldsymbol\xi=[\boldsymbol\rho;\boldsymbol\phi]$（平移在前）。
本章只用\emph{运动学}模型，输入即速度，无力矩/惯性项——这与 \cref{ch:unified-multimodal-mpc} 的全身动力学 $\mathbf M\dot{\mathbf v}+\mathbf h=\cdots$（\cref{eq:umm-core}）是\emph{两个抽象层次}，何时用哪个正是 \cref{sec:mmocp-why} 的主题。
\end{note}

% ════════════════════════════════════════════════════════════════
\section{为何运动学模型就足够}
\label{sec:mmocp-why}

\paragraph{动机：一台服务机器人去抓杯子。}
设想 Ridgeback 差速底盘驮着一台 UR5 六轴臂，要去够桌上的杯子。整个动作慢条斯理：底盘以零点几米每秒挪动，臂以温和的关节速度伸展。此时若我们老老实实写出完整刚体动力学，把每个电机力矩、每条轮胎的摩擦、抓取重物时基座的微晃都建模进去，优化器固然「更真」，但要付出沉重代价——而这些代价，几乎全被浪费。

\paragraph{核心观察：下层已经闭环。}
关键在于\emph{分层}。现实平台几乎总是：底盘自带速度控制器（你给 $v,\omega$，它内部闭环去跟踪），臂自带关节速度或位置伺服（你给 $\dot{\mathbf q}_a$，它内部闭环）。于是「力矩 $\to$ 加速度 $\to$ 速度」这一段动力学，已经被下层控制器吃掉了。留给上层最优控制的任务，不是「施加多大力」，而是\emph{在哪个方向、走多快}——即生成一条平滑、避障、可达的\textbf{速度轨迹}。这正是运动学模型刻画的对象。

\begin{insight}
\textbf{运动学 OCP 的定位。}\;它\emph{不}追求真实力矩，而是直接把「可执行速度」当决策变量来优化。这一刀切下去，状态维度骤降（无需把速度纳入状态）、流图变得线性化稀疏、求解极快——快到足以\emph{高频重规划}（数十至上百赫兹滚动优化）。低速移动操作、服务机器人恰是这一权衡的最佳受益者。
\end{insight}

\paragraph{运动学模型放弃了什么。}
诚实地讲，运动学模型\emph{无法}直接保证：电机力矩不饱和、轮胎不打滑、抓重物时基座不晃、动态接触力可行（\cref{tab:mmocp-kinpros}）。它把这些动态细节，\emph{委托}给了下层控制器与适度的输入边界（如限制 $v,\omega,\dot{\mathbf q}_a$ 的幅值，间接约束加速度需求）。

\begin{table}[H]\centering\small
\caption{运动学模型在移动操作 OCP 中的得与失。}
\label{tab:mmocp-kinpros}
\begin{tabular}{ll}
\toprule
得（优势） & 失（代价） \\
\midrule
状态维度低、流图稀疏 & 不预测力矩，不保证力矩不饱和 \\
线性化简单、求解快 & 不建模轮胎打滑 \\
适合高频滚动重规划 & 抓重物时基座晃动未建模 \\
易与全局导航/运动规划共存 & 动态接触力可行性无法直接保证 \\
\bottomrule
\end{tabular}
\end{table}

\begin{pitfall}[反事实：硬把全动力学塞进移动操作 OCP]
若坚持用完整多体动力学：状态须加入广义速度（维度近乎翻倍），输入从速度变为力矩，约束须加入驱动极限与接触力可行性。求解问题确实\emph{更真实}，但单步求解从毫秒级涨到数十毫秒，原本 $100\,\mathrm{Hz}$ 的滚动重规划直接超时。而对低速抓取，那些被「补上」的动态细节，本就由下层速度环处理得很好。结论：\textbf{统一塞进 MPC 不是更对，而是不划算}——抽象层次要与任务时间尺度、下层控制结构匹配。这与 \cref{ch:unified-multimodal-mpc} 的全身\emph{动力学} MPC 并不矛盾：那里处理的是腿足接触切换、动态平衡，动力学是任务的本质；这里处理的是低速、下层已闭环，动力学是可委托的细节。
\end{pitfall}

% ════════════════════════════════════════════════════════════════
\section{状态、输入与运动学流图}
\label{sec:mmocp-model}

\paragraph{状态与输入。}
取差速底盘 $+$ $n_a$ 轴臂，状态由底盘平面位姿与臂关节角拼成，输入由底盘线/角速度与臂关节速度拼成：
\begin{equation}
  \mathbf x=\begin{bmatrix}x_b\\ y_b\\ \theta_b\\ \mathbf q_a\end{bmatrix}\in\mathbb R^{3+n_a},
  \qquad
  \mathbf u=\begin{bmatrix}v\\ \omega\\ \dot{\mathbf q}_a\end{bmatrix}\in\mathbb R^{2+n_a}.
  \label{eq:mmocp-xu}
\end{equation}
注意状态\emph{不含}任何速度量——这是运动学模型与全动力学模型（其状态须含 $\dot{\mathbf q}$）的根本分野。

\paragraph{流图。}
底盘是差速驱动，在世界系里只能「沿当前朝向前进」加「原地转向」，不能横移；臂关节速度即输入本身。于是流图
\begin{equation}
  \dot{\mathbf x}=\mathbf f(\mathbf x,\mathbf u)=
  \begin{bmatrix}v\cos\theta_b\\ v\sin\theta_b\\ \omega\\ \dot{\mathbf q}_a\end{bmatrix}.
  \label{eq:mmocp-flow}
\end{equation}
前三行即 SE(2) 底盘 body 速度 $(v,0,\omega)$ 经朝向 $\theta_b$ 旋到世界系（承 \cref{ch:mm-kinematics} 的 body$\leftrightarrow$world 变换）；后 $n_a$ 行是单位映射。

\begin{insight}
\textbf{「$y$ 方向没有 $v_y$」就是非完整约束。}\;\cref{eq:mmocp-flow} 里底盘世界速度 $(\dot x_b,\dot y_b)$ 永远平行于朝向 $(\cos\theta_b,\sin\theta_b)$，即横向速度恒为零。写成 \cref{ch:constrained-dynamics} 的 Pfaffian 形式 $\mathbf A(\mathbf q)\dot{\mathbf q}=\mathbf 0$（\cref{eq:cd-pfaffian}），就是
\begin{equation}
  \underbrace{[\,-\sin\theta_b\quad \cos\theta_b\quad 0\,]}_{\mathbf A(\mathbf q_b)}
  \begin{bmatrix}\dot x_b\\ \dot y_b\\ \dot\theta_b\end{bmatrix}=0
  \quad(\text{底盘不能侧滑}).
  \label{eq:mmocp-pfaffian}
\end{equation}
由 Frobenius 判据（\cref{thm:cd-frobenius}）这是\emph{不可积}的——你无法把它化成对位形的等式约束，它真正限制的是\emph{速度方向}而非可达位形。这一约束的下游后果（「底盘绕圈」）见 \cref{sec:mmocp-stack}。
\end{insight}

\paragraph{线性化的稀疏性。}
SQP/iLQR 类求解器在每个节点需要流图的雅可比 $\mathbf A=\partial\mathbf f/\partial\mathbf x$、$\mathbf B=\partial\mathbf f/\partial\mathbf u$。逐项求导 \cref{eq:mmocp-flow}：

\begin{derivation}[流图雅可比]
对状态求导，唯一非平凡的耦合来自底盘朝向 $\theta_b$ 进入前两行：
\begin{equation}
  \frac{\partial\dot x_b}{\partial\theta_b}=-v\sin\theta_b,\qquad
  \frac{\partial\dot y_b}{\partial\theta_b}=v\cos\theta_b,
  \label{eq:mmocp-Adetail}
\end{equation}
其余状态偏导全为零（臂部分 $\dot{\mathbf q}_a$ 不依赖任何状态）。对输入求导：
\begin{equation}
  \frac{\partial\dot x_b}{\partial v}=\cos\theta_b,\quad
  \frac{\partial\dot y_b}{\partial v}=\sin\theta_b,\quad
  \frac{\partial\dot\theta_b}{\partial\omega}=1,\quad
  \frac{\partial\dot{\mathbf q}_a}{\partial\dot{\mathbf q}_a}=\mathbf I_{n_a}.
  \label{eq:mmocp-Bdetail}
\end{equation}
拼成矩阵（按 $\mathbf x=(x_b,y_b,\theta_b,\mathbf q_a)$、$\mathbf u=(v,\omega,\dot{\mathbf q}_a)$ 排序）：
\begin{equation}
  \mathbf A=\begin{bmatrix}
  0&0&-v\sin\theta_b&\mathbf 0\\
  0&0&\ \ v\cos\theta_b&\mathbf 0\\
  0&0&0&\mathbf 0\\
  \mathbf 0&\mathbf 0&\mathbf 0&\mathbf 0
  \end{bmatrix},
  \qquad
  \mathbf B=\begin{bmatrix}
  \cos\theta_b&0&\mathbf 0\\
  \sin\theta_b&0&\mathbf 0\\
  0&1&\mathbf 0\\
  \mathbf 0&\mathbf 0&\mathbf I_{n_a}
  \end{bmatrix}.
  \label{eq:mmocp-AB}
\end{equation}
$\mathbf A$ 几乎全零（只有两个非零元，且都正比于 $v$）；$\mathbf B$ 是一个 $2\times2$ 朝向块加一个单位块。
\end{derivation}

\begin{insight}
\textbf{稀疏 $=$ 快。}\;\cref{eq:mmocp-AB} 的极度稀疏，正是运动学移动操作 OCP 求解快的结构性原因：每个时间节点的线性化几乎免费，多重射击的带状 KKT 系统里每块都很瘦。这与 \cref{ch:unified-multimodal-mpc} 全身动力学那个稠密的 $\mathbf M(\mathbf q)$ 形成鲜明对照——后者每个节点都要算质量矩阵与其导数。
\end{insight}

% ════════════════════════════════════════════════════════════════
\section{末端 SE(3) 跟踪代价}
\label{sec:mmocp-cost}

\paragraph{为什么不在这里重推 SE(3) 误差。}
末端「去哪」这件事的几何，\cref{sec:umm-se3} 已彻底讲清：欧拉角差有万向锁、不连续、非最短路径三宗罪，正解是 SE(3) 对数误差 $\boldsymbol\xi_{\text{err}}=\mathrm{Log}(\mathbf T_{\text{ee}}^{-1}\mathbf T_d)^\vee$（\cref{eq:umm-xierr}），平移分量还经左雅可比逆与旋转耦合。本节\emph{不}重复这些，只关心一件运动学 OCP 特有的事：把这个误差写成代价后，它对\emph{状态} $\mathbf x$（而非全动力学里的 $\mathbf q$）的导数长什么样，以及 Hessian 怎么取。

\paragraph{末端位姿来自两段链。}
运动学模型里，末端世界位姿是「底盘位姿 $\times$ 臂正运动学」：
\begin{equation}
  \mathbf T_{\text{ee}}(\mathbf x)=\mathbf T_{wb}(x_b,y_b,\theta_b)\,\mathbf T_{be}(\mathbf q_a),
  \label{eq:mmocp-Twe}
\end{equation}
其中 $\mathbf T_{wb}$ 由 SE(2) 底盘位姿提升为平面内的 SE(3)，$\mathbf T_{be}$ 是臂从其基座到末端 frame 的正运动学。误差与代价：
\begin{equation}
  \mathbf e(\mathbf x)=\mathrm{Log}\big(\mathbf T_d^{-1}\mathbf T_{\text{ee}}(\mathbf x)\big)^\vee\in\mathbb R^6,
  \qquad
  \ell_{\text{ee}}(\mathbf x)=\tfrac12\,\mathbf e(\mathbf x)^\top\mathbf Q_{\text{ee}}\,\mathbf e(\mathbf x),
  \label{eq:mmocp-ee}
\end{equation}
权重 $\mathbf Q_{\text{ee}}=\mathrm{diag}(q_x,q_y,q_z,q_{\text{roll}},q_{\text{pitch}},q_{\text{yaw}})$ 按任务裁剪（抓取重位姿、指向只重姿态、到达只重位置），其阻抗解释见 \cref{sec:umm-se3}。位置与姿态量纲不同，权重不可盲目取等。

\begin{note}
\textbf{误差方向一致即可。}\;\cref{eq:mmocp-ee} 取 $\mathbf T_d^{-1}\mathbf T_{\text{ee}}$（当前拉回期望系），与 \cref{eq:umm-xierr} 的 $\mathbf T_{\text{ee}}^{-1}\mathbf T_{\text{ref}}$ 互为逆，小误差下近似差一个符号、大姿态误差下还差一个伴随的坐标表达。重点不是「选左还是右」，而是\textbf{整个实现自始至终用同一个方向}——误差定义、雅可比、权重必须配套。本书右扰动约定下，对 $\mathbf q$ 的求导按 \cref{eq:umm-Jerr} 取用，工业 CLIK 实现亦然。
\end{note}

\paragraph{代价对状态的梯度与 Hessian。}

\begin{derivation}[运动学末端代价的一阶/二阶展开]
记 $\mathbf J_e=\partial\mathbf e/\partial\mathbf x\in\mathbb R^{6\times(3+n_a)}$ 为误差对状态的雅可比。它由两部分链接而成：误差对末端位姿的对数雅可比（\cref{eq:umm-Jerr} 的 $\mathbf J_{\log}$），再乘以末端位姿对状态的雅可比——后者正是 \cref{ch:mm-kinematics} 的联合雅可比 $[\mathbf J_b\ \mathbf J_a]$（底盘块 $+$ 臂块）。对二次型代价 $\ell_{\text{ee}}=\tfrac12\mathbf e^\top\mathbf Q\mathbf e$，链式法则给一阶梯度
\begin{equation}
  \boldsymbol\ell_x=\mathbf J_e^\top\mathbf Q\,\mathbf e\ \in\mathbb R^{3+n_a}.
  \label{eq:mmocp-lx}
\end{equation}
精确 Hessian 含误差函数的二阶导（$\partial^2\mathbf e/\partial\mathbf x^2$ 的张量缩并），既贵又不保号。\textbf{Gauss--Newton 近似}丢弃这一项：
\begin{equation}
  \boldsymbol\ell_{xx}\approx\mathbf J_e^\top\mathbf Q\,\mathbf J_e\ \succeq\mathbf 0.
  \label{eq:mmocp-lxx}
\end{equation}
当残差 $\mathbf e$ 较小时被丢弃项可忽略，近似足够好；更重要的是 \cref{eq:mmocp-lxx} 天然\emph{半正定}，给 SQP 的二次子问题提供良态 Hessian。
\end{derivation}

\begin{pitfall}[为何不用完整 Hessian]
完整 Hessian 在残差大时可能\emph{不定}，使 SQP 子问题非凸、步长发散。Gauss--Newton 牺牲了远离目标时的二阶精度，换来全程半正定与数值稳定——这在滚动 MPC（每周期只迭代一两步）里尤其值得。代价是：当末端离目标极远时收敛偏慢，可由「随时间推进的末端参考轨迹」（而非一步到位的远目标）缓解。
\end{pitfall}

\paragraph{输入正则与关节正则。}
末端代价之外，运动学 OCP 还需两类「软」项。输入正则
\begin{equation}
  \ell_u(\mathbf u)=\tfrac12(\mathbf u-\mathbf u_{\text{ref}})^\top\mathbf R(\mathbf u-\mathbf u_{\text{ref}}),
  \qquad \mathbf R=\mathrm{diag}(R_v,R_\omega,R_{\dot q_1},\dots),
  \label{eq:mmocp-lu}
\end{equation}
其对角元编码\emph{运动偏好}：$R_v$ 太小底盘爱前后蹭，$R_\omega$ 太小底盘爱原地转，$R_{\dot q}$ 太小臂动作激进，全体太大则末端跟踪迟钝。想「先动臂、少动底盘」就抬高 $R_v,R_\omega$。关节正则 $\ell_q=\tfrac12(\mathbf q_a-\mathbf q_a^{\text{nom}})^\top\mathbf Q_a(\mathbf q_a-\mathbf q_a^{\text{nom}})$ 不是主任务，而是用\emph{冗余零空间}把臂拉向舒适构型——远离限位、避免肘部怪姿、减少多解跳变（冗余与零空间的机理见 \cref{sec:ak-redundancy}）。

\begin{insight}
\textbf{末端抖动先加输入权重，别加末端权重。}\;直觉上末端抖动似乎该「抓得更紧」（抬高 $\mathbf Q_{\text{ee}}$），但这往往\emph{加剧}抖动——末端在每个中间时刻死贴参考，把底盘与臂的自由度榨干。正确做法常是先抬输入权重 $\mathbf R$ 让轨迹更平滑（呼应 \cref{eq:mmocp-lu} 的平滑作用），或平滑参考轨迹本身。
\end{insight}

% ════════════════════════════════════════════════════════════════
\section{三类约束的恰当形态}
\label{sec:mmocp-constraint}

约束是移动操作 OCP 的安全与可行性所在。关键认知是：\emph{同一个物理限制，依它依赖输入还是状态、是「不理想」还是「不允许」，进入 OCP 的形态截然不同}。本节按这条线索梳理三类约束。

\paragraph{约束一：速度限位 $=$ 输入边界。}
关节速度 $\dot{\mathbf q}_a$ 与底盘速度 $v,\omega$ 都\emph{直接是输入分量}（\cref{eq:mmocp-xu}），故它们的限位是最干净的形态——对输入的箱式边界：
\begin{equation}
  \dot{\mathbf q}_{\min}\le\dot{\mathbf q}_a\le\dot{\mathbf q}_{\max},\qquad
  v_{\min}\le v\le v_{\max},\qquad
  \omega_{\min}\le\omega\le\omega_{\max}.
  \label{eq:mmocp-vlim}
\end{equation}

\begin{pitfall}[别漏掉底盘那两维]
输入向量是 $[v,\omega,\dot{\mathbf q}_a]$，速度边界数组的长度必须等于\emph{整个}输入维 $2+n_a$，前两维是底盘线/角速度。常见 bug 是只给臂关节填了上下界、漏了底盘两维：优化器于是认为底盘速度无界，仿真里底盘「窜出去」。自检：边界数组长度 $\overset{?}{=}$ 输入维。
\end{pitfall}

\paragraph{约束二：关节位置限位 $=$ 状态松弛障碍。}
关节位置 $\mathbf q_a$ 是\emph{状态}分量。仅限速度不限位置，长时间运行会缓慢漂到限位。位置约束 $\mathbf q_{\min}\le\mathbf q_a\le\mathbf q_{\max}$ 是真正的安全边界（「不允许」越界），宜用\emph{松弛障碍}软化——其完整定义、对数段在 $h\le\delta$ 处的二次延拓、以及为何要这样延拓（纯对数障碍在约束已被违反时返回 NaN 致求解崩溃），均见 \cref{eq:umm-relaxed} 与 \cref{ch:unified-multimodal-mpc} 的推导，此处不赘。设单侧约束 $h(\mathbf x)=q_{a,i}-q_{\min,i}\ge0$，惩罚项即 $B_\mu(h)$（\cref{eq:umm-relaxed}）。

\begin{note}
\textbf{松弛障碍的工程序。}\;调参常先用较大的 $\mu,\delta$（如 $\mu=0.1,\delta=0.01$）确保求解器能收敛，再逐步收紧使约束更贴边界。直接上小值，初值若违约则求解器可能找不到可行方向。「为什么需要在已违约时仍有有限梯度」这一动机，正是松弛障碍区别于纯内点障碍的要害（\cref{ch:unified-multimodal-mpc}）。
\end{note}

\paragraph{约束三：自碰撞 $=$ 状态非线性不等式。}
臂在底盘上方挥舞，可能撞自己的底盘、夹爪可能磕另一连杆。对每对可能碰撞的连杆 $(i,j)$，安全约束
\begin{equation}
  h_{ij}(\mathbf x)=d_{ij}(\mathbf q)-d_{\min}\ge0,
  \label{eq:mmocp-hcol}
\end{equation}
$d_{ij}$ 是两几何体最小距离，$d_{\min}$ 是安全裕度。其\emph{距离梯度}是把约束接进 SQP 的关键，已在 \cref{sec:umm-collision} 推得解析式（沿用该处约定：最近点 $\mathbf p_i\in\mathrm{link}_i$、$\mathbf p_j\in\mathrm{link}_j$，单位法向 $\hat{\mathbf n}=(\mathbf p_j-\mathbf p_i)/\|\mathbf p_j-\mathbf p_i\|$ 由 $i$ 指向 $j$）
\begin{equation}
  \nabla_{\mathbf q}\,d_{ij}=\hat{\mathbf n}^\top\big(\mathbf J_{p_j}-\mathbf J_{p_i}\big),
  \label{eq:mmocp-ddq}
\end{equation}
$\mathbf J_{p_i},\mathbf J_{p_j}$ 为两最近点的位置雅可比（几何来源、最近点到 frame 的平移见 \cref{sec:umm-collision}）。法向方向与雅可比之差的次序须\emph{配套}：$\hat{\mathbf n}$ 指向 $i\!\to\!j$ 时配 $(\mathbf J_{p_j}-\mathbf J_{p_i})$，二者同时反号才仍得正确的「距离增大方向」。约束 \cref{eq:mmocp-hcol} 同样以松弛障碍软化进 OCP。

\begin{pitfall}[碰撞对的两难：太少漏撞、太多超时]
自碰撞约束有两个工程陷阱。其一，\textbf{忘记排除相邻连杆}：相邻连杆装配上本就「贴着」，$d_{ij}<d_{\min}$ 恒成立，松弛障碍持续输出巨大代价，把整个 OCP 带偏（现象：一启动就「求解成功但行为诡异」，末端怎么都到不了）。须用机器人描述里的相邻关系排除天然相邻对，只留真正可能相撞的远端对（如夹爪 vs 底盘、肘部 vs 底盘）。其二，\textbf{碰撞对太多}：约束维度 $=$ 碰撞对数，每对每节点要做一次最近距离查询加一次雅可比；把对数从几个加到几十个，求解时间可能翻几倍，$100\,\mathrm{Hz}$ MPC 直接超时。实践只保留几何上真有可能相撞的少数对（常 $3$--$10$ 对），这是精度与实时性的折中，而非偷懒。
\end{pitfall}

\begin{note}
\textbf{距离函数的尖角。}\;自碰撞的难点\emph{不}在「距离公式不会写」，而在距离函数在\emph{穿透}与\emph{最近特征切换}（面-面跳到边-边）时不光滑，最近点法向 $\hat{\mathbf n}$ 会突变，使解析雅可比跳变、SQP 步长在碰撞临界处横跳。这类似 $|x|$ 在原点不可导。松弛障碍的二次延拓把约束推到「还没接触」就发力，意在让优化器\emph{永远走不到那个尖角}。
\end{note}

\paragraph{软约束谱：本章为何全用软约束。}
速度、位置、碰撞三类约束在 OCP 中可硬可软（谱见 \cref{tab:mmocp-constr}）。低速移动操作的在线 MPC 几乎全选\emph{软}：软约束最坏只是被违反但代价有限、几乎总能求出解；硬约束在初值不可行时直接判 infeasible，使 MPC 当周期无输出。工程常见路线是先用二次惩罚把跟踪调通，再把安全相关项换成松弛障碍，最后仅在必须精确处才考虑硬约束。

\begin{table}[H]\centering\small
\caption{移动操作运动学 OCP 的三类约束及其恰当形态。}
\label{tab:mmocp-constr}
\begin{tabular}{llll}
\toprule
约束 & 依赖 & 约束函数 & 恰当形态 \\
\midrule
关节/底盘速度限位 & 输入 $\mathbf u$ & $\dot{\mathbf q}_a,v,\omega$ 箱界 & 输入边界（\cref{eq:mmocp-vlim}） \\
关节位置限位 & 状态 $\mathbf x$ & $q_{a}-q_{\min}\ge0$ 等 & 状态松弛障碍（\cref{eq:umm-relaxed}） \\
自碰撞 & 状态 $\mathbf x$ & $d_{ij}-d_{\min}\ge0$ & 状态松弛障碍 $+$ 解析梯度（\cref{eq:mmocp-ddq}） \\
\bottomrule
\end{tabular}
\end{table}

% ════════════════════════════════════════════════════════════════
\section{与 base placement、分层导航的边界}
\label{sec:mmocp-stack}

\paragraph{这台 OCP 在系统里的位置。}
本章的运动学 OCP 是一个\emph{短时域局部优化器}：它在数秒的预测窗口内，生成一条平滑、避自碰撞、可达的速度轨迹，并高频滚动。它\emph{不}是全局规划器，也\emph{不}是底层伺服。一个典型的分层是「全局运动规划 $\to$ 本章局部 OCP $\to$ 底盘/臂速度环」三明治：全局层管「绕开房间里的大障碍、给出粗路径」，本章局部层管「沿粗路径连续重规划、实时微避障、生成可执行速度」，底层管「把 $v,\omega,\dot{\mathbf q}_a$ 真正跟踪出来」。三者在\emph{时间尺度}上互补，不是替代关系。

\paragraph{承 base placement：谁给初值。}
\cref{ch:mm-kinematics} 的 base placement 回答的是「底盘该停在哪、以什么朝向，臂才够得舒服」——一个\emph{准静态}的可达性/可操作性问题。它为本章 OCP 提供两样东西：一个好的\emph{初始猜测}（底盘大致目标位姿），以及底盘朝向正则的参考 $\theta_{\text{ref}}$。这正是治「底盘绕圈」的良方。

\begin{pitfall}[「底盘绕圈」多半不是 bug]
现象：差速底盘不直奔目标，反而绕大圈。许多人误判为权重 bug，疯狂调参。真相常是 \cref{eq:mmocp-pfaffian} 的非完整约束：差速底盘\emph{不能横移}，当末端目标在底盘侧方、且姿态权重又逼它对准某朝向时，优化器只能用「转向 $+$ 前进」的组合去逼近，外观就是绕圈。对策不是无脑调权重，而是：用 base placement 给更好的初始底盘位姿、适度抬高角速度权重 $R_\omega$、降低不必要的姿态权重、或给底盘一个合理的朝向参考 $\theta_{\text{ref}}$。把几何成因（非完整 $+$ 目标在侧方）认出来，比盲调省十倍时间。
\end{pitfall}

\paragraph{与全局层的两种配合。}
本章 OCP 与上游全局规划器有两种常见配合：其一，全局层生成粗路径，本章 OCP \emph{跟踪}该路径并实时避障、输出速度；其二，本章 OCP 直接优化短时域末端目标，全局层\emph{仅在}局部 OCP 陷入不可行或局部极小时才介入做全局重规划。前者把全局层当「领航」，后者把它当「兜底」。无论哪种，本章 OCP 都\emph{不}该独自承担全局避障——一个数秒预测窗口的局部优化器，看不见房间另一头的墙。把全局避障责任错压给它，是移动操作系统集成的高频设计错误。

\begin{insight}
\textbf{运行项与终端项分工。}\;滚动优化里，每个中间节点的末端代价（运行项）决定「路上贴不贴参考」，horizon 末端那一个（终端项）决定「最终收不收敛到目标」。终端末端权重通常设得比运行项更高，给优化器一个明确的「终点锚」；只有运行项、没有终端项时，有限窗口下优化器可能选「全程缓慢逼近、末端仍差一截」的低代价轨迹。这与本章把末端代价随时间展开成轨迹的视角一脉相承。
\end{insight}

% ════════════════════════════════════════════════════════════════
\section*{本章小结}
\addcontentsline{toc}{section}{本章小结}

本章把 \cref{ch:mm-kinematics} 的底盘臂联合运动学，沿时间展开成一个有限时域\emph{运动学}最优控制问题：因低速 $+$ 下层速度环已闭环，直接以可执行速度为决策变量即足够（\cref{sec:mmocp-why}）；状态/输入低维、流图 \cref{eq:mmocp-flow} 与其线性化 \cref{eq:mmocp-AB} 极度稀疏，故求解快、可高频重规划。末端「去哪」复用 \cref{sec:umm-se3} 的 SE(3) 对数误差，本章补齐其对状态的梯度与 Gauss--Newton 半正定 Hessian（\cref{eq:mmocp-lxx}）。三类约束按「依赖输入/状态、不理想/不允许」各取恰当形态：速度限位入输入边界，关节位置与自碰撞入状态松弛障碍（复用 \cref{eq:umm-relaxed} 与 \cref{sec:umm-collision} 的几何）。最后厘清本章 OCP 作为短时域局部优化器，与上游 base placement（给初值/朝向参考）、全局导航层（管全局避障）的职责边界，并把「底盘绕圈」还原为非完整约束的几何后果。核心结论汇于 \cref{tab:mmocp-summary}。

\begin{table}[H]\centering\small
\caption{本章核心公式速查。}
\label{tab:mmocp-summary}
\begin{tabular}{lll}
\toprule
对象 & 核心式 & 出处 \\
\midrule
状态/输入 & $\mathbf x=(x_b,y_b,\theta_b,\mathbf q_a)$，$\mathbf u=(v,\omega,\dot{\mathbf q}_a)$ & \cref{eq:mmocp-xu} \\
运动学流图 & $\dot{\mathbf x}=(v\cos\theta_b,\,v\sin\theta_b,\,\omega,\,\dot{\mathbf q}_a)$ & \cref{eq:mmocp-flow} \\
非完整约束 & $[-\sin\theta_b\ \cos\theta_b\ 0]\,\dot{\mathbf q}_b=0$（不侧滑） & \cref{eq:mmocp-pfaffian} \\
流图雅可比 & $\mathbf A,\mathbf B$ 极稀疏（$\mathbf A$ 仅两元 $\propto v$） & \cref{eq:mmocp-AB} \\
末端代价 & $\ell_{\text{ee}}=\tfrac12\mathbf e^\top\mathbf Q_{\text{ee}}\mathbf e$，$\mathbf e=\mathrm{Log}(\mathbf T_d^{-1}\mathbf T_{\text{ee}})^\vee$ & \cref{eq:mmocp-ee} \\
代价梯度/Hessian & $\boldsymbol\ell_x=\mathbf J_e^\top\mathbf Q\mathbf e$，$\boldsymbol\ell_{xx}\approx\mathbf J_e^\top\mathbf Q\mathbf J_e\succeq0$ & \cref{eq:mmocp-lx,eq:mmocp-lxx} \\
速度限位 & $\dot{\mathbf q}_{\min}\le\dot{\mathbf q}_a\le\dot{\mathbf q}_{\max}$ 等（输入边界） & \cref{eq:mmocp-vlim} \\
自碰撞 & $d_{ij}-d_{\min}\ge0$，$\nabla d=\hat{\mathbf n}^\top(\mathbf J_{p_j}-\mathbf J_{p_i})$ & \cref{eq:mmocp-hcol,eq:mmocp-ddq} \\
\bottomrule
\end{tabular}
\end{table}

\section*{符号表}
\addcontentsline{toc}{section}{符号表}

\begin{table}[H]\centering\small
\begin{tabular}{ll}
\toprule
符号 & 含义 \\
\midrule
$\mathbf x=(x_b,y_b,\theta_b,\mathbf q_a)$ & 状态：底盘 SE(2) 位姿 $+$ 臂关节角，$\in\mathbb R^{3+n_a}$ \\
$\mathbf u=(v,\omega,\dot{\mathbf q}_a)$ & 输入：底盘线/角速度 $+$ 臂关节速度，$\in\mathbb R^{2+n_a}$ \\
$\mathbf f(\mathbf x,\mathbf u)$ & 运动学流图（\cref{eq:mmocp-flow}） \\
$\mathbf A,\mathbf B$ & 流图对状态/输入的雅可比（\cref{eq:mmocp-AB}） \\
$\mathbf T_{\text{ee}}(\mathbf x),\mathbf T_d$ & 当前/期望末端位姿 $\in\SE(3)$ \\
$\mathbf e(\mathbf x)$ & SE(3) 末端误差 $\mathrm{Log}(\mathbf T_d^{-1}\mathbf T_{\text{ee}})^\vee\in\mathbb R^6$ \\
$\mathbf J_e$ & 误差对状态的雅可比 $\partial\mathbf e/\partial\mathbf x$ \\
$\mathbf Q_{\text{ee}},\mathbf R,\mathbf Q_a$ & 末端/输入/关节正则权重 \\
$\boldsymbol\ell_x,\boldsymbol\ell_{xx}$ & 代价对状态的一阶梯度 / Gauss--Newton Hessian \\
$h_{ij},d_{ij},d_{\min}$ & 自碰撞约束值 / 连杆对距离 / 安全裕度 \\
$\hat{\mathbf n},\mathbf J_{p_i}$ & 最近点连线单位法向 / 最近点位置雅可比 \\
$\mathbf A(\mathbf q_b)$ & 底盘非完整约束的 Pfaffian 行（\cref{eq:mmocp-pfaffian}） \\
$\theta_{\text{ref}}$ & 底盘朝向参考（来自 base placement，\cref{ch:mm-kinematics}） \\
\bottomrule
\end{tabular}
\end{table}

\section*{本章与相邻章节的关系}
\addcontentsline{toc}{section}{本章与相邻章节的关系}

本章向上承接 \cref{ch:mm-kinematics}（底盘臂联合运动学、联合雅可比、base placement）与 \cref{ch:unified-multimodal-mpc}（SE(3) 误差几何 \cref{sec:umm-se3}、自碰撞 $+$ 松弛障碍 \cref{sec:umm-collision}），并以 \cref{ch:constrained-dynamics}（Pfaffian 非完整约束）、\cref{sec:ak-redundancy}（冗余零空间）为横向支柱。它与 \cref{ch:unified-multimodal-mpc} 的全身\emph{动力学} MPC 互为对照：那里动力学是任务本质（接触切换、动态平衡），这里动力学可委托给下层（低速、速度环已闭环）——同一套 SE(3) 误差与松弛障碍工具，在两个抽象层次上各得其所。

\begin{practice}[本章练习]
\begin{enumerate}
  \item \textbf{（流图线性化）} 取 $n_a=2$。按 \cref{eq:mmocp-flow} 写出完整流图，手推 $\mathbf A=\partial\mathbf f/\partial\mathbf x$ 与 $\mathbf B=\partial\mathbf f/\partial\mathbf u$（应与 \cref{eq:mmocp-AB} 的结构一致），并指出 $\mathbf A$ 中两个非零元为何都正比于 $v$。
  \item \textbf{（Gauss--Newton 的边界）} 解释为何残差 $\mathbf e$ 大时 \cref{eq:mmocp-lxx} 偏离真 Hessian，以及为何这在滚动 MPC 里仍可接受。设计一个「随时间推进的末端参考」如何缓解远目标收敛慢。
  \item \textbf{（非完整后果）} 设末端目标恰在底盘正侧方 $1\,\mathrm{m}$ 处。结合 \cref{eq:mmocp-pfaffian} 说明：为什么差速底盘不能直接平移过去？给出两种\emph{不靠盲调权重}的对策。
  \item \textbf{（约束归类）} 现要新增一项「底盘朝向正则」$\ell_b=\tfrac12 w_\theta\,\mathrm{wrap}(\theta_b-\theta_{\text{ref}})^2$。它依赖状态还是输入？是「不理想」还是「不允许」？据此判断它该作代价还是约束、用哪种惩罚，并说明 $\theta_{\text{ref}}$ 从哪来（\cref{ch:mm-kinematics}）。
  \item \textbf{（碰撞对设计）} 对一台「差速底盘 $+$ 六轴臂 $+$ 夹爪」，列出你会保留的 $3$--$5$ 个自碰撞对，并说明为何\emph{不}把所有连杆两两组合。结合 \cref{eq:mmocp-ddq} 解释每加一对的求解开销。
\end{enumerate}
\end{practice}
        % ch:mm-ocp
\chapter{操作技能接口：从高层末端意图到全身可执行动作}\label{ch:mm-skill-interface}

% ════════════════════════════════════════════════════════════════
%  吸收 Tutorial 05_运动控制/30_复合/50_操作技能接口（理论与方法论部分）。
%  主线：把「高层操作策略 → 末端 SE(3) 意图 → 全身控制器 → 关节动作」之间的
%  接口当成一个可推导、可调试、可迁移的工程对象，而非一段消息格式。
%  记号承 \cref{ch:lie}（Exp/Log/SO3/SE3 右扰动）、\cref{ch:unified-multimodal-mpc}
%  （统一动力学与 SE(3) 末端代价）、\cref{ch:operational-space}（笛卡尔阻抗）。
% ════════════════════════════════════════════════════════════════

\begin{practice}[前置自测]
开始之前，先试答下列问题。若已能流畅作答，可只速读本章的两张接口表与限速公式；若卡住，正是本章要为你解决的。
\begin{enumerate}
  \item 一个视觉操作策略直接输出整机 $19$ 维关节目标，换到臂长、默认姿态、夹爪零位略有差异的另一台机器人后，同一组关节角还对应同一个抓取动作吗？为什么？
  \item 末端世界位姿满足 ${}^{W}\mathbf{T}_{ee}={}^{W}\mathbf{T}_{B}\,{}^{B}\mathbf{T}_{ee}$。若高层下发的是\emph{机体系}末端目标 ${}^{B}\mathbf{T}_{ee}^{\mathrm{cmd}}$，当四足行走使基座姿态 ${}^{W}\mathbf{R}_{B}$ 抖动时，杯子附近的末端目标会发生什么？
  \item 把两个姿态做“欧拉角逐分量相减”当成姿态误差，会在什么时候出问题？正确的 $\SO(3)$ 误差该怎么写？
  \item 末端目标在一个控制周期内跳变了 $30$ 厘米，为什么不能让低层全身控制器“忠实跟踪”？谁该负责把它裁掉？
  \item 机械臂逆运动学（IK）求出了到目标的一组关节解，是否就意味着四足或人形可以安全执行这个末端目标？
\end{enumerate}
\end{practice}

\section*{本章目标}
学完本章，你应当能够：
\begin{enumerate}
  \item \textbf{论证}为何高层操作策略应输出\emph{与本体无关}的末端意图（末端 $\SE(3)$ 位姿/twist、夹爪命令、接触模式），而非整机关节向量；并据此区分“任务相关变量”与“本体相关变量”；
  \item \textbf{推导} task-frame 与 body-frame 末端跟踪的坐标变换关系，证明 task-frame 目标对基座抖动\emph{一阶免疫}，并解释它对固定臂技能零样本迁移的意义；
  \item \textbf{掌握}在 $\SE(3)$/$\SO(3)$ 上对末端轨迹做限速平滑（位置 clip、姿态 $\boldsymbol\delta=\mathrm{Log}(\mathbf R_0^\top\mathbf R_1)$ 限幅、$\mathbf R'=\mathbf R_0\,\mathrm{Exp}(\boldsymbol\delta')$），并说清为何不能对四元数逐分量裁剪；
  \item \textbf{对比} RL-WBC、MPC-WBC、IK-QP 三类低层如何\emph{消费}同一末端命令，理解“几何可行 $\neq$ 动力学可行”，并写出阻抗/力控模式下的末端 wrench 与关节力矩映射 $\boldsymbol\tau_{ee}=\mathbf J_{ee}^\top\mathbf F$；
  \item \textbf{设计}抓取流水线的有限状态机（接近/对齐/预抓取/闭合/提升/搬运/放置），为每个阶段配上控制模式与安全门，并用滞回避免阶段抖动；
  \item \textbf{列出}固定臂操作技能迁移到移动基座所需的充分条件，并说明每个条件失效时的失败模式。
\end{enumerate}

\subsection*{本章知识导航}
本章是一条“\emph{把高层关心的‘手要去哪’，安全地翻译成移动机器人能执行的全身动作}”的主线。结构如下：

\begin{center}
\small
\begin{tikzpicture}[
  node distance=5mm and 7mm, >=Stealth,
  blk/.style={draw, rounded corners, align=center, font=\small, inner sep=3pt, fill=main!6, draw=main!60!black},
  grp/.style={draw, rounded corners, align=center, font=\small, inner sep=3pt, fill=second!6, draw=second!60!black},
  saf/.style={draw, rounded corners, align=center, font=\small, inner sep=3pt, fill=third!8, draw=third!60!black},
  ln/.style={->, thick, gray!60!black}]
\node[blk] (hi) {高层操作策略\\（任务意图）};
\node[grp, right=of hi] (frame) {task-frame\\目标表达};
\node[grp, right=of frame] (limit) {$\SE(3)$ 轨迹\\限速平滑};
\node[saf, below=7mm of frame] (filt) {可行性 / 安全\\过滤（限幅·门槛）};
\node[blk, right=of filt, fill=main!10, draw=main!60!black] (lo) {低层全身控制\\RL-WBC / MPC / IK-QP};
\node[saf, below=7mm of filt] (fsm) {抓取 FSM\\（阶段 + 模式）};
\draw[ln] (hi)--(frame); \draw[ln] (frame)--(limit);
\draw[ln] (limit.south) -- ++(0,-3mm) -| (filt.north);
\draw[ln] (filt)--(lo);
\draw[ln] (fsm.east) -| (lo.south);
\draw[ln] (fsm.north)--(filt.south);
\end{tikzpicture}
\end{center}

\noindent\textbf{推荐路径}：\cref{sec:si-why} 先讲清“为什么需要这层接口”（动机与反面），是全章的认识论地基；\cref{sec:si-taskframe} 是\emph{理论核心}——task-frame 变换与对抖动的免疫，移动操作的零样本迁移全靠它；\cref{sec:si-smooth} 给出 $\SE(3)$ 上的限速平滑，是把高层离散输出接到高频低层的安全垫；\cref{sec:si-lowlevel} 讲低层如何消费命令、为何要做可行性过滤、以及阻抗/力控如何并入；\cref{sec:si-grasp} 把连续跟踪组织成完整抓取任务，并落到迁移条件。\cref{ch:lie} 的 $\mathrm{Exp}/\mathrm{Log}$ 与 \cref{sec:si-smooth} 焊接，\cref{ch:operational-space} 的笛卡尔阻抗与 \cref{sec:si-lowlevel} 焊接。

\subsection*{前置知识桥接}
本章把若干前置章的结论“拉到接口层”使用，而不重证：
\begin{itemize}
  \item \cref{ch:lie} 已建立 $\SO(3)/\SE(3)$ 的指数/对数映射 $\mathrm{Exp},\mathrm{Log}$、右扰动 $\mathbf R\,\mathrm{Exp}(\delta\boldsymbol\phi)$ 与位姿复合/求逆的闭式——本章用它在群上做误差、插值与限速；
  \item \cref{ch:unified-multimodal-mpc} 给出复合机器人的统一动力学 \cref{eq:umm-core}，并把 $\SE(3)$ 末端误差 \cref{eq:umm-xierr} 写进了 MPC 代价——本章把“同一个末端误差”从 MPC \emph{内部}的代价，上提为高层与低层之间的\emph{接口约定}；
  \item \cref{ch:operational-space} 给出笛卡尔阻抗与“末端 wrench $\to$ 关节力矩”的对偶 $\boldsymbol\tau=\mathbf J^\top\mathbf f$——本章用它实现接口里的阻抗/力控模式；
  \item \cref{ch:mm-kinematics} 给出移动操作的联合运动学（底盘 $+$ 臂的联合雅可比 $\mathbf V_{ee}=[\mathbf J_b\ \mathbf J_a]\boldsymbol\nu$）——本章的低层把末端误差正是分配到“底盘/步态 $+$ 臂”这两组自由度上。
\end{itemize}

\subsection*{如果跳过本章会怎样}
\begin{itemize}
  \item 你会把“高层策略输出什么”当成 ROS 消息字段的命名问题，于是把一个整机关节向量直接喂给低层——系统短期像端到端，长期却\emph{难迁移、难调试、难保证安全}（\cref{sec:si-why} 的反面案例）；
  \item 你会在移动平台上直接复用固定臂学到的策略，却发现末端在物体附近反复抖——把锅甩给低层控制器，而真正的病根是\emph{目标表达在了晃动的身体上}（\cref{sec:si-taskframe}）；
  \item 你会让神经策略“自己学会”不输出跳变轨迹、自己学会不撞，于是在遮挡或分布外输入下，末端突然冲向远处——因为缺了一层\emph{确定性}的限速与可行性过滤（\cref{sec:si-smooth,sec:si-lowlevel}）。
\end{itemize}

\begin{note}
\textbf{本章记号与约定.}\;
基座（body）系记 $B$、世界系记 $W$、任务系（task-frame）记 $T$；位姿 ${}^{A}\mathbf{T}_{C}\in\SE(3)$ 读作“$C$ 系相对 $A$ 系的位姿”，满足复合律 ${}^{A}\mathbf{T}_{C}={}^{A}\mathbf{T}_{B}\,{}^{B}\mathbf{T}_{C}$。
末端（end-effector）记 $ee$，其世界位姿 ${}^{W}\mathbf{T}_{ee}$、机体位姿 ${}^{B}\mathbf{T}_{ee}$。
末端 $6$D 误差 $\boldsymbol\xi_{\mathrm{err}}\in\mathbb R^6$ 经 $\mathfrak{se}(3)$ 对数映射定义（沿用 \cref{eq:umm-xierr} 的左不变误差约定，$\boldsymbol\xi=[\boldsymbol\rho;\boldsymbol\phi]$，平移在前）。
末端 wrench $\mathbf F\in\mathbb R^6$，笛卡尔刚度/阻尼 $\mathbf K,\mathbf D$，末端雅可比 $\mathbf J_{ee}$。
高层下发的命令记 $(\cdot)^{\mathrm{cmd}}$；上标区分坐标系，例如 ${}^{T}\mathbf{T}_{ee}^{\mathrm{cmd}}$ 表示“在任务系中给定的末端目标”。
所有连续物理量取 SI 单位，姿态量取弧度。
\end{note}

% ════════════════════════════════════════════════════════════════
\section{为什么需要一层操作技能接口}
\label{sec:si-why}

\paragraph{动机：操作策略与全身控制解决的是两个不同问题。}
高层操作策略关心\emph{任务}：它从图像、语言、物体状态或人类演示中决定“末端该怎么动”——抓杯子、推门、擦桌、插孔。全身控制关心\emph{物理可行性}：它要在接触、摩擦、关节限位、力矩上界与平衡约束下，把这个末端意图真正执行出来。固定基座机械臂时，这两件事的耦合很弱，可以混在一起；一旦上了移动底盘、四足或人形，脚下接触不再固定、身体姿态会被末端动作反扰，二者必须\emph{显式分层}，中间用一层契约连接。

\paragraph{反面：高层直接输出整机关节目标会把麻烦留给未来。}
设想一个视觉策略直接输出 $19$ 自由度（四足 $12$ $+$ 臂 $6$ $+$ 夹爪 $1$）的关节目标，并在某个仿真里训得很好。它至少有三处脆弱：
\begin{enumerate}
  \item \textbf{跨平台不迁移}：换一台臂长、默认姿态、夹爪零位略不同的机器人，同一组关节角不再对应同一个末端行为。甚至同一台机器人，只要基座姿态有差，同一组臂关节角在世界系里的末端位置也变了——高层被迫去学\emph{本应由低层吸收}的本体差异。
  \item \textbf{失败难定位}：关节动作没有任务语义。任务失败时，你很难判断是视觉认错、目标不可达、腿没站稳，还是某个关节超限。若高层输出的是末端 $\SE(3)$ 目标，则“目标是否合理 / 是否可达 / 是否被低层跟住”可以分别检查。
  \item \textbf{安全无处插手}：把限速、自碰撞、支撑裕度这些安全逻辑塞进一个端到端关节策略里，既难训也难验证。
\end{enumerate}

\paragraph{历史：从“先到位、再抓取”到“边走边操作”。}
传统工业臂在固定基座上工作，操作规划只需考虑臂的工作空间。移动操作早期把问题切成两层：底盘导航到物体附近\emph{停稳}，再由机械臂规划抓取。四足臂与人形把问题再推进一步——机器人不是停稳后操作，而是在移动、平衡与接触中操作。一类有代表性的思路（如把固定臂上学到的技能搬到腿式平台的工作，UMI-on-Legs\cite{ha2024umi}）的关键洞见是：\emph{操作技能可以只输出任务系中的末端轨迹，由移动平台的低层全身控制器负责把它执行出来}。这就把“操作技能”与“移动平衡”解耦——Diffusion Policy、ACT、VLA 或遥操作只要遵守末端轨迹接口，就能接入同一套低层。

\begin{insight}[接口的本质是任务变量的选择]
设全身状态 $\mathbf x=(\mathbf q,\mathbf v)$。高层若输出关节动作，等于选了关节空间任务变量 $\mathbf y=\mathbf q_{\mathrm{act}}$；若输出末端轨迹，等于选了操作空间任务变量 $\mathbf y={}^{T}\mathbf{T}_{ee}\in\SE(3)$。二者由正运动学 ${}^{W}\mathbf{T}_{ee}=f_{\mathrm{kin}}(\mathbf q)$ 相连。\textbf{好的接口不是把信息传得越多越好，而是把任务相关与本体相关的变量分开}：高层表达任务意图，低层承担本体差异与物理可行性。抓杯子关心的是夹爪相对杯子的位姿，而不是第 $3$ 个肩关节转了多少度。
\end{insight}

\paragraph{接口的三层含义。}
只发一个 \verb|target_pose| 是不够的。一份完整的操作接口契约至少跨越三层——几何（目标在哪）、时间（何时到）、物理（能否执行）：

\begin{table}[H]\centering\small
\caption{操作技能接口的三层契约。只有几何层是不够的。}
\label{tab:si-contract}
\begin{tabular}{lll}
\toprule
层次 & 回答的问题 & 典型字段 \\
\midrule
几何接口 & 目标在哪个坐标系、什么位姿 & frame、$\SE(3)$ pose、twist \\
时间接口 & 何时到达、是轨迹第几帧 & 时间戳 stamp、时域 $H$、步长 $\Delta t$ \\
物理接口 & 是否允许接触、是否可执行 & 速度上界、控制模式、刚度、有效标志 \\
\bottomrule
\end{tabular}
\end{table}

\begin{pitfall}[思维陷阱：把接口当成软件格式问题]
新手常以为接口只是 ROS message 或 Python dict 的字段命名，于是先定消息格式、再补语义。实际上\textbf{接口选择先于格式}：它决定了任务变量、控制分工与迁移能力。应当\emph{先}从控制理论与任务语义确定“传末端 $\SE(3)$ 还是传关节”“目标定义在哪个系”“是否需要力控模式”，\emph{再}去写消息字段。另一个相关误区是“端到端就是没有接口”——即便用端到端视觉策略，真实系统仍需要时间戳、限幅、可行性检查与回退；端到端可以延伸到“感知 $\to$ 目标生成”，但控制执行始终需要明确边界。
\end{pitfall}

% ════════════════════════════════════════════════════════════════
\section{Task-frame 与 body-frame 末端跟踪}
\label{sec:si-taskframe}

\paragraph{动机：目标应绑定任务，而不是绑定晃动的身体。}
固定臂的基座不动，“末端相对基座的目标”与“末端相对世界的目标”几乎等价。移动机器人不同：四足或人形行走时，基座有平移、转动与高频晃动。但杯子、门把手、抽屉、桌面\emph{不会}跟着机器人身体晃。于是操作目标更自然地定义在\emph{任务系}（task-frame）——它可以是物体系、桌面系、门系，或一个世界局部系。低层负责把任务系目标换算成“当前基座下的可执行目标”。

\begin{derivation}[body-frame 目标对基座抖动不免疫]
设高层要求末端在基座前方固定位置 ${}^{B}\mathbf p_{ee}^{\mathrm{cmd}}=[0.5,\,0,\,0.2]^\top$（米）。其世界系目标为
\begin{equation}
  {}^{W}\mathbf p_{ee}^{\mathrm{cmd}}={}^{W}\mathbf p_{B}+{}^{W}\mathbf R_{B}\,{}^{B}\mathbf p_{ee}^{\mathrm{cmd}}.
  \label{eq:si-bodyworld}
\end{equation}
考察基座绕 $\mathbf z$ 轴（yaw）的小抖动 ${}^{W}\mathbf R_{B}\leftarrow{}^{W}\mathbf R_{B}\,\mathrm{Exp}(\delta\psi\,\mathbf e_z)$。对 \cref{eq:si-bodyworld} 在 $\delta\psi=0$ 处取一阶（用 $\mathrm{Exp}(\delta\psi\,\mathbf e_z)\approx\mathbf I+\delta\psi\,\mathbf e_z^\wedge$）：
\begin{equation}
  \delta\big({}^{W}\mathbf p_{ee}^{\mathrm{cmd}}\big)\approx{}^{W}\mathbf R_{B}\,(\delta\psi\,\mathbf e_z^\wedge)\,{}^{B}\mathbf p_{ee}^{\mathrm{cmd}}
  =\delta\psi\;{}^{W}\mathbf R_{B}\big(\mathbf e_z\times{}^{B}\mathbf p_{ee}^{\mathrm{cmd}}\big).
  \label{eq:si-jitter}
\end{equation}
其量级正比于基座抖动 $\delta\psi$ 乘以末端到基座的水平力臂 $\|\mathbf e_z\times{}^{B}\mathbf p_{ee}^{\mathrm{cmd}}\|$。这里力臂约 $0.5$ m，故\emph{毫弧度级}的基座抖动就会在末端目标上产生\emph{毫米级}的世界系晃动。低层会忠实跟踪这个晃动目标，结果是末端在杯子附近来回摆——这不是低层不够好，而是\emph{目标表达错了}。
\end{derivation}

\paragraph{task-frame 跟踪的核心变换。}
记基座世界位姿 ${}^{W}\mathbf{T}_{B}$（来自状态估计）、任务系世界位姿 ${}^{W}\mathbf{T}_{T}$（来自视觉/SLAM/物体检测）。Body-frame 控制给定 ${}^{B}\mathbf{T}_{ee}^{\mathrm{cmd}}$，其世界目标 ${}^{W}\mathbf{T}_{ee}^{\mathrm{cmd}}={}^{W}\mathbf{T}_{B}\,{}^{B}\mathbf{T}_{ee}^{\mathrm{cmd}}$ 随基座变化。Task-frame 控制改为给定 ${}^{T}\mathbf{T}_{ee}^{\mathrm{cmd}}$，其世界目标
\begin{equation}
  {}^{W}\mathbf{T}_{ee}^{\mathrm{cmd}}={}^{W}\mathbf{T}_{T}\,{}^{T}\mathbf{T}_{ee}^{\mathrm{cmd}}
  \label{eq:si-taskworld}
\end{equation}
只要任务系绑在桌面或物体上、不随机器人晃动，就\emph{与基座姿态无关}。低层每个周期把它换算回当前基座系下的等效目标：
\begin{equation}
  \boxed{\;{}^{B}\mathbf{T}_{ee}^{\mathrm{cmd}}=\big({}^{W}\mathbf{T}_{B}\big)^{-1}\,{}^{W}\mathbf{T}_{T}\,{}^{T}\mathbf{T}_{ee}^{\mathrm{cmd}}.\;}
  \label{eq:si-task2base}
\end{equation}
\cref{eq:si-task2base} 是 task-frame 跟踪的核心：基座一抖，${}^{B}\mathbf{T}_{ee}^{\mathrm{cmd}}$ 会\emph{反向}随之变化，恰好抵消，从而保持世界系目标稳定——抖动被“吸收”进了低层对臂关节的补偿里。

\begin{proposition}[task-frame 对基座运动的不变性]\label{prop:si-invariance}
若任务系世界位姿 ${}^{W}\mathbf{T}_{T}$ 与任务系下的命令 ${}^{T}\mathbf{T}_{ee}^{\mathrm{cmd}}$ 均不依赖基座位姿 ${}^{W}\mathbf{T}_{B}$，则由 \cref{eq:si-taskworld} 给出的世界系末端目标 ${}^{W}\mathbf{T}_{ee}^{\mathrm{cmd}}$ 对 ${}^{W}\mathbf{T}_{B}$ 的任意变化保持不变；从而由 \cref{eq:si-task2base} 算得的基座系目标随基座做\emph{精确}补偿，使世界系跟踪目标不受基座抖动影响（理想状态估计下成立）。
\end{proposition}
\begin{proof}
直接由 \cref{eq:si-taskworld}：右端不含 ${}^{W}\mathbf{T}_{B}$，故 $\partial\,{}^{W}\mathbf{T}_{ee}^{\mathrm{cmd}}/\partial\,{}^{W}\mathbf{T}_{B}=\mathbf 0$。而 \cref{eq:si-task2base} 是同一个世界系目标在当前基座系下的表达 $({}^{W}\mathbf{T}_{B})^{-1}\,{}^{W}\mathbf{T}_{ee}^{\mathrm{cmd}}$，故世界系目标恒定。注意结论依赖“${}^{W}\mathbf{T}_{B}$ 准确”这一前提：状态估计延迟或漂移会把这份不变性打折扣（见下方陷阱）。
\end{proof}

\paragraph{末端误差：必须用 $\SE(3)$ 对数，不能逐分量减。}
无论在世界系还是基座系算误差，姿态部分都要走群上的对数映射。沿用 \cref{eq:umm-xierr} 的左不变误差，
\begin{equation}
  \boldsymbol\xi_{\mathrm{err}}=\mathrm{Log}\!\big(({}^{W}\mathbf{T}_{ee})^{-1}\,{}^{W}\mathbf{T}_{ee}^{\mathrm{cmd}}\big)^{\vee}\in\mathbb R^6,
  \label{eq:si-xierr}
\end{equation}
它表达在当前末端（body）系中、其微分轨迹与参考无关，正是 MPC 与状态估计偏好的性质。\cref{ch:unified-multimodal-mpc} 把它放进了 MPC 代价 $J_{ee}=\tfrac12\boldsymbol\xi_{\mathrm{err}}^\top\mathbf Q_{ee}\boldsymbol\xi_{\mathrm{err}}$；本章把“同一个 $\boldsymbol\xi_{\mathrm{err}}$”从 MPC 内部上提为高层接口约定——高层下发的不是“某个关节角”，而是“某个任务系中的 $\SE(3)$ 轨迹”，误差始终在群上定义。

\begin{table}[H]\centering\small
\caption{常见操作任务的 task-frame 选择。task-frame 不是唯一的，可分阶段切换。}
\label{tab:si-taskframe-choice}
\begin{tabular}{lll}
\toprule
任务 & 推荐 task-frame & 原因 \\
\midrule
抓杯子 & 杯子系或桌面局部系 & 抓取几何与物体绑定 \\
推门 & 门铰链或门把手系 & 法向/切向清晰，便于分解力位 \\
擦桌子 & 桌面系 & 法向力与切向轨迹自然分离 \\
搬运物体 & 世界局部系或物体初始系 & 目标不随机器人晃动 \\
遥操作 & 操作者手柄初始系 & 手势相对运动稳定 \\
\bottomrule
\end{tabular}
\end{table}

\noindent task-frame 可分阶段切换：抓取前用物体系，抓住后切到夹爪或搬运目标系，放置时再切到桌面放置区域系——切换时必须保持轨迹连续，否则末端会在切换瞬间跳变。

\paragraph{task-frame 不替代 WBC，只是给它一个更稳的目标。}
task-frame 只定义“目标在哪”。低层全身控制仍要处理末端跟踪、支撑脚不滑、基座姿态稳定、关节限位、自碰撞、力矩限制与接触模式切换。task-frame 给低层一个\emph{更稳定、更任务相关}的目标，低层仍可拒绝或修正不可行目标：目标过远时 saturate；目标导致自碰撞时降级为保持当前末端；目标需要基座前移时，把末端误差转化为底盘或步态命令——后者正是 \cref{ch:mm-kinematics} 联合雅可比 $\mathbf V_{ee}=[\mathbf J_b\ \mathbf J_a]\boldsymbol\nu$ 在做的事：末端误差在“底盘 $+$ 臂”的零空间里分配。

\begin{insight}[task-frame 与零样本迁移的反事实]
若用 body-frame 承接固定臂学到的策略：固定臂训练时基座永远不动，策略看到的是“目标相对基座静止”的分布；移动部署时基座不断动，目标相对基座持续变化——策略被迫在\emph{训练分布之外}工作，零样本迁移随之破裂。task-frame 的作用，正是把基座抖动吸收进低层，让高层始终停留在“物体附近怎么动手”这个\emph{相对稳定}的问题上。这也是固定臂技能能迁到腿式平台的数学前提。
\end{insight}

\begin{pitfall}[左右乘顺序写反 / 误以为 task-frame $=$ world-frame / 误以为 task-frame 能消除一切抖动]
其一，$\SE(3)$ 变换\emph{不交换}：把 ${}^{W}\mathbf{T}_{B}\,{}^{B}\mathbf{T}_{ee}$ 写成 ${}^{B}\mathbf{T}_{ee}\,{}^{W}\mathbf{T}_{B}$ 会让目标方向完全错乱。务必在变量名里写清上下标（如 \verb|T_W_B|、\verb|T_B_ee|），让“相邻下标对消”成为可读的检查。其二，task-frame 不等于 world-frame——world-frame 只是 task-frame 的一种特例，物体系、桌面系、门系、工具系、手柄初始系都是合法 task-frame。其三，\cref{prop:si-invariance} 的不变性依赖“${}^{W}\mathbf{T}_{B}$ 准确且与 ${}^{W}\mathbf{T}_{T}$ 同步”：若 ${}^{W}\mathbf{T}_{B}$ 是当前时刻、${}^{W}\mathbf{T}_{T}$ 是 $100$ ms 前的视觉估计，二者\emph{不能直接相乘}——状态估计延迟、任务系漂移、低层带宽不足都仍会带来抖动。task-frame 只是第一步，仍需时间同步、滤波与阻抗调节。
\end{pitfall}

% ════════════════════════════════════════════════════════════════
\section{\texorpdfstring{$\SE(3)$}{SE(3)} 轨迹的限速与平滑}
\label{sec:si-smooth}

\paragraph{动机：高层是离散、低频、可能跳变的，低层是连续、高频、要求平滑的。}
Diffusion Policy、VLA 的推理频率（如 $10$ Hz）通常远低于低层控制频率（如 $100$ Hz）；遮挡、分布外输入、采样噪声还会让高层偶尔吐出突变目标。若让低层忠实跟踪一次 $30$ 厘米的跳变，末端会冲向远处。因此在高层与低层之间需要一层\emph{确定性}的限速平滑，把传给低层的目标约束在物理合理的单步变化内。它\emph{不替代}低层控制器，只保证喂进去的目标不跳变。

\paragraph{位置：欧氏空间内直接 clip。}
设上一周期末端位置 $\mathbf p_{\mathrm{prev}}$、本周期目标 $\mathbf p_{\mathrm{cmd}}$、速度上界 $v_{\max}$、步长 $\Delta t$。限幅后的位置增量
\begin{equation}
  \Delta\mathbf p=\operatorname{clip}\big(\mathbf p_{\mathrm{cmd}}-\mathbf p_{\mathrm{prev}},\,v_{\max}\Delta t\big)
  =\big(\mathbf p_{\mathrm{cmd}}-\mathbf p_{\mathrm{prev}}\big)\cdot\min\!\Big(1,\;\tfrac{v_{\max}\Delta t}{\|\mathbf p_{\mathrm{cmd}}-\mathbf p_{\mathrm{prev}}\|}\Big),
  \label{eq:si-poslimit}
\end{equation}
新位置 $\mathbf p'=\mathbf p_{\mathrm{prev}}+\Delta\mathbf p$。注意按\emph{向量模长}缩放而非逐分量裁剪，否则会改变运动方向。

\begin{derivation}[姿态：在 $\SO(3)$ 上限速，而非对四元数逐分量裁]
姿态\emph{不能}像位置那样逐分量减。给定当前姿态 $\mathbf R_0$ 与目标姿态 $\mathbf R_1$，二者之差是一个\emph{相对旋转}，其轴角即李代数元素
\begin{equation}
  \boldsymbol\delta=\mathrm{Log}\big(\mathbf R_0^\top\mathbf R_1\big)\in\mathbb R^3,
  \label{eq:si-rotdelta}
\end{equation}
$\|\boldsymbol\delta\|$ 是两姿态间的测地（最短）转角。限制单步转角不超过 $\omega_{\max}\Delta t$：
\begin{equation}
  \boldsymbol\delta'=
  \begin{cases}
    \boldsymbol\delta, & \|\boldsymbol\delta\|\le\omega_{\max}\Delta t,\\[2pt]
    \boldsymbol\delta\,\dfrac{\omega_{\max}\Delta t}{\|\boldsymbol\delta\|}, & \text{otherwise},
  \end{cases}
  \label{eq:si-rotlimit}
\end{equation}
再映回群得到新姿态
\begin{equation}
  \mathbf R'=\mathbf R_0\,\mathrm{Exp}(\boldsymbol\delta').
  \label{eq:si-rotnew}
\end{equation}
\cref{eq:si-rotnew} 的更新形如 $\mathbf R_0\,\mathrm{Exp}(\cdot)$，正是 \cref{ch:lie} 的\emph{右扰动}——增量定义在 $\mathbf R_0$ 的局部系中，沿测地线匀速插值，结果始终是合法旋转。
\end{derivation}

\begin{pitfall}[为什么不能对四元数逐分量 clip]
单位四元数位于 $S^3$ 球面上。对其四个分量分别做线性裁剪，几乎一定会让结果\emph{离开球面}，得到非单位四元数、对应非正交“旋转矩阵”，姿态随之缓慢漂移；同理，欧拉角逐分量限幅会在接近万向锁时给出毫无意义的“增量”，并可能产生 $\pm\pi$ 跳变。正确做法只有一条：把姿态差化为 $\SO(3)$ 上的测地增量 $\boldsymbol\delta$（\cref{eq:si-rotdelta}），在切空间里限幅（\cref{eq:si-rotlimit}），再用 $\mathrm{Exp}$ 映回（\cref{eq:si-rotnew}）。这也解释了 \cref{sec:si-taskframe} 为何坚持用 $\SE(3)$ 对数定义误差——群上的运算天然避开了这些坑。
\end{pitfall}

\paragraph{轨迹窗口与缓冲：用预瞄补低层与基座运动。}
只发当前单点目标，低层不知道目标接下来怎么变，只能被动追赶。更好的做法是发一段未来轨迹窗口
\begin{equation}
  \mathcal Y_{t:t+H}=\Big\{\big({}^{T}\mathbf{T}_{ee}^{\mathrm{cmd}}(t+k\Delta t),\;g(t+k\Delta t),\;m(t+k\Delta t)\big)\Big\}_{k=0}^{H-1},
  \label{eq:si-window}
\end{equation}
其中 $g$ 是夹爪命令、$m$ 是控制模式。低层可取第一帧作为当前跟踪点，也可把整窗输入策略做预瞄——移动平台据此\emph{提前}补偿基座运动（这正是 UMI-on-Legs 给低层未来 EE 窗口的用意）。由于高层 $10$ Hz、低层 $100$ Hz，工程上维护一个命令缓冲：高层每次写入一段未来轨迹，低层按当前时间从中插值取目标（位置线性插值、姿态用 \cref{eq:si-rotnew} 式的 $\SO(3)$ 插值）；高层短暂超时则继续执行缓冲中的安全轨迹，缓冲耗尽则进入保持或回退（\cref{sec:si-grasp}）。

\begin{codebox}[text]{接口层:末端命令的限速平滑（确定性，置于高层与低层之间）}
function limit_command(p_prev, R_prev, p_cmd, R_cmd, v_max, w_max, dt):
    # 位置:按模长缩放，不改方向（式 \ref{eq:si-poslimit}）
    dp        = p_cmd - p_prev
    p_new     = p_prev + dp * min(1, v_max*dt / (norm(dp)+eps))

    # 姿态:在 SO(3) 上限速（式 \ref{eq:si-rotdelta}-\ref{eq:si-rotnew}）
    delta     = Log(R_prev^T @ R_cmd)          # 测地增量(轴角) in R^3
    angle     = norm(delta)
    if angle > w_max*dt:
        delta = delta * (w_max*dt / angle)     # 切空间内限幅
    R_new     = R_prev @ Exp(delta)            # 右扰动映回群

    return p_new, R_new                        # 始终为合法 SE(3)
\end{codebox}

% ════════════════════════════════════════════════════════════════
\section{低层全身控制如何消费末端命令}
\label{sec:si-lowlevel}

\paragraph{动机：低层不是轨迹播放器。}
若低层只是回放关节轨迹，移动平台上的操作很容易失败：脚下接触在变、基座姿态在变、末端目标可能被高层延迟污染、物体接触会产生外力。低层必须在\emph{每个}控制周期重新解释末端命令——决定如何把末端误差分配到腿、腰、臂与基座，并决定不可行时如何降级。一个共同框架是：先用统一动力学 \cref{eq:umm-core}
\[
  \mathbf M(\mathbf q)\dot{\mathbf v}+\mathbf h=\mathbf S^\top\boldsymbol\tau+\textstyle\sum_{i\in\mathcal C}\mathbf J_{c,i}^\top\boldsymbol\lambda_i+\mathbf J_{ee}^\top\mathbf f_{ee}
\]
描述“腿 $+$ 臂 $+$ 浮基”的全部刚体物理（$\mathbf S=[\,\mathbf 0\ \mathbf I\,]$ 前 $6$ 列为零，对应不可直接驱动的浮基），低层在此约束下消费末端命令。

\begin{table}[H]\centering\small
\caption{三类低层对同一末端命令的消费方式。它们可共享同一接口、甚至混合使用。}
\label{tab:si-lowlevel}
\begin{tabular}{p{2.0cm}p{3.3cm}p{2.6cm}p{2.9cm}}
\toprule
低层类型 & 末端命令进入的方式 & 输出 & 取舍 \\
\midrule
RL-WBC & 把 $\boldsymbol\xi_{\mathrm{err}}$ 或未来窗口加入观测 & 关节位置目标 & 频率高、能学耦合；约束靠奖励/终止隐式表达 \\
MPC-WBC & 写成代价 $\ell_{ee}=\tfrac12\|\boldsymbol\xi_{\mathrm{err}}\|_{\mathbf Q_{ee}}^2$ & 接触力/基座轨迹/臂速度参考 & 约束显式、可解释；实时求解成本高 \\
IK-QP & 速度层最小二乘跟踪 $\dot{\mathbf x}_{ee}^{\mathrm{cmd}}$ & 关节速度/位置 & 简单可解释；不含动力学/接触，仅查几何可达 \\
\bottomrule
\end{tabular}
\end{table}

\paragraph{RL-WBC 的消费。}
RL 低层把末端命令并入观测：可加入当前误差 $\boldsymbol\xi_0=\mathrm{Log}(({}^{W}\mathbf{T}_{ee})^{-1}\,{}^{W}\mathbf{T}_{ee}^{\mathrm{cmd}}(t))^\vee$，也可加入未来窗口 $[\boldsymbol\xi_0,\boldsymbol\xi_1,\dots,\boldsymbol\xi_{H-1}]$（让策略知道末端将去哪）。策略输出全身关节位置目标，奖励里含末端跟踪、基座稳定、动作平滑与安全项。优点是高频；代价是约束靠训练学习与终止惩罚\emph{隐式}表达，不如 QP 显式。

\paragraph{MPC-WBC 的消费。}
MPC 把末端命令写成代价 $\ell_{ee}=\tfrac12\|\mathrm{Log}(({}^{W}\mathbf{T}_{ee})^{-1}\,{}^{W}\mathbf{T}_{ee}^{\mathrm{cmd}})\|_{\mathbf Q_{ee}}^2$，同时保留质心动量、摩擦锥、接触模式与自碰撞约束（皆见 \cref{ch:unified-multimodal-mpc}）。它输出较低频参考（接触力、基座轨迹、臂速度），由高频 WBC 转成关节力矩。更可解释，但四足臂/人形的实时求解很重，工程上常把 MPC 当慢速规划、WBC 或 RL 当快速执行。

\paragraph{IK-QP 的消费。}
IK-QP 在速度层求解
\begin{equation}
  \min_{\dot{\mathbf q}}\ \big\|\mathbf J_{ee}\dot{\mathbf q}-\dot{\mathbf x}_{ee}^{\mathrm{cmd}}\big\|^2+\lambda\|\dot{\mathbf q}\|^2
  \quad\text{s.t.}\quad \dot{\mathbf q}_{\min}\le\dot{\mathbf q}\le\dot{\mathbf q}_{\max},
  \label{eq:si-ikqp}
\end{equation}
其中阻尼项 $\lambda\|\dot{\mathbf q}\|^2$ 即 \cref{thm:ak-dls} 的阻尼最小二乘（DLS），在奇异/冗余处给出有界解（冗余与零空间见 \cref{sec:ak-redundancy}、奇异见 \cref{sec:ak-singularity}）；可再加关节居中、自碰撞近似与基座保持项。IK-QP 对固定臂常用，但对四足臂单靠 IK\emph{不足以}保证足端接触与基座稳定，更适合做高层命令的\emph{可达性预检}或站立慢速操作的基线。

\begin{pitfall}[“IK 成功”不等于“全身可行”]
最常见的误判：只要臂能 IK 到目标，就认为整机任务可执行。IK 只回答\emph{几何}可行性（关节空间存在解），它\emph{不考虑}足端摩擦锥、支撑多边形、接触力分配与基座动力学稳定。一个 IK 完全可达的末端目标，可能要求质心移出支撑域而让机器人摔倒。正确的层级是：IK 查几何可达 $\Rightarrow$ 全身控制（WBC/MPC）查动力学可行。因此低层必须有权\emph{拒绝}危险目标，并通过接口把 status（被限幅/拒绝/降级）回传给高层——不可行目标\emph{不必}强行追踪。
\end{pitfall}

\paragraph{执行前的可行性过滤。}
低层在调用策略\emph{之前}应先过滤命令——策略不该直接面对未过滤的高层输出：

\begin{table}[H]\centering\small
\caption{低层消费末端命令前的可行性过滤。失败时温和降级，而非非黑即白。}
\label{tab:si-feasibility}
\begin{tabular}{lll}
\toprule
检查项 & 方法 & 失败处理 \\
\midrule
工作空间 & 末端距离 / IK 残差 & 限幅或拒绝 \\
关节限位 & 预测 IK 解 & 投影到安全范围 \\
自碰撞 & 几何最近距离 & 停止臂运动或绕行 \\
支撑裕度 & 质心到支撑边界距离 & 降低末端跟踪权重 \\
速度限制 & $\|\Delta\mathbf x\|/\Delta t$ & 轨迹限速（\cref{eq:si-poslimit}） \\
目标年龄 & 当前时间 $-$ stamp & 丢弃过旧命令 \\
置信度 & 高层 confidence & 保持或减速 \\
\bottomrule
\end{tabular}
\end{table}

\paragraph{阻抗与力控模式：很多操作不是纯位置跟踪。}
推门要沿法向施力，擦桌要保持法向压力，插孔要低刚度以免卡死。故接口须含模式 $m$ 与刚度 $\mathbf K$。笛卡尔阻抗目标（承 \cref{ch:operational-space}）
\begin{equation}
  \mathbf F_{\mathrm{cmd}}=\mathbf K(\mathbf x_{\mathrm{cmd}}-\mathbf x)+\mathbf D(\dot{\mathbf x}_{\mathrm{cmd}}-\dot{\mathbf x})+\mathbf F_{\mathrm{ff}},
  \label{eq:si-impedance}
\end{equation}
低层再经静力学对偶把末端 wrench 转成关节广义力或 WBC 任务：
\begin{equation}
  \boldsymbol\tau_{ee}=\mathbf J_{ee}^\top\mathbf F_{\mathrm{cmd}}.
  \label{eq:si-tau}
\end{equation}
\cref{eq:si-tau} 是 $\boldsymbol\tau=\mathbf J^\top\mathbf f$ 在末端的实例——它正是 \cref{eq:umm-core} 里 $\mathbf J_{ee}^\top\mathbf f_{ee}$ 项的“主动施力”一侧。对 RL 低层，刚度 $\mathbf K$ 可作为观测，让策略在低刚度模式更柔顺、高刚度模式更精确。

\begin{insight}[力--位混合：选择矩阵活在 task-frame 里]
有的任务需要“某些方向力控、其余方向位控”（擦桌：法向力控、切向位控；开抽屉：拉出方向力控、其余对齐）。经典做法是选择矩阵 $\mathbf S\in\{0,1\}^{6\times6}$（对角，$S_{ii}=1$ 为力控）：
\[
  \mathbf F_{\mathrm{cmd}}=\mathbf S\,\mathbf f_{\mathrm{des}}+(\mathbf I-\mathbf S)\big(\mathbf K(\mathbf x_{\mathrm{cmd}}-\mathbf x)+\mathbf D(\dot{\mathbf x}_{\mathrm{cmd}}-\dot{\mathbf x})\big).
\]
关键在于：$\mathbf S$ 的方向应定义在 \emph{task-frame}（如桌面系 $z$ 轴为法向）而非世界系或机体系——这把 \cref{sec:si-taskframe} 的 task-frame 选择从“影响位置跟踪”推广到“同时决定力--位分解的方向”。选错 task-frame 会让力控与位控方向混淆。力--位混合的完整谱系（含 Mason/Raibert--Craig/Hogan 与无源性）见 \cref{ch:force-wbc}；而\emph{无接触时启用力控极其危险}——力控器会持续增大位移去够到目标力，前方没有物体时末端会一直前冲直到撞上或到达关节极限，故力控必须在确认接触（如末端力 $>2$ N）后才激活。
\end{insight}

% ════════════════════════════════════════════════════════════════
\section{抓取流水线与跨平台迁移}
\label{sec:si-grasp}

\paragraph{动机：抓取不是一个目标点。}
“把夹爪移到物体位置然后闭合”是过度简化。真实抓取要：接近物体、对齐抓取轴、控速逼近（别撞飞物体）、闭合并判断是否抓住、提升并保持平衡、搬运到目标、放置并松开。每个阶段的控制模式、容错与安全检查都不同。因此抓取应由\emph{有限状态机}（FSM）组织——FSM 不是退回传统规划，而是把各阶段的离散事件与安全约束\emph{显式化}，让连续控制器只负责当前阶段的物理执行；每阶段的目标仍可由高层策略给出。

\begin{table}[H]\centering\small
\caption{抓取流水线的典型阶段、控制模式与主要风险。}
\label{tab:si-grasp-stages}
\begin{tabular}{llll}
\toprule
阶段 & 目标 & 控制模式 & 主要风险 \\
\midrule
Search & 找到物体与任务系 & 感知更新 & 目标漂移 \\
Approach & 末端到预抓取位 & 位置控制 & 速度过快 \\
Align & 夹爪姿态对齐抓取轴 & 姿态控制 & 关节限位 \\
Pregrasp & 逼近接触前位 & 低速位置 & 撞到物体 \\
Close & 夹爪闭合 & 夹爪力/位置 & 夹空或夹碎 \\
Lift & 提升物体 & 位置 $+$ 稳定 & 负载扰动 \\
Transport & 搬运 & task-frame 跟踪 & 基座晃动 \\
Place & 放置 & 低速位置/阻抗 & 撞桌面 \\
Release & 松开 & 夹爪控制 & 物体滑落 \\
Retreat & 后撤 & 位置控制 & 二次碰撞 \\
\bottomrule
\end{tabular}
\end{table}

\paragraph{预抓取位：别把目标直接设成物体中心。}
直接瞄物体中心通常会撞。应先到一个\emph{预抓取位}，再沿接近方向低速贴近：
\begin{equation}
  \mathbf{T}_{\mathrm{pre}}=\mathbf{T}_{\mathrm{obj}}\,\mathbf{T}_{\mathrm{grasp}}\,\mathbf{T}_{\mathrm{offset}},
  \label{eq:si-pregrasp}
\end{equation}
其中 $\mathbf{T}_{\mathrm{grasp}}$ 是物体相对夹爪的理想抓取变换，$\mathbf{T}_{\mathrm{offset}}$ 是沿接近方向（如夹爪 $x$ 轴）后退一小段（如 $10$ cm）。低层先到 $\mathbf{T}_{\mathrm{pre}}$，再沿接近方向低速靠近——这能显著减小撞击。

\paragraph{抓取确认：别只看时间或夹爪是否闭合。}
夹爪闭合\emph{不等于}抓住：可能夹空、夹偏或夹住后滑落。应综合多信号判断——夹爪位置（是否到闭合目标）、夹爪速度（是否被物体挡住）、电机电流/力（是否产生夹持力）、物体运动（是否被带动）、末端力（接触是否异常）。无力传感器时，可用“夹爪位置误差 $+$ 电机电流”估计：完全闭合却无阻力多半夹空；闭合很少却电流很高多半夹到硬边或卡住。更稳妥的判据还要看\emph{提升后}物体是否随动。

\begin{codebox}[text]{抓取状态机骨架（阶段逻辑;每阶段目标可由高层策略给出）}
state := SEARCH
loop every control step:
    perc, rs, now := sense()                      # 感知 / 机器人状态 / 时间
    switch state:
        SEARCH:    if perc.object_found:               state := APPROACH
        APPROACH:  if rs.ee_pos_err < 0.08:            state := ALIGN
        ALIGN:     if rs.ee_rot_err < 0.15:            state := PREGRASP
        PREGRASP:  if rs.ee_pos_err < 0.02:            state := CLOSE
        CLOSE:     if rs.grasp_confirmed:              state := LIFT
                   elif rs.close_timeout:             state := FAILSAFE
        LIFT:      if rs.lift_height > 0.10:           state := TRANSPORT
        # ... TRANSPORT / PLACE / RELEASE / RETREAT 同理，各带安全门
    cmd := make_command(state, perc, rs, now)     # 生成本阶段模式的 EE 命令
    cmd := safety_filter(cmd, rs, now)            # 限速 + 可行性 + 滞回
    send_to_lowlevel(cmd)
\end{codebox}

\begin{pitfall}[阶段切换无滞回会抖 / 状态机不会削弱学习的泛化]
其一，若“误差 $<$ 阈值就进下一阶段、$>$ 阈值就退回”，传感器噪声会让状态在阈值附近来回穿越、反复抖动。修法是滞回（进/出用两个不同阈值）、最短停留时间与置信度累计。其二，有人担心 FSM 是手工规则、会限制策略——实际上 FSM 只表达\emph{安全与阶段结构}，高层仍可学习每阶段的连续目标；学习负责连续决策，FSM 负责离散安全边界，二者互补而非互斥。
\end{pitfall}

\paragraph{失败恢复：回退不等于急停。}
急停在接触/搬运任务里可能更危险（搬运中突然松开会掉落，推门时突然撤力可能反弹）。应按失败类型与任务阶段选择回退：高层短暂超时则保持上一安全轨迹并减速；目标置信度低则冻结末端、等重检；预抓取不可达则请求高层移动基座或重采样；基座不稳则收臂到安全位、优先站稳；接触力过大则降刚度、沿反向退让；夹持失败则回到 pregrasp 重试；自碰撞风险则停臂、保持腿稳。一个实用的组织方式是把安全状态嵌套在操作 FSM 之上：NORMAL（正常执行）$\to$ CAUTION（降速、提高稳定权重）$\to$ SAFE\_HOLD（停操作、收臂）$\to$ EMERGENCY（硬件安全流程），并用置信度、支撑裕度与接触力作为升降级触发量。

\begin{insight}[安全过滤要进训练闭环；回退要分级]
安全过滤器会\emph{改变}高层命令的分布。若高层训练时没见过这种过滤、部署时才加，策略会不适应而行为变差——故应把同一套过滤逻辑放进仿真训练闭环，让高层逐渐学会输出\emph{可执行}的命令；对 RL 低层，还应把过滤返回的 status/scale 作为观测，否则低层不知道目标为何突然变慢或被保持。与此呼应，安全阈值不应是单一硬阈值（过严会让机器人频繁拒绝任务、根本干不成活），而应分\emph{警告}（降速）、\emph{降级}（降权/限幅）、\emph{硬停止}三级——用一个 $\mathrm{scale}\in[0,1]$ 把“执行 / 拒绝”的二选一，扩成连续的温和降级。
\end{insight}

\paragraph{从单臂到双臂：耦合约束。}
人形的双臂操作（搬箱、叠衣、开瓶）中，两臂末端不再独立——它们通过共同操作的物体\emph{刚性耦合}。协作搬运模式下，接口不再下发两个独立末端目标，而是下发\emph{物体}目标位姿 $\mathbf{T}_{\mathrm{obj}}^{\mathrm{cmd}}$ 加左右抓取偏移，由低层解出每臂目标：
\begin{equation}
  \mathbf{T}_{\mathrm{left}}^{\mathrm{cmd}}=\mathbf{T}_{\mathrm{obj}}^{\mathrm{cmd}}\,\mathbf{T}_{\mathrm{obj}\to\mathrm{left}}^{\mathrm{grasp}},\qquad
  \mathbf{T}_{\mathrm{right}}^{\mathrm{cmd}}=\mathbf{T}_{\mathrm{obj}}^{\mathrm{cmd}}\,\mathbf{T}_{\mathrm{obj}\to\mathrm{right}}^{\mathrm{grasp}}.
  \label{eq:si-bimanual}
\end{equation}
这与多机协作的协同运动学在数学上同构（\cref{sec:ma-coopkin}），区别仅在于双臂同处一机、通信延迟近零；其内力管理（不做无用功的内力）与抓取矩阵理论分别见 \cref{sec:ma-internalforce} 与 \cref{sec:cm-grasp}，移动双臂的整机协同规划见 \cref{ch:arm-dualarm-planning}。

\paragraph{跨平台迁移的充分条件。}
固定臂技能迁到移动基座\emph{不是}无条件成立，至少要满足下列五条；任一失效都对应一类典型失败：

\begin{table}[H]\centering\small
\caption{固定臂操作技能迁移到移动基座的五个条件与失效后果。}
\label{tab:si-transfer}
\begin{tabular}{p{3.4cm}p{5.0cm}p{4.0cm}}
\toprule
条件 & 含义 & 失效后果 \\
\midrule
任务系一致 & 高层轨迹相对任务对象表达（task-frame） & 用 body/关节表达则不迁移 \\
夹爪几何相近 & 抓取宽度与接触点相似 & 抓取姿态不适用、夹不住 \\
工作空间覆盖 & 移动平台能把末端送到目标范围 & 够不到、需重定位基座 \\
低层跟踪精度足 & 末端误差落在策略容忍内 & 高层要 $1$ cm、低层只给 $10$ cm 则失败 \\
感知分布相近 & 相机视角/光照/遮挡不过度偏移 & 分布外、识别错或目标漂移 \\
\bottomrule
\end{tabular}
\end{table}

\noindent 这也解释了接口选择对 sim-to-real 的影响：关节角接口的 gap 最大（关节动力学差异直接传到末端），末端位姿接口居中（低层 IK/WBC 吸收部分差异），task-frame 末端接口最小——因为它定义在物体/环境上、不依赖本体，即便仿真动力学与真机有差，“手到杯子的相对位姿”在物理上仍一致。换言之，接口层不是仿真与真机之间的“透明通道”，而是一个能\emph{吸收一部分 sim-to-real gap} 的适配层。

% ════════════════════════════════════════════════════════════════
\section*{本章小结}

本章把“高层操作意图 $\to$ 全身可执行动作”之间的接口当成一个可推导、可调试、可迁移的工程对象。核心结论：

\begin{table}[H]\centering\small
\caption{本章要点速查}
\label{tab:si-summary}
\begin{tabular}{p{2.6cm}p{5.4cm}p{4.4cm}}
\toprule
知识点 & 核心结论 & 关键式 / 落点 \\
\midrule
接口分层 & 高层表达任务意图、低层承担物理可行性；接口是任务变量的选择 & \cref{tab:si-contract} 三层契约 \\
task-frame & 目标应绑任务对象而非晃动的身体；对基座抖动一阶免疫 & \cref{eq:si-task2base}、\cref{prop:si-invariance} \\
$\SE(3)$ 误差 & 姿态误差走对数映射，不可逐分量减 & \cref{eq:si-xierr} \\
轨迹限速 & 位置按模长 clip、姿态在 $\SO(3)$ 上限幅再 $\mathrm{Exp}$ 映回 & \cref{eq:si-poslimit,eq:si-rotnew} \\
低层消费 & RL/MPC/IK-QP 共享同一接口；几何可行 $\neq$ 动力学可行 & \cref{tab:si-lowlevel}、\cref{eq:si-tau} \\
抓取流水线 & 多阶段 FSM $+$ 阶段安全门 $+$ 滞回；回退分级非急停 & \cref{tab:si-grasp-stages} \\
跨平台迁移 & 迁移性来自 task-frame 表达与低层跟踪精度 & \cref{tab:si-transfer} \\
\bottomrule
\end{tabular}
\end{table}

\begin{table}[H]\centering\small
\caption{本章符号表}
\label{tab:si-symbols}
\begin{tabular}{ll}
\toprule
符号 & 含义 \\
\midrule
${}^{A}\mathbf{T}_{C}\in\SE(3)$ & $C$ 系相对 $A$ 系的位姿；$W$/$B$/$T$ 为世界/基座/任务系 \\
${}^{B}\mathbf{T}_{ee}^{\mathrm{cmd}}$ & 由 task-frame 命令换算到基座系的末端目标（\cref{eq:si-task2base}） \\
$\boldsymbol\xi_{\mathrm{err}}\in\mathbb R^6$ & $\SE(3)$ 末端误差，$=\mathrm{Log}(\mathbf T_{ee}^{-1}\mathbf T_{ee}^{\mathrm{cmd}})^\vee$（\cref{eq:si-xierr}） \\
$\boldsymbol\delta=\mathrm{Log}(\mathbf R_0^\top\mathbf R_1)$ & 两姿态间的测地增量（轴角），姿态限速用（\cref{eq:si-rotdelta}） \\
$v_{\max},\omega_{\max}$ & 末端线速度/角速度上界（限速用） \\
$\mathbf F,\mathbf K,\mathbf D,\mathbf F_{\mathrm{ff}}$ & 末端 wrench、笛卡尔刚度/阻尼、前馈力（\cref{eq:si-impedance}） \\
$\mathbf J_{ee}$ & 末端雅可比；$\boldsymbol\tau_{ee}=\mathbf J_{ee}^\top\mathbf F$（\cref{eq:si-tau}） \\
$\mathbf S\in\{0,1\}^{6\times6}$ & 力--位混合选择矩阵（定义在 task-frame） \\
$g,m$ & 夹爪命令、控制模式（位置/阻抗/力控） \\
$H,\Delta t$ & 轨迹窗口长度、采样步长（\cref{eq:si-window}） \\
\bottomrule
\end{tabular}
\end{table}

\paragraph{故障诊断速查。}
\begin{itemize}
  \item \textbf{末端目标方向反了} $\Rightarrow$ 多为 $\SE(3)$ 左右乘顺序或 frame 写反：打印 ${}^{W}\mathbf{T}_{B}$、做单位变换自检、可视化坐标轴（\cref{sec:si-taskframe}）。
  \item \textbf{世界系末端抖} $\Rightarrow$ 用了 body-frame 目标或任务系漂移：固定 task-frame、对比 world/body 目标、查状态估计延迟（\cref{prop:si-invariance} 陷阱）。
  \item \textbf{低层突然冲向远处} $\Rightarrow$ 高层目标跳变未限速：打印目标速度、启用 \cref{eq:si-poslimit,eq:si-rotnew} 的限速器、查视觉置信度。
  \item \textbf{高层卡顿后仍执行旧动作} $\Rightarrow$ 缺时间戳/过期检查：打印 $\mathrm{now}-\mathrm{stamp}$、设 max\_age、缓冲耗尽进保持（\cref{sec:si-smooth}）。
  \item \textbf{IK 可达却摔倒} $\Rightarrow$ 几何可行但动力学不可行：查支撑裕度、降末端权重、加支撑裕度门（\cref{sec:si-lowlevel} 陷阱）。
  \item \textbf{抓取阶段来回跳} $\Rightarrow$ FSM 无滞回：加最短停留时间、双阈值、平滑检测信号（\cref{sec:si-grasp} 陷阱）。
  \item \textbf{固定臂策略迁移失败} $\Rightarrow$ 数据不在 task-frame 或低层精度不足：查演示坐标系、训练注入低层误差、核对 \cref{tab:si-transfer} 五条件。
\end{itemize}

\begin{practice}[本章综合练习]
\begin{enumerate}
  \item \textbf{（task-frame 推导）} 仿 \cref{eq:si-jitter}，推导 body-frame 目标在基座 yaw 抖动 $\delta\psi$ 下的世界系\emph{姿态}变化一阶近似，并说明为何姿态抖动与力臂无关、只与 $\delta\psi$ 同阶。
  \item \textbf{（变换自检）} 随机生成 ${}^{W}\mathbf{T}_{B}$ 与 ${}^{T}\mathbf{T}_{ee}^{\mathrm{cmd}}$，验证由 \cref{eq:si-task2base} 算得的 ${}^{B}\mathbf{T}_{ee}^{\mathrm{cmd}}$ 满足 ${}^{W}\mathbf{T}_{B}\,{}^{B}\mathbf{T}_{ee}^{\mathrm{cmd}}={}^{W}\mathbf{T}_{T}\,{}^{T}\mathbf{T}_{ee}^{\mathrm{cmd}}$（数值误差应 $<10^{-6}$）。
  \item \textbf{（姿态限速）} 证明 \cref{eq:si-rotlimit,eq:si-rotnew} 给出的 $\mathbf R'$ 与 $\mathbf R_0$ 之间的测地距离恰为 $\min(\|\boldsymbol\delta\|,\omega_{\max}\Delta t)$；并据此说明为何对四元数逐分量 clip 不具备这一性质。
  \item \textbf{（低层可行性）} 写出 \cref{eq:si-ikqp} 的扩展形式，加入“关节居中”与“基座姿态保持”两项软约束，并讨论 $\lambda$ 与这两项权重对奇异附近行为的影响（联系 \cref{thm:ak-dls}）。
  \item \textbf{（抓取 FSM）} 为“拿起杯子并放到桌面右侧”写出完整 FSM：每阶段的 task-frame、控制模式、进入/退出阈值（带滞回）与回退策略；指出哪一步必须从物体系切到桌面放置区域系。
\end{enumerate}
\end{practice}

\paragraph{承上启下。}
本章把 \cref{ch:unified-multimodal-mpc} 的统一动力学与 SE(3) 末端代价、\cref{ch:operational-space} 的笛卡尔阻抗、\cref{ch:mm-kinematics} 的联合运动学，统一收束到“高层意图 $\to$ 末端 SE(3) 轨迹 $\to$ 全身控制”这一条接口上。其中双臂协作的耦合约束 \cref{eq:si-bimanual} 指向 \cref{ch:arm-dualarm-planning} 的移动双臂协同规划；而当高层命令与低层控制\emph{都}来自人类演示或动捕数据时，如何从参考运动、对抗先验、技能潜空间与重定向训练出可部署的全身行为，则归入运动模仿与全身技能学习部。
               % ch:mm-skill-interface
